\subsubsection*{Starting the automorphic construction at the assigned levels}
\label{subsec:automorphic-realization}
\hypertarget{source-automorphic-realization}{}

The field, levels and weight functions have been chosen at a translate. We now give the Hecke and character calculations used in Construction~\ref{hyp:family-prescribed-realization}. The first set of signed tests selects the characters. The second calculation determines their actual spaces and local orders on the rebuilt member.

\begin{paramdefinition}[Fields, frames and levels at a translate]\label{hyp:initial-adelic-frame}
Fix the auxiliary CM field \(E\) of Definition~\ref{def:family-original-host},
one translate, its two weight functions, and all the
primes assigned to the sequences. Denote the small field and the two coefficient
fields by \(K_s\subset K_0\subset K\subset\mathbb C\), in
their distinguished embeddings. The underlying Hermitian spaces,
the marked split identifications, the two deletion characters
\(\delta,\delta'\) of Lemma~\ref{lem:split-deletion-characters}, and the levels at the primes assigned to the sequences belong to the initial data.
Let \(\Xi\) denote the finite set of tuple characters. The successive field requests are finite sets of algebraic numbers with
\[
 S_1(K_s)\subseteq K_0,\qquad S_2(K_s,K_0,\Xi)\subseteq K .
\] Combine requests at each stage by taking their union with the preceding arguments fixed. The prescribed rational prime set includes the two type primes, \(d\), and every prime at most \(X''\).
\end{paramdefinition}

To meet these finite requests, first fix the embedded small field, then adjoin the union of the first sets, and only then evaluate the second sets in that field and take their compositum. Every extension in these two steps is finite. The field cost for the geometric requests on the initial frames is proved in Lemma~\ref{lem:marked-original-geometric-field}.
\begin{lemma}[Initial characters and support over the smaller field]
For Definition~\ref{hyp:initial-adelic-frame} with trivial auxiliary characters,
\(\Xi=\{1,\delta,\delta'\}\). For every \(a\in\mathcal O_{K_0}\), the support test factors through contraction:
\[
 v(a)>0\quad\Longleftrightarrow\quad a\in v\cap\mathcal O_{K_0}.
\]
Thus any support selected by these tests is a union of fibres of the contraction map.
\end{lemma}
\begin{proof}
The reference contributes \(1\). At a coefficient place the
deletion rule uses \(\delta'\) over the first deletion prime
and \(\delta\) otherwise. The three distinguished tuples realize
the three listed values. Thus the final finite field request
\(S_2(K_s,K_0,\Xi)\) is determined before localization in \(K\);
adjoining it requires no new word level. For \(x\in\mathcal O_{K_0}\),
\[
 v(x)>0\quad\Longleftrightarrow\quad x\in v\cap\mathcal O_{K_0}.
\]
The positivity tests coincide at two places exactly when their
contractions coincide. Apply this to the prospective nonbase support.
\end{proof}

Begin with the untwisted saturated form of the marked adelic datum.
Its localization set is empty and its prescribed orders are
\[
                 \eta_j(v)=0\qquad(1\le j\le m,\ \text{all selected }v).
\] The two
weight comparisons used in making the choice are the lower
unitary comparison and the comparison at the residue characteristic
of the lift. The first division lattice at this step uses the zero
orders, the unit coordinate frame and the full word-column lattices.
It gives a starting row before the local arithmetic orders are
inserted. The row stage comes from the lowest discrete series at
the specified archimedean weights and the finite types on that
adelic datum. \label{cor:automorphic-realization-1}

\paragraph{Forming the adelic row on the prescribed group.}
Retain the original Hermitian frame, the finite level group and
the marked split half. Let \(P\) be the finite rational prime set in
Definition~\ref{hyp:initial-adelic-frame}. Before restricting the
support, allow the entire set of coefficient places
\[
             \{v:v\mid p\text{ for some }p\in P\},
\]
which is finite and stable under conjugation. The split-half condition is used even
when two different tuple conductors contain places at the same
rational prime. Both must use the same chosen half there. The
place of each real subfield is the one below the specified
column place, so the pinned choice gives this condition for
all the conductors at once.

Use the active tuple characters and their conductor sets in the
level datum. The distinguished tuple has character one. At the
two deletion characteristics use the prescribed split side, and
carry the original ramification test at every other place. If \(S_{\rm type}\) is the set of chosen type places and \(\bar G_v\) the matrix group given by reduction there, the quotient is identified place by place:
\[
 G_{\rm type}\simeq\prod_{v\in S_{\rm type}}\bar G_v .
\] Use the diagonal self-dual
frames when they have been supplied by the Hermitian frame.

In
the standard lattice case the coordinate frames with unit scales
are self-dual. These give the frames for the row and each of the
lifts. The finite group and character comparison then supplies
the type projectors on this product, together with the inert place
and tuple assignments. The identity map of the finite group can
serve as the preliminary label map. The Fock and lowest-type
models are to be built for these very groups and characters.
\label{cor:automorphic-realization-2}


\begin{proposition}[Adelic realization on the prescribed frame]
Use the fields and levels of Definition~\ref{hyp:initial-adelic-frame}
and the marked self-dual frames at the chosen split halves. Take the
level volumes in their common Haar measure, the regular discrete-series
parameters and their specified lowest compact types, and the definite
standard lattices with \(n_V\ge 8fr+4\) in the matched geometric model.
Then the initial row admits the prescribed real
intertwiners and theta lifts on the chosen finite types. At the distinguished tuple, if \(\theta\) is the theta series of a rational input and \(f\) its algebraic class sum, the normalized comparison is
\[
 f=\theta .
\]
The phase is one for this normalization.
\end{proposition}
\begin{proof}
The finite-place characters take values in the fixed coefficient
subfield on ideles with trivial archimedean component. The analytic construction
uses these character values. The proof uses the volume tie on this
same group to supply the core, its coefficient algebra and these
finite character values. The real models at every place use
the regular parameter of the prescribed weight and its
specified lowest compact type. The discrete series of the
determinant cover give these models. We use Harish--Chandra's
discrete-series theorem for the regular parameters and the
Casselman--Wallach theorem to extend the admissible
Harish--Chandra-module intertwiners to the smooth
moderate-growth globalizations.

For the Harish--Chandra modules \(V,W\) in this comparison and their
smooth moderate-growth globalizations \(V^\infty,W^\infty\), the
Casselman--Wallach equivalence identifies
\[
 \operatorname{Hom}_{G}(V^\infty,W^\infty)
       \simeq\operatorname{Hom}_{(\mathfrak g,K)}(V,W).
\]
Thus the extension preserves the intertwining
identities on the dense compact-finite vectors. With the
chain and the finite-action identities on these models, take
the theta map from the lowest-type Fock coefficient map.
Rapid decay of the cusp input gives the convergent integrals
for the intermediate lifts. The final projected lift uses the
good projector built on the chosen finite quotient. This is
the role of the intertwiner and lift existence assertions in
the fine datum. \cite{harish-chandra-degree,casselman-wallach}\label{cor:automorphic-realization-3}

For the geometric realization use the standard definite lattices.
The dimension inequality required there is
\[
                          n_V\ge8fr+4 .
\]
Use the same frame for the volume normalization and the algebraic
identifications. At each lift the finite comparison has one scalar for all allowed tuples. Determine it at the distinguished tuple in this common frame.

To determine the scalar relating the two forms, take the
theta series \(\theta\) of a rational input there, and let \(f\)
be its algebraic class sum. For a fixed class \(x\), the definite theta integral is a finite class sum. 

Let \(K_f\) be its finite level, \(h\) range over the double cosets, and \(\Gamma_h\) be the stabilizer in that class (with the same central quotient in the measure and the class sum). Write \(k(x,h)\) for the already normalized kernel times the input, paired at infinity. Splitting the integral into the individual level orbits gives
\[
 \theta(x)=\sum_{[h]}\frac{\operatorname{vol}(K_f)}{|\Gamma_h|}\,k(x,h)=f(x).
\]
Here each orbit has the indicated mass by the finite stabilizer formula; the second equality uses the same character, kernel and measure in the algebraic class sum. Evaluation over all classes therefore gives
\[
  \bigl({\rm ev}_x(f-\theta)\bigr)_x=0,\qquad
  \bigcap_x\ker({\rm ev}_x)=0 .
\]
Thus \(f=\theta\) with phase one. The weight scalars and geometric
indices transport this equality on the same levels, finite
group, tuple characters and selected places. The construction
keeps the intertwiners, group action and lattices together
in the normalized adelic realization. \label{cor:automorphic-realization-4}
\end{proof}

\begin{proposition}[The same multiplicity on the linked real and finite types]\label{prop:linked-type-multiplicity}
For the linked unitary groups with the same finite level and
the real and split actions in the frames of
Definition~\ref{hyp:initial-adelic-frame}, write \(m_H(\sigma)\) and \(m_{H'}(\sigma')\) for the multiplicities of matching types in the two groups. At either initial weight the trace comparison gives
\[
 m_H(\sigma)=m_{H'}(\sigma').
\]
\end{proposition}

\begin{proof}
The multiplicity assertion compares the ranks of spaces of
finite-adelic intertwiners into cuspidal forms of fixed real
discrete-series type. It compares the row group with the group
of the indicated lift. Their odd ranks agree, their real
parameters agree at every real place with adjacent gaps at
least \(100r^4\), and the finite group isomorphism identifies
the lattice frames modulo the type places. The selected type
has the no-fixed-vector condition along the prescribed proper
parabolic radicals. The representations in the comparison are
irreducible admissible and have a level-fixed vector surviving
the type projector. Denote the two groups by \(H,H'\) and the
transport of a finite representation \(\sigma\) by \(\sigma'\).
With the real models included in the multiplicities, the required identity is
\[
                         m_H(\sigma)=m_{H'}(\sigma').
\]
\label{cor:automorphic-realization-5}

\proofstep{Match the projected trace at the prescribed weight.}
At each real place use the pseudo-coefficient normalized to
select the specified discrete series for the trace. At the
finite places use the type tests and an auxiliary split place
outside their support. The constant terms along proper
rational parabolics vanish for these tests. The prescribed
split test allows only regular semisimple elements on the
diagonal kernel, and the nonfield terms vanish in the class
integrals. The trace expansion is then a finite sum over
the regular polynomial classes with field centralizer. If \(T\) is the common torus identifying the two centralizers, use a single measure \(dt\) on it. For the corresponding identifications \(i,i'\) the measures on the two sides are
\[
 d\mu_{\mathrm{cent}}=i_*dt,\qquad d\mu'_{\mathrm{cent}}=i'_*dt .
\]
Transport the finite Haar measures with the frame isometries.
This identifies the finite orbital factors, including the
normalizations of the lattice stabilizers.
\label{cor:automorphic-realization-6}

At a real place where the torus is compact, write its distinct
eigenvalues as \(z_1,\ldots,z_r\). The real conjugacy classes
are the choices of which eigenlines have positive sign, with
the number of signs determined by the real form. Write \(W_K\subset S_r\) for the stabilizer of the sign pattern.
The classes provide the disjoint decomposition
\[
                         S_r=\coprod_y y^{-1}W_K .
\] For ordered integral
Harish--Chandra parameters \(\lambda_j\), put
\(a_j=\lambda_j+(r-1)/2\). The discrete-series character
on the class indexed by \(y\) is the alternating sum over
\(W_K\) on the reordered eigenvalues divided by their
Vandermonde. Its overall factor \((-1)^{pq}\) is one
because \(p+q=r\) is odd.

Moving each denominator to the
same ordering contributes the sign of \(y\), so the sum of the
quotients is
\[
   \frac{\det(z_i^{a_j})_{1\leq i,j\leq r}}
        {\det(z_i^{r-j})_{1\leq i,j\leq r}}.
\]
Indeed, expanding the compact alternating sums on the cosets gives
\[
 \sum_y\sum_{w\in y^{-1}W_K}\operatorname{sgn}(w)
         \prod_{i=1}^r z_i^{a_{w(i)}}
 =\sum_{w\in S_r}\operatorname{sgn}(w)\prod_{i=1}^r z_i^{a_{w(i)}}
 =\det(z_i^{a_j}).
\] These are the Weyl character
coordinates of highest weight \(a_j-r+j\). In the orbital
integral the formula is conjugated and divided by the same
torus volume. At a real place with noncompact centralizer
the orbital integrals of the chosen pseudo-coefficients
vanish. This computes the two real stable factors in terms
of the same eigenvalues and weights even when the signatures
are different. Thus, writing \(h'\) for the transport of the finite test \(h\), the equality assembled from the orbital integrals is
\[
 \operatorname{tr}_{\rm cusp}(h)=\operatorname{tr}'_{\rm cusp}(h').
\] \label{cor:automorphic-realization-7}

Choose a common compact open subgroup on the finite side
for the tests, and transport it to \(H'\). The fixed-vector
spaces in the corresponding discrete-series multiplicity
spaces have finite dimension by cuspidal spectral finiteness.
Make one finite list \(\sigma_1,\ldots,\sigma_d\) of distinct
irreducible admissible representations covering both spectra
seen by the tests and include the given \(\sigma\). The
spectral trace and the preceding equality give, for every
allowed finite test \(h\),
\[
   \sum_{i=1}^d a_i\operatorname{tr}\sigma_i(h)=0,
   \qquad a_i=m_H(\sigma_i)-m_{H'}(\sigma_i').
\]
Let \(e_v\) be the local convolution idempotent for the chosen
type at a type place, the normalized Iwahori indicator at
the auxiliary split place \(u\), and the normalized maximal
compact indicator at the remaining places. Their pure tensor
acts on \(\sigma_i\) by a finite-rank idempotent \(P_i\).

The resulting projector is nonzero on the stipulated \(\sigma\). For a finite-adelic
element \(g=(g_v)_v\), set
\[
 b_v(x)=e_v(g_v^{-1}x),\qquad f_v=e_v*b_v.
\]
Since \(b_v*e_v=b_v\), the test at a type place satisfies
\(e_v*f_v*e_v=f_v\). Away from the type places and \(u\)
it is bi-invariant under the maximal compact. The family
agrees with the unit Hecke family almost everywhere. At \(u\)
it is left Iwahori-invariant. Replacing only \(f_u\) by
\(1_{ItI}\), for each \(t\) in the strict dominant torus
cone, preserves all the conditions for the trace comparison.
The original pure tensor acts by \(P_i\sigma_i(g)P_i\).
\label{cor:automorphic-realization-8}

For a local isomorphism class \(\lambda\) at \(u\), collect all
the factors away from \(u\) in the coefficient
\[
 b_\lambda=\sum_{i:\,\sigma_{i,u}\simeq\lambda}
             a_i\prod_{v\ne u}{\rm tr}\,\sigma_{i,v}(f_v).
\]
Classes without Iwahori invariants contribute zero. For \(t\) in
the strict dominant cone the trace relation therefore reads
\[
 0=\sum_\lambda b_\lambda{\rm tr}\,\lambda(1_{ItI})
   ={\rm vol}(I)\delta_B(t)^{-1}
       \sum_\lambda b_\lambda\sum_{\xi\in J(\lambda)}\xi(t).
\]
Here \(J(\lambda)\) lists the unnormalized Jacquet exponents with
multiplicities (normalized exponents give the common factor
\({\rm vol}(I)\delta_B(t)^{-1/2}\)). The cited exponent separation
on the local classes in this comparison forces \(b_\lambda=0\).
Reinsert \(f_u\). Cyclicity of finite-rank trace gives
\({\rm tr}(P_i\sigma_i(g)P_i)={\rm tr}(\sigma_i(g)P_i)\),
so \(\sum_i a_i{\rm tr}(\sigma_i(g)P_i)=0\) for all \(g\).
\cite{zelevinsky-induction}\label{cor:automorphic-realization-9}
The commutant of \(\sigma_i\) is \(\mathbb C\): an endomorphism
has an eigenvector on a nonzero finite-dimensional fixed space,
and that vector generates the irreducible module.

For a prescribed
index \(k\), apply Jacobson density to bases of the finite spaces
\({\rm im}\,P_i\) in the distinct modules. There is a group
algebra element \(a\) such that
\[
 \sigma_i(a)|_{{\rm im}\,P_i}=\delta_{ik}\,{\rm id},\qquad
 0=\sum_i a_i{\rm tr}(\sigma_i(a)P_i)=a_k\operatorname{rk}P_k .
\]
Since \(P_k\ne0\) at the selected type, its multiplicities agree.
The identical finite and real actions transport this equality at
both initial weights. \label{cor:automorphic-realization-10}
\end{proof}

\begin{lemma}[Rational separation by common row tests]\label{lem:rational-row-separation}
Use the prescribed tests of Definition~\ref{hyp:fitted-pointwise-nonvanishing} and their nonvanishing in Lemma~\ref{lem:fitted-completed-nonvanishing}. For the rational row maps \(B_t:A_t\to E_{\mathbb C}\) at the allowed tuples, one finite coefficient request supplies complex-linear functionals \(\lambda_a\) with
\[
 \bigcap_a\ker\lambda_a=0,\qquad\lambda_a\bigl(B_t(A_t)\bigr)\subseteq K\quad\text{for every }a,t,
\]
where \(K\) denotes its specified embedded copy.
The fixed-field and support comparisons used at rank zero
preserve these maps under restriction of the coefficient
places.
\end{lemma}

\begin{proof}
Write \(K\) for the large coefficient field with its fixed
embedding, \(A_t\) for a rational row of tuple \(t\), and
\(B_t(x)=\Theta_t(j_t x)\) for its image in the
common complex row space \(E_\mathbb C\). We need complex-linear
functionals \(\psi_u\) with zero common kernel on
\(E_\mathbb C\) and every \(\psi_u(B_t(x))\) in the
embedded \(K\). This is the field supplied by the finite
coefficient request on the actual tuple set.

For each allowed finite test type \(\alpha\), let
\(C_\alpha\) be the column space and \(c_{\alpha,u}\) its fitted
spanning columns. The common theta integral identifies the row
tests with the Bessel pairings for the specified real vector and
covector. For a row \(y\) and \(c=\sum_u a_uc_{\alpha,u}\), convergence
and linearity give
\[
 {\cal B}_\alpha(c,y)=\sum_u a_u\,{\cal B}_\alpha(c_{\alpha,u},y).
\]
If all the column tests vanish on \(y\), linearity kills every pairing on every \(C_\alpha\). On the given base multiplicity space Proposition~\ref{prop:fitted-grid-base-separation} chooses one real test and a basis of the dual from these allowed tests. Evaluation against this basis gives \(y=0\).
A finite family suffices: if the current common kernel has positive
dimension, choose one further test nonzero there. Each choice
strictly lowers that dimension. \label{cor:automorphic-realization-11}

For the field assertion first apply the rational good
projector to the input \(x\), giving \(x'\). If \(P\) denotes the good projector, equivariance and the good target give
\[
 B_t(x')=B_t(Px)=P B_t(x)=B_t(x).
\] Write \(C_t=\prod_{i,j}c_{ij,t}\) for
the product of the nonzero fitted square constants and
\(0\) for the distinguished tuple. The finite Fourier and
class sum reads \(v_{t,u}(x')\) in the embedded coefficient
field. Write \(\iota\) for the coefficient embedding, \(\zeta_t\)
for the coefficient-field root of unity, and \(\kappa\) for
the single complex scalar for the chosen test family.
On the good input the formula reads
\[
   \iota(\zeta_t) C_t\,\psi_u(B_t(x))
     = \kappa v_{t,u}(x').
\]
The bottom toric comparison in its nonvanishing integer range gives
\(C_t/C_0=\prod_{i,j}(c_{ij,t}/c_{ij,0})\in\iota(K)^\times\).
If \(\kappa\ne0\), set \(\psi'_u=(C_0/\kappa)\psi_u\). Then
\[
 \psi'_u(B_t x)=\frac{v_{t,u}(x')}{\iota(\zeta_t)(C_t/C_0)}
             \in\iota(K),\qquad \bigcap_u\ker\psi'_u=0.
\]

If \(\kappa=0\), the displayed Fourier identity gives
\(\psi_u(B_tx)=0\), already in the field. Make the request on the whole
tuple family so that these comparisons use the same field.
For a complex row with conjugate scalar action, precompose
with its conjugate-linear identification and conjugate the
output. The resulting functionals are complex-linear, and
their values are in the conjugate coefficient embedding.
If \(X_1,X_2\) are the cutoffs for the test and field assertions, impose them simultaneously by taking
\[
 x\ge \max\{X_1,X_2\}
\]
and increasing the scale further whenever the application requires it. The Hida size setting is supplied
on a bare state with the same group, so the tests pull back
to the original state along its complex row isomorphism.
Their values and field are preserved. This is the second
assertion at the residue-characteristic weights. It is used
jointly with multiplicity on the restricted datum.
\label{cor:automorphic-realization-12}

With respect to a \(K\)-basis \((a_i)\) of \(A_t\), the normalized
tests have coefficient matrix
\[
               M_t=\big(\iota^{-1}\psi_u(B_ta_i)\big)_{u,i}.
\]
The finite test map \(\Psi=(\psi_u)_u\) is injective, hence
\[
 \ker\Theta_t=\ker(\Psi\Theta_t)=j_t(\mathbb C\otimes_K\ker M_t).
\]
The last equality follows by row reduction over \(K\). Equivariance
of \(\Theta_t\) makes this descended kernel stable under the level group.

To pass from the full allowed place set to the actual nonbase
support, restrict to the places detected by the index tests,
together with their conjugates. The other places have trivial
auxiliary character. Restrict the tuple-indexed data along
the inclusion of the smaller set, then update the weight
scalars on the restrictions of the same group and geometric
maps. The two index tests commute with that restriction.
They depend on the place readings retained on the restriction
and are unchanged by the accompanying scalar update.

Denote restriction to the detected places and their conjugates by \(R\), and the support recovered from the two tests by \(S\). The preceding compatibility gives
\[
 S(RD)=S(D)
\]
for the datum \(D\) under consideration. Thus a second restriction discards no further places. The residue-characteristic weight comparison is
applied \emph{after} this step. Its hypotheses include equality
with the recovered support, the definite weight condition,
the core and group conditions and membership of all residue
characteristics in the allowed prime set. It returns the
comparison used to construct the first row-lift family.
\label{cor:automorphic-realization-13}

The fixed-field part of the support argument uses an earlier fine
adelic datum on the same pre-level frame. For the fixed small field
one requests a finite list of algebraic coefficients in \(K_0\)
and a finite list of residue characteristics. The uniformity here
is over the generator choices and auxiliary numerical parameters
at the prescribed levels. Let \(p\) be a type prime or the prime
assigned to a sequence. The norm relation gives the same bound at every embedding:
\[
              \pi\bar\pi=p\quad\Longrightarrow\quad
               |\sigma(\pi)|=\sqrt p\qquad(\sigma:E\hookrightarrow\mathbb C).
\] In the
Minkowski embedding the possible \(\pi\)'s lie in the finite set
\[
 \mathcal O_E\cap\prod_{\sigma:E\hookrightarrow\mathbb C}
                          \{z:|z|\le\sqrt p\}.
\]

For each occurring tuple of generators at the finitely many word
and type primes, choose one fine reference realization with its
lower-unitary comparison. Its two zero-content index tests compare
full lattices in a common space. For each pair choose \(N>0\) clearing both lattice inclusions.
Localizing gives
\[
 NL\subseteq L',\quad NL'\subseteq L
        \quad\Longrightarrow\quad L_v=L'_v\quad(v\nmid N).
\] Request the finitely many resulting
residue primes and their conjugates, and an embedded rational basis
of the reference coefficient field. The comparison at fixed generators
preserves the lower-unitary maps and these zero-content tests when
the other numerical parameters vary. Thus the union over reference
tuples is finite before \(K_0\) is chosen.

The chosen reference coefficient field embeds into \(K_0\).
Transport its support with zero added content to \(K\).
In the enlarged datum this inherited support is supplemented by
places with nonzero numerical content at the prescribed scale and
their conjugates. The latter residue characteristics are among the
type primes, the auxiliary prime and primes at most \(X''\), so they are
already allowed. The inherited residue characteristics were included
with the finite list. For either part of the support let \(S\) denote its set of places. Saturation over \(K_0\) means precisely
\[
 w\cap K_0=w'\cap K_0\quad\Longrightarrow\quad(w\in S\ \Longleftrightarrow\ w'\in S).
\]

For the inherited part, two primes over the same \(K_0\)-prime
restrict to the same prime of the earlier coefficient field, and
the comparison identifies membership with the test at that prime.
For the numerical content part the test on a lift is a function
of the rational residue characteristic. It is invariant under
the same contraction. This gives the saturation on the
actual datum used to initialize the zero row orders. The fixed-field
comparison uses an enlargement of the fine datum carrying the
lower-unitary maps, with the same finite group and rational readings.
\label{cor:automorphic-realization-14}
\end{proof}

\begin{proposition}[Ray splitting and Galois representations on the initial row]\label{prop:initial-row-arithmetic-readings}
For the fixed initial row and its marked finite types, the ray
splitting test identifies the pinned ray field over the marked
quadratic base. The local type and the Galois representation attached to the same
row agree on the prescribed inertia data. For a continuous de Rham representation of rank \(a\) at a coefficient place, let \(h_{\tau,i}\) be the multiplicity of jump \(i\) at an embedding \(\tau\) inducing that place. Then
\[
 \sum_i h_{\tau,i}=a\qquad\text{for each such }\tau.
\]
\end{proposition}

\begin{proof}
\proofstep{Identify the two principal prime ideals and the pinned ray field.}
Keep the split places, the old row representation for each
finite type \(\sigma\), and the two central actions on the
ambient and multiplicity spaces. Let \(E\) be the auxiliary field,
and let \(g_0,g_1\) be the two primes used for the finite
types. Proposition~\ref{prop:two-type-primes} supplies generators
\(\pi_k\) in the integers of its distinguished quadratic subfield.
With \(Q_S=\operatorname{lcm}(2^{e+2},f_\psi)\) as in
\eqref{eq:host-character-compositum}, they satisfy
\[
   \pi_k\overline{\pi_k}=g_k,
   \qquad \pi_k\equiv 1\pmod{Q_Sqd}.
\]

Under the fixed embedding \(j_E\) of \(E\), both
\(j_E(\pi_k)\) and \((j_E(\pi_k)-1)/(Q_Sqd)\)
are integral over \(\mathbb Z\). Extend complex conjugation on \(E\) to an automorphism of the algebraic closure. The
product of \(j_E(\pi_k)\) and its image under this
automorphism is \(g_k\). This is the ray-prime test used
for splitting in the ray field of the marked base. To
see the splitting, put \(c=Q_Sqd\).

In the imaginary
quadratic base the two conjugate factors are integral
and each is one modulo \(c\). Thus \(g_k\equiv1\pmod c\)
and \(g_k\nmid c\). For the pinned ray field \(M\), the
ray-field ramification comparison gives
\[
              g\mid\operatorname{disc}(M/\mathbb Q)
                  \ \Longrightarrow\ g\mid c
              \qquad(g\text{ a rational prime}).
\]
The quadratic base has discriminant \(-d\), and \(d\)
divides \(c\). At \(g\nmid c\) the base is unramified and
the ray extension is unramified at every prime above \(g\)
by reciprocity for the modulus \((c)\). Multiplicativity of
ramification indices in the tower gives the comparison for
the absolute discriminant. Apply it with \(g=g_k\).
The two principal ideals have
norm \(g_k\) and multiply to \((g_k)\). They are distinct
because their product is \((g_k)\) and \(g_k\) is unramified.
Each has trivial ray class
and splits in the ray class field by reciprocity.
\label{cor:automorphic-realization-15}

Conversely let \(\mathfrak p\) be a degree-one prime away from \(c\)
splitting in the ray field. For the modulus \((c)\), reciprocity gives
\[
 {\rm Frob}_{\mathfrak p}=1
 \ \Longleftrightarrow\ [\mathfrak p]_{(c)}=1
 \ \Longleftrightarrow\
 \mathfrak p=(\alpha),\quad \alpha\equiv1\bmod c.
\]
In the imaginary quadratic base the norm of \(\alpha\) is the
positive rational prime below \(\mathfrak p\). These are the
generator and congruence conditions under the fixed embedding of \(E\). The ray
field, its marked conjugation and the old finite type actions are
transported with that embedding. \label{cor:automorphic-realization-16}
Indeed, if every marked ray-test prime splits completely in a finite
extension of the base, so does every such degree-one prime.

Degree-one primes suffice here: the primes of the imaginary quadratic
base of higher residue degree have Dirichlet density zero. Chebotarev's
split-prime inclusion criterion puts the extension inside the ray
field, which is therefore maximal for this test. For the embedded field \(K\) the corresponding ring is recovered from integrality alone:
\[
 \mathcal O_K=\{z\in K:f(z)=0\text{ for some monic } f\in\mathbb Z[X]\}.
\] Transport also preserves the integral row
exponents at the type primes and the valuation zero in the
lower-left comparison. \label{cor:automorphic-realization-17}
\proofstep{Use the fixed-vector comparison on the actual split local component.}
At a marked place over \(g_k\), let \(K_1\) be the first
principal congruence subgroup and \(\pi_v\) the split component
of the old row of type \(\sigma\). The split comparison embeds
\(gK_1g^{-1}\) into both bottom levels, where both level
characters are trivial. A nonzero bottom-level vector therefore
gives, by local factorization,
\[
 \pi_v^{gK_1g^{-1}}\ne0,\qquad
       \pi_v(g^{-1})\pi_v^{gK_1g^{-1}}=\pi_v^{K_1}.
\]
This uses the two level maps on the actual bottom frame.
It supplies the same depth condition for all finite types and
marked places in that row. At these places the completion
is \(\mathbb Q_{g_k}\). Depth-zero compatibility of local Langlands gives trivial wild inertia for the Weil--Deligne representation \(r_v\):
\[
 r_v(P_v)=\{1\},\qquad P_v\subset I_v\ \text{the wild inertia subgroup}.
\] Keep the local reciprocity parameter
at every split place and the Satake multiset at the spherical
inert places. \label{cor:automorphic-realization-18}

Use a tame inertia generator coupled to the fixed residue
generator at each type place. The Green labelling gives the
series idempotents in these coordinates. A local unit
representing this inertia generator can be compared with
its rational residue representative by reciprocity. Testing
the character on that representative gives the Teichm\"uller value
of the inverse finite character in the allowed orientation.

If
\(\alpha_1,\ldots,\alpha_r\) are the tame inertia eigenvalues
listed with multiplicity and \([\cdot]\) denotes Teichm\"uller lift,
the residual support polynomial is
\[
 P_\sigma(T)=\prod_{h=1}^r(T-\overline{\alpha_h}),\qquad
 [\overline{\alpha_h}]=\alpha_h .
\]
The reductions here are at the marked residue valuation, with the
chosen tame eigenvalues in the prime-to-residue-characteristic
roots of unity. Its degree is \(r\) because the parameter is tame. Transport the
place on the closure and its generator by the marked conjugation;
the residue generator law transports the polynomial with them, on the same Galois place frame and with the two original finite type actions.
\label{cor:automorphic-realization-19}
On this parameter family the character and cuspidal
comparison realizes the local parameters on the automorphic
side of the bare level.

This comparison uses local reciprocity, the
finite-dimensional fixed-level automorphic comparisons and
the character lifting at the selected ray field. In the ray
stage replace the auxiliary parameter records by those of
the split ray class field. For every chosen residue label \(u\), Teichm\"uller lifting followed by reduction at the specified ray place returns that very label:
\[
 [u]\longmapsto\overline{[u]}=u .
\]
Distinct roots in the reduction field stay distinct because
their orders are prime to the residue characteristic.
The ray modification preserves the two central actions,
the series idempotents and the marked places. Throughout
these operations the localization index set is empty and
the prescribed order data are carried along. The arithmetic
enrichment is applied before the later specialization of
the orders. \label{cor:automorphic-realization-20}

Let \(\iota_G\) identify the old and new finite groups using the
indexed type places and their matrix coordinates. The series
idempotents have matching coefficients. For either row comparison
\(J\) and central element \(z=\sum_g z_g[g]\), equivariance gives
\[
 J\rho(z)=\sum_g z_g J\rho(g)
         =\sum_g z_g\rho'(\iota_G g)J=\rho'(\iota_G z)J .
\]
The old row fits and geometric and finite isotypic readings
therefore use the same central projectors. The underlying host,
rank, empty index set, algebraic closure with its host embedding,
prime valuations and type-place indices are fixed. The places
above the real subfield still project to the specified list.
\label{cor:automorphic-realization-21}
The algebraic-closure frame extends the chosen embedding
of the host.

On the same Weil--Deligne realization impose
continuity of each \(\ell\)-adic matrix representation in
the de Rham condition. Write \(f:A\to B\) for either coefficient map. Continuity of the transported representation follows from entrywise composition:
\[
 G\xrightarrow{\ \rho\ }\mathrm{Mat}_a(A)\xrightarrow{(a_{ij})\mapsto(f(a_{ij}))}\mathrm{Mat}_a(B).
\] Thus the old
continuous local lifts still satisfy the strengthened
condition. The unramified and monodromy tests use the
original de Rham condition. \proofstep{Read the dimension of the de Rham filtration at the chosen embedding.}
Apply the Hodge--Tate
comparison to a continuous rank-\(a\) representation
that satisfies the local de Rham condition at \(w\).
For an embedding \(\tau\) into the \(\ell\)-adic
coefficient closure inducing \(w\), extend it to the
algebraic closure with the prescribed valuation ring.
One obtains the extension by conjugating a field
embedding over the host, since the two valuation
rings restrict to the same ring on the host. The continuous
representation descends to a finite extension of \(\mathbb Q_\ell\)
containing the local field. In the labelled de Rham space \(D_\tau\) this gives
\[
 \sum_i\dim\operatorname{gr}^i D_\tau
 =\sum_i(\dim\operatorname{Fil}^i D_\tau-
               \dim\operatorname{Fil}^{i+1}D_\tau)=\dim D_\tau=a .
\]
The jumps are counted with multiplicity at each \(\tau\) inducing
\(w\), equivalently \(x\in\mathfrak m_w\) if and only if
\(|\tau(x)|_\ell<1\) on host integers. \label{cor:automorphic-realization-22}
\end{proof}

\begin{proposition}[Isotypic span in the rational row]\label{prop:rational-row-isotypic-span}
At a lift \(1\le j\le m\), write \(V_{j,\lambda}\) for the space with prescribed label \(\lambda\). Its internal isotypic decomposition is
\[
 V_{j,\lambda}=\bigoplus_{\sigma\ {\rm occurring},\ {\rm label}(\sigma)=\lambda}I_{j,\sigma}.
\] In the rational
model, the allowed split Hecke operators preserve the
indicated coefficient lattices and furnish eigenvalues
on the full good simultaneous eigenspaces.
\end{proposition}

\begin{proof}
The lower-unitary multiplicity equality applies on the
same frame at every state with those pins. To see which
finite types it compares, use the good projector for the
actual two type primes. Its coefficients at \(g\) and
\(g^{-1}\) agree by the finite type inversion identity.
Inversion takes an irreducible character to the character
of its dual. In the expansion of the good projector in primitive central idempotents, inversion interchanges the labels \(\sigma\) and \(\sigma^\vee\). Independence of irreducible characters now gives
\[
 \sigma\in\operatorname{supp}(e_{\rm good})\quad\Longleftrightarrow\quad\sigma^\vee\in\operatorname{supp}(e_{\rm good}).
\] The
character comparison descends through the injective
coefficient embedding. This is the duality stability
condition in the multiplicity comparison.
\label{cor:automorphic-realization-23}

\proofstep{Compare the dimension of each label with the sum of its cuspidal isotypic dimensions.}
Let \(I_{j,\sigma}\) be the row's complex isotypic space
at a lift \(1\le j\le m\), and let \(n(\sigma)\) be
the dimension of the good finite-type space of \(\sigma\)
on the compact base row. Transport through the lift
pairing and the tensor-product identification reads
\(I_{j,\sigma}\) as the intersection of the global
\(\sigma\)-part with the fixed level, level character and
good projector of the lift. In the finite cuspidal
decomposition, the multiplicity equality identifies its
dimension with the compact counterpart. The compact
comparison identifies that dimension with the base count
\(n(\sigma)\). Hence
\[
                      \dim I_{j,\sigma}=n(\sigma).
\]
Only the fixed locus, the level character and the same
good projector enter this dimension comparison.
\label{cor:automorphic-realization-24}

The rational reading available at this point consists of
the rational coordinate isomorphisms for the sharp and
base forms, the two complex readings, their metrics at
the other embeddings, and the lift maps from rational
coordinates to the complex rows. No order comparison
is needed for these maps. The bottom Hecke action assigns
each \(\sigma\) a constituent label \(c(\sigma)\).
The theta adjoint identity on the rational generators
extends to arbitrary test vectors.

For a theta map
\(\Theta\), project a source vector to
\((\ker\Theta)^\perp\). Both sides of the adjoint identity vanish
on the discarded kernel; linearity extends it from the generators.
On a labelled image the adjoint identity gives
\[
 \Theta\Theta^*y=\lambda y,\qquad
 \lambda\|y\|^2=\|\Theta^*y\|^2>0\quad(y\ne0).
\] For the common Hecke label \(c\)
and its lift space \(V_{j,c}\) the dimension comparison gives
\[
 \bigoplus_{\sigma:\,c(\sigma)=c}I_{j,\sigma}\subset V_{j,c},
 \qquad \dim V_{j,c}=\sum_{\sigma:\,c(\sigma)=c}n(\sigma).
\]
Distinct cuspidal projectors give the direct sum, and
\(\dim I_{j,\sigma}=n(\sigma)\). Thus the inclusion is an equality
at every lift. \label{cor:automorphic-realization-25}

\proofstep{Read the operators and eigenvalues on the same integral coefficient model.}
At the arithmetic level these decompositions give the
span test for the good simultaneous Hecke eigenspaces.
The rational forms come from the full geometric lattices
and are read through \(K_s\hookrightarrow K_0\).
At each type prime \(g_k\), the inverse of the
chosen primitive \(g_k\)-th root of unity defines
a nontrivial additive character of \(\mathbb F_{g_k}\).
Compose it with the marked linear combination of the
superdiagonal entries of the upper unipotent matrices.

Write \(L(u)\) for the marked linear combination of the superdiagonal entries and \(\zeta\) for the chosen primitive root. The level character is the inverse unipotent character because
\[
 (\zeta^{-1})^{L(u)}=(\zeta^{L(u)})^{-1}.
\] The split
Hecke transform and its rational realization give the operators
on the same rational forms, with the index-zero operator
equal to the identity. The integral comparison gives
stability for the indicated coefficient lattices.
Use the isotypic span and the good finite-type condition
to choose their eigenvalue arrays. The finite set of
excluded places contains the indicated ramified and level
places. These rational operators act on the same integral lattices
whose characters and Hecke span were just determined. \label{cor:automorphic-realization-26}
\end{proof}

\paragraph{Theta series and marked points on the rational row.}
Proposition~\ref{prop:rational-row-isotypic-span} supplies the
integral operators and their simultaneous eigenvalues.
The cuspidal parameter comparison uses the same good lower
types that occur in the rational row. At split primes outside
its finite exclusions, the rational arrays read the elementary
symmetric polynomials in the normalized Satake parameters.
The index-zero identity fixes the normalization of the
Hecke operators. The matching of Frobenius polynomials on
these split primes, with semisimplicity and Chebotarev, gives
the coefficient-field condition for the cuspidal polynomials
used in the local multiplicity models. The comparison also
uses the fixed lower-left valuation and the reciprocal
character twist in passing back to the original row.
\label{cor:automorphic-realization-27}


\begin{proposition}[A common grid for the coefficient and geometric models]
At the prescribed grid size, depth and finite list of levels, the cuspidal forms, toric coefficients, and marked-point evaluations can be chosen with common grid characters, excluded places and signed indices. If \(J\) identifies the original row with its twisted copy in the identity multiplicity frame, every Hecke operator \(T\) outside the common exclusions satisfies
\[
 T^{\rm new}J=JT^{\rm old}.
\]
Use the multiplicity comparison of Proposition~\ref{prop:linked-type-multiplicity}
and the coefficient-ring and marked-point identifications for this
grid. \label{cor:automorphic-realization-28}
\end{proposition}
\begin{proof}
Take a grid on this row at the prescribed size and depth.
The cusp multiplicity construction and the class-group
coefficient-ring construction each specify a finite
coefficient extension at every indicated residue prime and
lift. Use their compositum. Choose one finite set \(S_x\)
of places of \(E\) containing the exclusions of the rational
Hecke action and the allowed conductors of the grid.
The twisted coefficient rings can then be chosen with
excluded place sets containing \(S_x\). The two coefficient
extensions compare in these very rings. In particular the
coefficient matrix of the cuspidal row and the finite-order
toric character can both be read there. This common choice
is needed because the toric construction compares at the
same grid entries and away from the same exclusions.
\label{cor:automorphic-realization-29}

For the cyclotomic coefficient valuation ring \(O_v\), the
integrality test in the signed comparison is
\[
 a/b\in O_v\quad\Longleftrightarrow\quad v(a)\ge v(b)
       \qquad(a,b\in O_v,\ b\ne0).
\]

The toric-character construction takes the
grid characters with their Teichm\"uller lifts and
their values in the fixed field. It uses the lower-left
unit and the inclusion of their conductor in the excluded
set. The toric ratio comparison identifies the coefficient
reading of these characters. Apply the multiplicity
comparison both at the trivial character and at the
selected toric characters, using the same full
coefficient-module basis. The settings obtained this way
carry the rational Hecke reading and the complete grid.
\label{cor:automorphic-realization-30}

Let \(J\) identify the old row with its twisted copy in the
identity multiplicity frame. The exclusion sets at corresponding
places agree, so for an allowed Hecke operator \(H\) one has
\[
 H^{\rm new}J=JH^{\rm old},\qquad
 H^{\rm old}v=\lambda v\ \Longrightarrow\
             H^{\rm new}(Jv)=\lambda Jv .
\]
The excluded residue characteristics obey their previous bound.
The level identification transports the unipotent character and
the grid with the coefficient forms. The series comparison on
these same forms supplies the designated lift support and the
signed level condition. \label{cor:automorphic-realization-31}

For the geometric step from the series, the marked
points detect the coefficient algebra in the framed
space. The signed comparison descends the support
condition to the geometric model. Let \(k_H(x)\) count the allowed kernels of a split Hecke correspondence \(H\) at a framed point \(x\). Outside the excluded characteristics the local level gives a single count: 
\[
 k_H(x)=k_H^{\rm level}=k_H(x')\qquad(x,x'\text{ in the prescribed frame}).
\]
 Its quotient read over the
complete field is integral in the chosen geometric
coordinates. The integral correspondence and the series support outside the selected range give the marked moduli-point model. Transport of level preserves the grid characters, their field values and their roots. \label{cor:automorphic-realization-32}

In particular the character assignment to coefficient places uses an index map \(v\mapsto i(v)\) satisfying
\[
 i(\bar v)=i(v)\qquad\text{at each assigned place }v .
\] Apply reciprocity to its finite-order Hecke
character, valued in the roots of unity of the specified
prime-power order. Let \(C_v\) be the set of conductor places of the character chosen
at \(v\), and \(B_{i(v)}\) the exclusions for its index. The two
support conditions are
\[
 C_v\cap B_{i(v)}=\varnothing,\qquad
 C_v=\varnothing\quad(v\text{ outside the twist support}).
\] Use the
integral ring-class descent at the two readings to put
the twist into the coefficient model. The rank test
for each reading is the level comparison at its actual
signed place. Equality of the characters gives equality
of their ramification supports, so the two conjugate
choices agree after the signed transport. The construction
at a fixed finite set of levels makes these choices for
the whole list of prescribed indices. Its spectral
specialization must keep the chosen twist labels while
inserting the actual row orders.
\label{cor:automorphic-realization-33}

\end{proof}

\begin{proposition}[Localization of the integral Hecke algebra]\label{prop:marked-hecke-localization}
At a prescribed coefficient place, tuple, and lift, a finite
extension of the local coefficient field splits the occurring
full Hecke eigensystems and contains any prescribed finite list
of scalar values. Over its integers \(O\), the Hecke order \(\mathbb T\) decomposes into residual-character blocks. If \(V_\sigma\) is a full eigenspace with character \(\lambda_\sigma\), the summand with residual character \(\bar\lambda\) is described by
\[
 \mathbb T=\prod_{\bar\lambda}\mathbb T_{\bar\lambda},\qquad V_{\bar\lambda}=\bigoplus_{\sigma:\,\overline{\lambda_\sigma}=\bar\lambda} V_\sigma.
\] The construction preserves the original
integral action under restriction to the block.
\end{proposition}

\begin{proof}
\proofstep{Split the commuting spherical operators over one finite extension.}
Let \(v\) be a marked coefficient place of residue characteristic
\(\ell\), and keep its comparison \(\iota_v\colon
\overline{\mathbb Q}_\ell\simeq\mathbb C\). The strict-henselian
character choice fixes the tuple, level and conductor on the definite
multiplicity spaces. For each lift \(j\), use the existing rational
model of the corresponding space and its integral lattice. At the
allowed split primes outside the conductor, \(\ell\) and the other
exceptional places, the normalized spherical double cosets act by
rational operators preserving that lattice. On each complex isotypic summand \(V_\sigma\) of the definite decomposition, the hyperspecial fixed line of the local spherical representation has a Hecke character \(\lambda_\sigma\). The spherical operator \(T\) acts on each copy through this same character, so
\[
 T|_{V_\sigma}=\lambda_\sigma(T)\,\mathrm{id}_{V_\sigma}.
\] Equivalently, the convolution adjoint is obtained
by inversion and conjugation, so the commuting spherical operators
are simultaneously diagonalizable for the invariant pairing.

A basis for the rational Hecke algebra involves only finitely many
of the allowed double cosets. Let \(T_1,\ldots,T_s\) be generators
so obtained, and choose \(F_{j,v}/\mathbb Q_\ell\) finite containing
the marked coefficient field and their eigenvalues. The simultaneous
eigenspace for a character \(\lambda\) is
\[
                 \bigcap_{a=1}^s\ker(T_a-\lambda(T_a)I) .
\]
The equations are defined over \(F_{j,v}\), so they split the entire
action there. Take the compositum
with the fields for the local parameter comparisons and the finite
character values. Adjoin the requested complex coefficients in the specified embedding.

For a finite set \(S_v\subset\mathbb C\) this amounts to the extension
\[
 F_{j,v}\ \subseteq\ F_{j,v}(\iota_v^{-1}(S_v)).
\] All local factors over \((j,v)\)
thus have the required coefficient inclusion with the original
embedding. The chosen finite character takes values in units, and
its auxiliary factor is trivial on inertia at the conductor primes
where ramification of the marked factor is required. The product
retains those ramified primes. Its conductor is away from \(\ell\)
and the excluded level primes by the character choice.
\label{cor:automorphic-realization-34}

\proofstep{Identify maximal ideals with reductions of full characters.}
Put \(O=\mathcal O_{F_{j,v}}\). On the base-changed lattice the integral
algebra \(A\) acting faithfully is finite over \(O\), as a submodule
of the integral endomorphism ring. If the space is zero, the index
list is empty. Otherwise
\[
 A\otimes_O F_{j,v}=\prod_{\alpha}F_{j,v},
 \qquad A\subset\prod_{\alpha} O ,
\]
where \(\alpha\) runs through the occurring full characters.
Each eigenvalue belongs to \(O\) by the monic characteristic
polynomial of its integral operator. Write \(k=O/\mathfrak m_O\).
Since \(\prod_\alpha O\) is finite over \(A\), lying over gives
\[
 \operatorname{Max}A
   =\{\ker(A\xrightarrow{\alpha}O\longrightarrow k):\alpha\}.
\]
Two residual characters with this kernel induce the same map
to \(k\), since both fix \(k\) on diagonal scalars. Completeness
lifts the residual factors to orthogonal idempotents \(e_{\mathfrak m}\).
Thus for the original lattice \(L\),
\[
 (e_{\mathfrak m}L)\otimes_O F_{j,v}
   =\bigoplus_{\ker\bar\alpha=\mathfrak m} (L\otimes_O F_{j,v})[\alpha].
\]
The decomposition intertwines every integral Hecke operator.

\proofstep{Recover the complex isotypic summands from the same character.}
Write \(v=\sum_\sigma v_\sigma\) in the complex isotypic
decomposition. For an allowed generator \(H\),
\[
 (H-\alpha(H))v=0
 \quad\Longleftrightarrow\quad
   (\alpha_\sigma(H)-\alpha(H))v_\sigma=0\quad\text{for every }\sigma .
\]
Intersecting these conditions over \(H\) gives precisely the
summands with full character \(\alpha\). The hyperspecial factor
has trivial level character at an allowed split prime. A nonzero
summand there is spherical, so its Hecke polynomial reads the
unramified parameter with the rank-\(a\) twist
\(\lvert\det\rvert^{(1-a)/2}\). The tuple, lift and multiplicities
are common across the coefficient factors. At the two marked
type primes away from \(\ell\), only the prescribed residue
character acts on inertia, with its inverse at the conjugate
place. The Weil--Deligne and infinity comparisons use the
designated embeddings. \label{cor:automorphic-realization-35}

The block lists are determined by the normalized split polynomials
on the constituents used here. For a block of length \(d\) centred on a conjugate self-dual cusp
\(\pi\) of rank \(a\), its terms, with multiplicity, are
\[
           \pi |\det|^{(d-1)/2-k},\qquad 0\le k<d .
\] In total rank \(r\) the common
Galois normalization adds \((1-r)/2\). Relative to the normalization
in rank \(a\), the requisite algebraic shift is \(u+(a-r)/2\).
The allowed infinity parameters are regular with the block's
parity, so the shifted cuspidal Galois realizations in this comparison
are supplied by the regular algebraic cuspidal input. The norm twist
changes the polarizing character by a norm power. For a regular
algebraic polarized cusp over a CM field the attachment is
continuous, with the normalized local polynomial at each unramified
place away from the residue characteristic. This is the part of
Barnet-Lamb, Gee, Geraghty and Taylor's attachment theorem used here.

To pass from the split Frobenius elements to the Galois group, compare each coefficient of the two direct-sum polynomials. The difference is continuous, so
\[
 \{g:\det(1-X\rho_1(g))=\det(1-X\rho_2(g))\}
\]
is closed and contains a dense set. Here \(\rho_1,\rho_2\) denote the two direct sums in the comparison. In particular there is equality of the normalized
polynomials at all but finitely many primes. Cancel the common
normalization and apply Jacquet and Shalika's uniqueness for the
isobaric sums with their real norm twists. Unitarity and the
central characters distinguish the real exponent of each
cuspidal term. Fix one cusp; let \(e\) be the largest exponent
left and \(m(e)\) its multiplicity. Its longest remaining blocks
have length \(2e+1\), with one occurrence of \(e\) per block.
Update all exponent multiplicities by
\[
 m_{\rm new}(s)=m(s)-m(e)\,\mathbf 1_{\{e,e-1,\ldots,-e\}}(s).
\]
Induction on the number of terms recovers even the blocks with
overlapping strings. \cite{blggt-attachment,jacquet-shalika-ii}\label{cor:automorphic-realization-36}

For each row type and residue place which has an integral Galois
realization, reduce one such realization and take the direct sum of
its composition factors. The kernel is open since it contains
the open kernel of the reduction. Any two admissible realizations
of this type at the same place read the same split polynomials
through the common coefficient embedding. Reduce the coefficient
identities in that marked valuation ring outside the finite union
of their bad primes.

For two such residual representations \(\bar\rho_1,\bar\rho_2\),
the coefficient functions are locally constant. Equality on the dense
split Frobenius conjugates therefore gives, for every Galois element \(g\),
\[
 \det(X-\bar\rho_1(g))=\det(X-\bar\rho_2(g)).
\] The polynomial is unchanged by
semisimplification, by the block triangular determinant formula.
If there is no such realization the comparison has an empty source
and the auxiliary residual representation may be taken trivial.
This supplies the semisimple data for arbitrary prescribed
\(S_v\). The comparison is coefficientwise also at the small residue
characteristics. Its assertion about the polynomials is independent
of the degree of the coefficient enlargement.
\label{cor:automorphic-realization-37}
\end{proof}

\begin{proposition}[Simultaneous coefficient requests and local row orders]\label{prop:joint-row-orders}
At each prescribed tuple and lift, finitely many conditions
on the compatible semisimple readings may be imposed at once
by adjoining the union of their coefficient requests. Write \(l_\alpha\) for the ordinary \(O\)-lengths of the indicated localizations. On these specializations
\[
 \eta_j(v)=\frac{\sum_\alpha l_\alpha}{\operatorname{length}_O(O/\ell O)};
\]
in the untwisted case all the lengths and the order vanish. The comparison uses the same good
projector and is unchanged by enlarging the coefficient DVR
at the given place.
\end{proposition}

\begin{proof}
\proofstep{Meet the field requests by union before applying specialization.}
Use the localized order of Proposition~\ref{prop:marked-hecke-localization}.
Every coefficient condition here has the form of a finite set
\(S_v\) at each original place \(v\), followed by an assertion
about every compatible semisimple realization whose marked embeddings
contain \(S_v\). For the two conditions with requests \(S_v^{(1)}\) and \(S_v^{(2)}\), impose the single request
\[
 S_v=S_v^{(1)}\cup S_v^{(2)}\qquad\text{at each original place }v .
\]

The coefficient construction just given works for that union and
keeps the conductor, level and tuple unchanged. For a compatible specialization the two implications can therefore be applied at the same homomorphism. For an occurring eigencharacter, choose a nonzero eigenvector \(u\) in the geometric row. The resulting specialization is the homomorphism
\[
    \mathbb T\longrightarrow O,\qquad T\longmapsto\lambda_u(T),
                         \qquad Tu=\lambda_u(T)u .
\]
Here the row is the one constructed in Lemma~\ref{lem:constructed-signed-coefficient-row}. Its integrality follows from the stable ball of Lemma~\ref{lem:marked-semistable-local-ball}. At a new ray type, occurrence on the tested Hecke algebra is proved in Lemma~\ref{lem:new-ray-type-occurrence}. In particular the order identity,
which is proved on specializations satisfying the full coefficient
ring conditions, can be imposed together with those conditions.

At the fixed tuple and weight only finitely many systems of split
Hecke polynomials occur. Enlarge their exceptional set to contain
the conductors and excluded level places. For each realized system
fix one continuous integral representative
\(\rho:G_E\to{\rm GL}_a(\overline{\mathbb Q}_\ell)\).
Writing the algebraic closure as a countable union of finite
extensions \(F_b\), one has
\[
 G_E=\bigcup_b\rho^{-1}({\rm GL}_a(F_b)).
\]
These are closed subgroups of the compact group \(G_E\); by Baire
one is open. Adjoin the matrix entries of finitely many coset
representatives to obtain a finite field of definition \(F_\rho\).
There are finitely many extensions of \(F_\rho\) of degree
\(\le a\) in the chosen closure. Their compositum contains every
root of \(\det(T-\rho(g))\). By split Frobenius density any other
realization of the system has the same characteristic polynomials.
Request generators of these finitely many composita. An unrealized
system imposes no condition.

To impose the integral block condition, first specialize with the
rational eigenvalues, split polynomials and decomposition of the
good locus. Every occurring good block lies over a good residual
block here by descent of residual irreducibility. For each such
block use its separating Galois element and the matrices at its
finitely many characteristic-zero points. At each original \(v\)
make the \emph{single} request
\[
 S_v=S_v^{\rm split}\ \cup\
    \bigcup_{\text{these blocks and points at }v}
       \{\text{roots of the chosen characteristic polynomials}\}.
\]
Integral base change carries those matrices into the final blocks.
Thus all the root tests and order comparisons hold on the one
specialization obtained after the union. \label{cor:automorphic-realization-38}

For the row orders, fix a lift \(j\) and a place \(v\) with
nonreference tuple. At this tuple the rational definite space
carries an integral lattice \(L\). Let \(p_G\) denote its rational good projector.
Put \(G=\operatorname{im}p_G\) for the good space and \(B=\ker p_G\)
for its complement in the ambient row. Use the matched theta reading,
so \(B\) is the kernel of its adjoint. The tuple identification,
twisted vacuum identity and rational class-sum comparison give the
following pair of geometric lattices in this row:
\[
 L_{\rm upper}=p_G L=G\cap(L+B),\qquad L_{\rm lower}=L\cap G .
\]
For the middle equality write \(a=p_Ga+(1-p_G)a\) for \(a\in L\).
Conversely a vector of \(G\) in \(L+B\) projects to itself.
The lower identity follows by pulling \(L\) back along the injective
good-space embedding. These identifications give the same lattice
quotient in the two realizations. \label{cor:automorphic-realization-39}

The good projector in the localized definite model is the sum of
the good point projectors. For the finite type check on the
original row the supported lifts at a marked split factor are
\((g,(g^{-1})^{\mathsf t})\) in the paired integral coordinates.
The principal congruence kernels fix the forms at both marked
row levels. If \(e=\sum_g c_g[g]\) is the finite idempotent, choose a lift \(\tilde g\) of each element by surjectivity.

On the representation factoring through reduction one obtains
\[
 \sum_g c_g\,\rho(\tilde g)=\sum_g c_g\,\bar\rho(g)=\bar\rho(e).
\] A nonzero vector in the good part of an isotypic
summand retains that idempotent after projection onto an
irreducible copy in which it has a nonzero component. One full
good product of finite series projectors therefore acts
nontrivially in its local components. The tame type dictionary
identifies the degree-\(r\) support polynomials with their
inertia polynomials including multiplicity. Conjugation gives
the dual series, and the good set is dual stable. This gives
the original good-label test carried into the localization.
\label{cor:automorphic-realization-40}

For the matched theta range, fix a localized index and denote the
upward theta map by \(T\). On every nonzero lower isotypic summand the adjoint identity reads
\[
             T^*T=\lambda_\sigma I,\qquad \lambda_\sigma>0.
\]
At a split prime of norm \(q\), remove the common normalization.
The lifted root multiset is the smaller one together with the
specified string of length \(t\). Its largest modulus is
\(q^{(t-1)/2}\), attained exactly once. For a block of length \(d<t\)
with cuspidal root \(\alpha\),
\[
 |\alpha|q^{(d-1)/2}<q^{d/2}\le q^{(t-1)/2}.
\]

If a block of length \(d\ge t\) has centre of rank \(a\), its rank
alone consumes \(ad\) dimensions. Thus
\[
                ad\le r+t<2t\le2d,\qquad a=1 .
\] That center is unitary, so its largest modulus forces
\(d=t\), and uniqueness of the top root forces a single such
block. Dividing that root by \(q^{(t-1)/2}\) recovers its center.
Reciprocity and split Frobenius density identify the global
character. Remove one copy of this block, with multiplicities.
At a good point \(x\),
let \(U_x\) be the sum of the lower summands with the corresponding
parameters, and \(V_x\) the point eigenspace. The parameter
comparison gives \(TU_x\subset V_x\). Any other lower summand lifts orthogonally to \(V_x\).

Indeed, writing \(C\) for the common tuple block and \(L,L'\) for the lower parameter multisets, cancellation gives
\[
 L\sqcup C=L'\sqcup C\quad\Longrightarrow\quad L=L' .
\] Consequently \(T^*V_x\)
is contained in \(U_x\). The downward nonvanishing condition
is injectivity of \(T^*\) on the whole \(V_x\) for a good point.
Section~\ref{subsec:selected-point-adjoint} gives injectivity of
\(T^*|_{V_x}\) before the row orders are inserted, for the specified
rank-one string \(t\ge r+5\), good finite type and marked moduli
model. For \(v\in V_x\) one has
\[
 v\perp TU_x\ \Longrightarrow\
 T^*v\in U_x\cap U_x^\perp=0\ \Longrightarrow\ v=0.
\]
The inclusion \(TU_x\subset V_x\) is therefore an equality.
The class-sum identity and the good-type incidence rule now
identify the sum of these ranges with the range of the good
projector. The remaining simultaneous eigenspaces give its
complement. \label{cor:automorphic-realization-41}

Localize and extend scalars to the marked coefficient DVR.
The integral double cosets preserve \(L\), with the same conductor
in both its rational and spectral realizations. It decomposes
integrally as \(\bigoplus_\alpha L_\alpha\) over the occurring
residual blocks, and \(G=\bigoplus_\alpha G_\alpha\),
\(B=\bigoplus_\alpha B_\alpha\) after scalar extension.
The two lattices in the quotient
\[
          \frac{G\cap(L+B)}{G\cap L}
\]
therefore split by the same residual idempotents. Write \(Q_\alpha\) for the resulting localized quotient, and
\(O_v\) for the coefficient DVR at \(v\mid\ell\). The normalized total
order of this quotient is
\[
              \frac{\sum_\alpha\operatorname{length}_{O_v} Q_\alpha}
                   {\operatorname{length}_{O_v}(O_v/\ell O_v)} .
\]
\proofstep{Normalize length by the original coefficient valuation.} This convention is invariant under extension of DVRs
at the same place. Flatness carries the intersections and images
to the extended comparison, and a composition factor over the
old ring becomes a module of length the ramification index.
For an extension \(O_v\to O_w\) of ramification index \(e\), this
is the cancellation of Lemma~\ref{lem:retained-dvr-normalized-length}:
\[
 \frac{e\operatorname{len}_{O_v}(M)}{e v_v(\ell)}
          =\frac{\operatorname{len}_{O_v}(M)}{v_v(\ell)}
\]
for the index quotient \(M\). At an untwisted place \(v\) there are no localization terms, and for every lift one obtains
\[
 \sum_{\text{localizations at }v}\operatorname{ord}(M)=0=\eta_j(v).
\] This proves
the identity with the prescribed \(\eta_j(v)\) at every place
and every \(1\le j\le m\), with no factor from the local degrees
or the number of localizations. The equality uses the good
projector and integral comparison on the same chosen tuple.
\label{cor:automorphic-realization-42}

\proofstep{Carry the order to the level-three row specialization.}
Apply the archimedean and type comparison to the localization that
meets the united request. It uses the Lie and Galois comparisons
on its compatible local parameter data and obtains the normalized
character, local type and level data on that specialization.
The row-lift family there is read in the same finite group and
marked adelic frame. Write \(o_j(v)\) for the integral localization order on this row. The input to the normalized family is the pointwise equality
\[
 o_j(v)=\eta_j(v)\qquad(1\le j\le m,\ v\text{ a coefficient place}).
\] Transporting the row to the normalized theta family preserves
the original weight functions and character. The grid list transports by the
preserved marked readings at each allowed depth. The initial
pass and the rebuilt pass apply this argument to their respective
strict-henselian inputs. Their resulting row specializations are identified by the same
row and character comparisons.
\label{cor:automorphic-realization-43}
\end{proof}

\begin{lemma}[Type of the universal theta series]\label{lem:universal-theta-type}
Fix an index, a lift, and a sign in the prescribed row.
Let \(M_0\subset M\) be the submodule with the marked level
type. If \(\Theta=\sum_{g\in\Delta}a_g[g]\) is the universal theta series in \(M\otimes_O O[\Delta]\), then
\[
 a_g\in M_0\quad(g\in\Delta),\qquad \Theta_{\rm ordinary}\in M_0.
\] After an admitted specialization the
series lies in the scalar-extended type submodule.
\end{lemma}

\begin{proof}
Fix an index, a lift \(1\le j\le m\), and a sign.
Write \(M\) for the coefficient module of this row
and \(M_0\) for the submodule satisfying the marked level
type. Writing \(T_h\) for a level operator and \(A_h\) for its prescribed
type action in these coordinates, the defining equations are
\[
                    M_0=\bigcap_h\ker(T_h-A_h)\subseteq M.
\] The injection into the complex model intertwines these
operators, so their kernel conditions may be tested there. There the theta transformation law
gives the marked type. The same law holds for each admitted
twisted generator and for every admitted universal coefficient
viewed in the original row. The real comparison uses the
uniform lift reading and the identical complex models
of the actual row. \label{cor:automorphic-realization-44}

\proofstep{Extract the coefficient for each element of the finite twist group.}
The universal series lies in \(M\otimes_O O[\Delta]\)
for the finite twist group \(\Delta\) and the coefficient
ring \(O\) at this index. On the group-algebra basis the extraction of one coefficient is
\[
  (1\otimes[\delta]^*)\!
    \left(\sum_{d\in\Delta}m_d\otimes[d]\right)=m_\delta .
\] In the admitted case its
class-sum formula is a finite sum over terms of that label.
Each term belongs to the prescribed type submodule \(M_0\) of \(M\), so their sum is in
\(M_0\) as well. Off the specified range the universal
series is zero.

Since the coefficient maps recover
every coefficient of any expansion in the group basis,
this proves the type assertion for every expansion
and every label. The ordinary theta series itself lies
in \(M_0\) by the same faithful type test. Thus every
datum in the designated lift range may be used for the
base equivariance of the theta data. For an admitted
specialization the corresponding type assertion holds
in the scalar-extended type submodule.
\label{cor:automorphic-realization-45}
\end{proof}

\begin{established}[The real row and measured columns at the prescribed scale]\label{hyp:real-row-measured-models}
Fix an arithmetic row specialization of Proposition~\ref{prop:joint-row-orders}. Let \(F_u^0\) be
its degree-zero column form model for the word \(u\), including
the chosen column values, the level at \(T_u\) and the two
type markings. If \(V_{\rm good}\subset V_{\rm full}\) are the two rows recorded here, the dimension conditions used in the replacements are
\[
 \dim F_u^0>0\quad\text{for each word }u,\qquad\dim V_{\rm good}\le\dim V_{\rm full}.
\]
The local matrix models, Petersson measure and quotient measures on the flag fibres can all be computed in the same Hermitian coordinates. In these coordinates the unit lowest vectors satisfy the absolute Bessel bounds used for the finite term expansion.
\end{established}
\begin{proof}
The dimension formula of Proposition~\ref{prop:marked-column-dimension} is a product of positive volume, weight and finite marking factors. The other dimension inequality follows from inclusion of vector spaces. Keep this row and these columns while making the following two
replacements of the spectral data. In the first replacement use
the scale of the actual Haar data of the form, the discrete-series
parameters obtained from the labelled Harish--Chandra weights
and signatures, and the numerical discriminant of the fixed
scale context.

For each finite place the local matrix model is
obtained from the completion of the marked quadratic algebra
with its involution. The adjoint form in the local matrix model comes from scalar extension of the same Hermitian form. For a Gram matrix \(H\) in the row or column coordinates, the adjoint is read locally by
\[
 A\longmapsto H^{-1}\bar A^{\,t}H ,
\] This gives the placewise
isomorphisms for the replaced spectral data, with the same
exponent and size context. Division lattices, their height
functions, the chosen character, and the localization grid
are transported by the row identifications. The transport for
the initial state carries the origin comparison, and that for
the rebuilt state carries the origin and geometric norm
comparisons.
\label{cor:real-row-realization-1}

Let \(X\) be the adelic unitary model of this Hermitian space.
The measures in \(X\) are chosen together with the standard
additive character. Start with the self-dual lattice from
the reference form \(F_{u_*}^0\). At the marked split primes keep the fixed frames. For each finite place \(v\), take the self-dual Haar measure on the additive unipotent coordinates for the standard character \(\psi_v\). In the Fourier convention this means
\[
 \widehat{\widehat f}(x)=f(-x)\qquad(f\text{ a Schwartz--Bruhat function}).
\]

At the real places use the prescribed unit
lowest vectors and the Bessel functions of the discrete
series of the selected weight. The local choice is made
for the weight gaps of the scale context. The stabilizer
quotient measures come from the Gram fibres. Use the
convergent unitary Siegel--Weil identity for the two
stabilizers and the rational quotient density to normalize
their covolumes simultaneously. In this normalization the
Petersson measure and the product measures on the flag
fibres refer to the same ambient Haar, as in the fibre
comparison of Section~\ref{sec:absolute}. The gap bound
gives the weighted absolute integrability of the Bessel
functions through the matrix-coefficient estimate of Proposition~\ref{prop:bessel-sign-gap-decay} and the local Bessel projection in Section~\ref{sec:absolute}. These are the measure and decay
conditions used in the finite term expansion.
\label{cor:real-row-realization-2}
\end{proof}

\begin{proposition}[Identification of the measured columns and their fibres]\label{prop:measured-columns}
For the Hermitian column models and measure of Lemma~\ref{hyp:real-row-measured-models}, match the ordered
frames, the distinguished line of norm \(-T_u\), and the characters
at both type primes. The resulting finite adelic transport induces an
isometry of the degree-zero column spaces. For a nonempty term fibre \(\mathcal F\), write \(\mathcal F/U\) for its lift classes modulo flag unipotents, \(S\) for the marked target stabilizer and \(X_0\) for the target isotropic space. There is an injection
\[
 \mathcal F/U\lhook\joinrel\longrightarrow\operatorname{im}\bigl(S\to\operatorname{Aut}(X_0^\perp/X_0)\bigr),
\]
with finite target. Corank and the existence of a scalar lift are
independent of the representatives of a pair with equal row reading.
\end{proposition}

\begin{proof}
\proofstep{Matching the local lattices.}
For each sequence, two ordered frames have the same Gram matrix.
Their hyperbolic pairs have pairing one, and the distinguished
line has norm \(-T_u\). If \(A\) is the matching between their spans and \(G,G'\)
are their Gram matrices in any bases, the Gram equalities give
\(A^*G'A=G\). This extends by Witt's theorem. In the
orthogonal complement the two degree-zero lattices are
self-dual off \(T_u\). At an unramified inert place split off a unit-norm line at each step. If a chosen vector has unit norm \(a\), surjectivity of the norm on units gives a unit \(c\) with
\[
 N(c)=a^{-1},\qquad \langle cv,cv\rangle=N(c)a=1 .
\] At a split place the two halves of a self-dual
lattice are paired duals. At \(T_u\) the dual quotient on
each split factor has length one. Indeed the primitive
distinguished line splits rationally from its orthogonal
complement in a self-dual middle lattice.

Projection gives
the two dual quotients the same length, and the norm
\(-T_u\) gives that length at the split factors. In split
coordinates the resulting lattice chain is a full lattice
and the inverse image of a hyperplane modulo a uniformizer.
A change of basis \(A\) on one factor forces the inverse transpose on the dual factor, since for paired coordinate vectors
\[
 (Ax)^t(A^{-t}y)=x^ty .
\] At the two type
primes a diagonal unit correction in the common finite
reduction matches the superdiagonal characters in the
chosen frames. Conjugation by the inverse of
\(\operatorname{diag}(1,b,b^2,\ldots)\) multiplies all
first-superdiagonal entries by \(b\). In reversed dual
coordinates inverse transpose negates their sum. The
primitive residue characters are related by a nonzero
residue scalar, so one chooses \(b\) to match the inverse
character in the sharp reading and lifts its entries to units.

\proofstep{Transport of the class sets.}
These local comparisons are identity at almost all places.
They give a finite adelic transport identifying the
column levels and compact models. For the induced bijection
\(a\mapsto a'\) on double quotients, let \(\mathcal I_a\) be the
unitary value comparison in the compact model at \(a\). It sends a
class function \(F\) to \(F'\) with \(F'(a')=\mathcal I_a F(a)\).
Covariance extends this identity to all representatives, and
\[
          \sum_{a'}\|F'(a')\|^2=\sum_a\|F(a)\|^2
\]
uses the bijection and the finer-level trivial stabilizers. The chosen unit vectors in the compact
highest-weight representations identify the invariant
norms. At the finer level the stabilizers are trivial by
the split type-prime argument of
Section~\ref{subsec:selected-word-level}. Indeed the class bijection preserves the values in the identified compact representations. With trivial stabilizers the norm comparison is simply reindexing: for a column function \(f\) and class bijection \(b\),
\[
 \sum_C\|f(C)\|^2=\sum_{C'}\|f(b^{-1}C')\|^2 .
\] It identifies the actual column forms
entering the Bessel and Gram computations.

The finite
term fibres use these very classes and their shorter
and complete flag representatives. Fix one lift \(s_0\) of a term in a nonempty fibre. Write
\(X_0\) for the isotropic space of its marked target flag and
\(\Lambda'\) for its target lattice. Modulo the two flag unipotents the transition for any other lift is \(s_0 c\). If \(b_i\) are the marked flag basis vectors and \(v\) the marked vector, then \(c\) satisfies
\[
 c\Lambda'=\Lambda',\qquad c b_i=b_i\ (\text{each }i),\qquad c v=v .
\] Choose one such \(c\) per lift. Its
action on \(X_0^\perp/X_0\) lies in the finite image of the marked
lattice stabilizer. Here the definite lattice automorphism bound
applies with the signature at stage zero and the positive norm
parameter \(T_{u'}\) of the target word. Write \(\bar c\) for the induced endomorphism of
\(X_0^\perp/X_0\). If \(\bar c=\bar c'\), put \(u=c^{-1}c'\).

The three conditions defining the flag-unipotent equivalence are
\[
 u|_{X_0}=1,\qquad u\Lambda'=\Lambda',\qquad
 (u-1)X_0^\perp\subseteq X_0 .
\]
The last inclusion is just \(\bar c^{-1}\bar c'=1\); the first
two follow from the marking and the integral lattice. Such \(c\) also normalizes
the target flag unipotents, so the two lifts represent the same
class. The chosen map into the finite image is injective.

\proofstep{Independence of representatives.}
Consider also a pair of lifts having the same row reading, and a
second pair with the same respective terms and again equal row
readings. In the transport just used, equality of the two marked
embedding readings identifies the induced Levi actions. The neat
level makes their rational stabilizer representatives equal. In particular the relative target transporter induces the identity on the flag quotient:
\[
 (c_2c_1^{-1})v-v\in X_0\qquad(v\in X_0^\perp).
\] Conjugate it by the first lift. The new relative
matrix differs from the old one by this conjugate and the
lattice-preserving source and frame unipotents. They carry the
space defining the scalar \(s_0(\cdot,\cdot)\) isomorphically
onto the new space, preserving its codimension, and carry any
scalar lift to a scalar lift. Interchange the two pairs to get
the reverse implication. This proves representative independence
before using absolute convergence to reindex the expansion.
\cite{companion-row-count}\label{cor:real-row-realization-3}
\end{proof}

The product discrete series exists with the prescribed
labelled real parameters. Its formal degrees are computed
in the actual product Haar. In the real tie their product
is the one at the true scale, and the finite factor is the
value of the actual counting projector at the identity.
The comparison of quotient volume with mass unfolds the
same product measures and uses the definite genus
decomposition. Hence the identity contribution in the
measured model is exactly the predicted \(D_*\) of
Section~\ref{sec:count}. In particular it uses the
finite group quotient and level index of this row. The
finite good operator is fixed under the conjugate reading
in the matched frames. The local series projector has
its coefficients paired by dualizing the supported tame
labels and reversing their order. This is the coefficient
identity which compares the two readings of the operator.
\label{cor:real-row-realization-4}

\begin{proposition}[The isometry of the refined row]\label{prop:refined-real-row}
Let \(J\) be the sharp geometric level, \(n_P\) its level index
in the normalization of the Petersson form, and \(M(J)\) the full
multiplicity space. Evaluation at the chosen unit lowest vectors
identifies \(M(J)\) with the conjugate measured row, with one positive
squared-norm ratio \(c\). It respects the good projector. If \(j\) is the normalized composite on the arithmetic good row and \(V_{\rm full}\) the row before projection, then
\[
 \|jv\|^2=\|v\|^2\quad(v\text{ in the arithmetic good row}),\qquad \dim V_{\rm full}=D_{\mathrm{all}}.
\]
\end{proposition}

\begin{proof}
\proofstep{The two finite characters.}
Let \(J\) be the geometric level defining the refined row space and \(K_2\) the
corresponding subgroup in \(X\). Use the simultaneous
finite reduction of the type frames to first match
principal kernels. In the common finite group the two
marked sharp levels are conjugate. Lift the finite
conjugator inside the compact group and correct the
finite translate. If \(\iota\) is the group identification,
the corrected translate \(t\) satisfies
\[
 t^{-1}\iota^{-1}(K_2)t=J,\qquad
 \chi_J(t^{-1}\iota^{-1}(k)t)=\overline{\chi_X(k)}
 \quad(k\in K_2).
\]
\proofstep{Inverting lowest-vector evaluation.}
The central idempotents are unchanged by this conjugation
of frames. The real representation comparison uses the
contragredient discrete series at the same infinitesimal
character in the conjugate Hilbert space. Choose the
geometric unit lowest vectors corresponding to those
of \(X\). For an intertwiner in the full sharp multiplicity
space \(M(J)\), evaluation at these lowest vectors after
the translate \(t\) defines a map \(L\) onto the conjugate
row space of \(X\). Automorphy and the displayed level
and character conditions give the prescribed finite
equivariance. At infinity, for a unit lowest vector \(w\)
in a discrete series, let \(P_{v,w}\) be the operator
of its normalized matrix coefficient taking the column
\(w\) to \(v\). Schur orthogonality gives the rank-one matrix-unit identities
\[
 P_{w,w}^2=P_{w,w},\qquad P_{v,w}P_{w,w}=P_{v,w}.
\]
A row function \(f\) in the image of \(P_{w,w}\) gives the
intertwiner \(v\mapsto P_{v,w}f\) by covariance under
real translation of the matrix coefficients.

Evaluation
at \(w\) recovers \(f\). Apply the product of these
operators, using the contragredient identification and
the inverse finite translation. This gives surjectivity
and the inverse to evaluation. The construction goes
through the smooth discrete-series vectors, using the
matrix-coefficient projection with the actual formal
degrees. Rapid decay of cusp forms makes the evaluations
square integrable.
Petersson integration on the two identified quotients has
a single positive ratio \(c\). With the real vectors of
norm one this gives
\[
 \|LS\|^2=c\|S\|^2\qquad(S\in M(J)).
\]
The map is complex linear into the space of conjugated
functions with its complex structure there. The
finite action on both sides now intertwines with
conjugation of its coefficients. Thus \(L\) carries the
good part of \(M(J)\) onto the conjugate good subspace,
as seen by applying the central good projector inside
the full evaluation isomorphism.
\label{cor:real-row-realization-5}

\proofstep{Normalization on the good subspace.}
The geometric realization of the arithmetic row
\(E_{\rm row}\) is an isomorphism onto the linearized
good multiplicity space at \(J\). In the multiplicity evaluation, write \(d\mu\) for the measure before sharp-level normalization. The Petersson measure on the row is
\[
 d\mu_{\rm sharp}=n_P^{-1}d\mu .
\]
The good space at \(J\) is contained in the principal
model. Invert the faithful principal realization on
this subspace to obtain the geometric row embedding.
Its composition \(e_E\) with multiplicity realization
has image the whole good space of \(M(J)\), by surjectivity
of this geometric realization. Using the same unit tuple
in both evaluations gives \(\|e_E x\|^2=n_P\|x\|^2\).

It follows that
\[
             j=(c n_P)^{-1/2}L e_E
\]
is the isometry onto the conjugate good row space of the
measured model. The full evaluation bijection identifies
its full dimension with the recorded \(D_{\mathrm{all}}\).
This dimension is from the full space at the refined level before applying
the good projector. The two positive factors in \(j\) cancel
exactly in the pairing for every host and every allowed
word level.
\label{cor:real-row-realization-6}
\end{proof}

For the second replacement choose the spectral datum
using the row with its normalized Petersson measure and the comparison map \(j\). This choice fixes
\(t,c,j\) before any replacement of the division lattices
and heights on the word spaces. If a replacement retains
the degree-zero levels and readings, the term fibres are
carried by the two induced bijections of class labels.
The fibre of a term is the image of its original finite
fibre under those bijections, with equal function values.
Write \(b\) for the isometry induced by reindexing the good space. The transported row map is the composite, so for every row vector \(v\),
\[
 j_{\rm new}=b\circ j,\qquad\|j_{\rm new}v\|=\|jv\|=\|v\|.
\]

To apply
this comparison here take the original division lattices
and their height function. Transport first through the
symmetrized column form model, which preserves the
degree-zero values, then through the equivalence to the
actual rebuilt carriers of the same columns. It follows
that the new comparison of the row spaces at the refined level uses the original
column levels and actual column forms. The full
dictionary, the real tie, the discriminant and the
transported order and character data constitute the
level-four output for this pass.
\label{cor:real-row-realization-7}

\begin{lemma}[The ambient lattices in the used tuples]\label{lem:real-ambient-tuples}
At the sharp and bottom levels separately, inclusion of the used
tuples in the full tuple list identifies their coefficient spaces,
integral ambient lattices and positive Hermitian forms. Writing \(b\) for this identification at one of the two levels, \(\Lambda\) for the lattice, and \(H\) for the form, one has at each used tuple \(\mu\)
\[
 b\Lambda_{\mu,{\rm used}}=\Lambda_{\mu,{\rm full}},\qquad b^*H_{\mu,{\rm full}}=H_{\mu,{\rm used}}
\]
at every embedding where the form is positive. These identifications use the
same degree-zero columns as Proposition~\ref{prop:refined-real-row}.
\end{lemma}

\begin{proof}
The ambient lattices at the sharp and bottom levels use
the rational coordinates of this row specialization. Fix one of these levels and a tuple \(\mu\). Let
\(V_\mu\) be its coefficient-field coordinate space, with
lattice \(\Lambda_\mu\) and Hermitian matrix \(H_\mu\).
Scalar extension at the marked embedding sends its coordinate vectors to a
basis at the marked embedding, and carries \(\Lambda_\mu\)
onto the marked ambient lattice. At a positive embedding
\(\sigma\), choose an orthonormal frame for
\(\sigma(H_\mu)\).

Let \(F\) send the marked basis coordinates to coordinates in this orthonormal frame at \(\sigma\). Thus \(F^*F=\sigma(H_\mu)\), and
\[
 \|F\sigma(x)\|^2=\sigma(x)^*F^*F\sigma(x)=\sigma(x)^*\sigma(H_\mu)\sigma(x).
\] This is the positive form
in the comparison of the ambient norms. At the marked
embedding the original Hermitian and operator identities
give the designated complex form and adjoint under the
reading. Finite generation and spanning of the ambient
lattice follow from those of \(\Lambda_\mu\) and the two
identities for the basis and lattice.

The used tuple list maps into the full tuple list by the
same tuple of local character labels. The prescribed tuple
map agrees with this inclusion because the row indices are
defined by those labels. At either level the coordinate map on the selected tuples and the restriction of the full coordinate map agree on every basis vector. If these maps are \(f,g\),
\[
 f\!\left(\sum_i a_i e_i\right)=\sum_i a_i f(e_i)=\sum_i a_i g(e_i)=g\!\left(\sum_i a_i e_i\right).
\] The identity on coordinates, with the
dimension identified through that equality of indices,
therefore identifies the rational readings, integral
lattices and positive forms at every embedding where the
form is positive. This constructs the two ambient
identifications consumed by the division construction.
The sharp and bottom matrices and lattices are taken from
their respective levels.
\label{cor:real-row-realization-8}
\end{proof}

Both passes have this construction at their own
level-three states. The first level-four state is
completed to a built member before taking its signed
export. The rebuilt pass makes the corresponding choices
on the carried specialization. Their division and theta
constructions consume the corresponding real isometry and
ambient identifications on that pass.
\label{cor:real-row-realization-9}

\paragraph{The theta maps on this row.}
Fix one specialization and its levels at the primes assigned to the sequences. Denote the
unitary spaces in its chain by \(Y_0,\ldots,Y_m\), identifying
the definite top with \(Y_m\) by its given isometry. For a used tuple \(\mu\), take the multiplicity space \(M_j\) in the row at the refined level on \(Y_j\). The downward maps form the chain
\[
 M_m\xrightarrow{D_{m-1}}M_{m-1}\longrightarrow\cdots\longrightarrow M_1\xrightarrow{D_0}M_0 .
\] We need a simultaneous realization of the \(D_j\) by theta
integrals with the Petersson pairings on these very spaces.
This is stronger than existence of each local intertwiner.

\begin{proposition}[Isometric realization of the theta chain]\label{prop:isometric-refined-theta-chain}
At the fixed specialization and assigned levels, choose
isometric readings \(R_j\) of every multiplicity space on
a used tuple. With the stated Petersson measures and chosen
Weil and level characters, write \(\Theta_j^{\downarrow}\) for the downward theta integral transported back to the multiplicity rows by these isometries. On every edge of the chain, including the last,
\[
 \Theta_j^{\downarrow}=D_j.
\] On unused tuples the zero row realizes the
prescribed zero maps.
\end{proposition}

\begin{proof}
\proofstep{Fix the inner product multiplier at each stage.}
Use the identifications of the row spaces at the refined level for both the multiplicity
pairing and the downward integral. At zero these include \(M_0\)
in the base row, at intermediate indices in the cuspidal
multiplicity spaces, and at the last index in the definite one.
For the sharp index \(n_P\), the stage-zero measure is
\(n_P^{-1}\nu_0\) with pairing multiplier \(c_0=1\).
At the intervening stages the original quotient measures
have \(c_j=n_P^{-1}\). At the top the local definite factors \(M_{X,\tau}\) from the row enter with the common level divisor:
\[
 c_m=\frac{\prod_\tau \dim M_{X,\tau}}{n_P}.
\] The latter
identity is Schur orthogonality at the definite factors.
The same lower multiplier \(c_j\) occurs when the downward
model integral is paired against a vector of \(M_j\).
\label{cor:level-four-core-1}

\proofstep{Choose a common evaluation tensor serving both edges at each interior stage.}
For each edge its nonzero real intertwiner has a nonzero
matrix coefficient on a polynomial test and a pair of pure
smooth tensors, since such tensors span the algebraic tensor
product and separate its elements by the inner product. Fix
one such polynomial and a pair of tensors giving a nonzero coefficient
on each edge, including the last. Vary each tensor factor in the
finite real span of the corresponding chosen factors.

Let \(P_j(\mathbf t)\) be the resulting coefficient on edge \(j\) in the
common real coordinate space. At the pair chosen for edge \(j\) its coefficient is nonzero, so
\(P_j\not\equiv0\), even with the unused coordinates adjoined.
The polynomial ring has no zero divisors, so
\[
 \prod_j P_j(\mathbf t)\overline{P_j(\mathbf t)}\not\equiv0 .
\] At a point where the product is nonzero,
normalize each factor to unit norm and scale the test on each
edge to make its coefficient one. Thus both occurrences of a
stage use one and the same evaluation tensor. The resulting
polynomial test has nonzero Schwartz function by the model
intertwining identity.

\proofstep{Scale each row isometrically and then normalize its edge integral.}
Write \(\kappa_j>0\) for the required quotient multiplier on
\(Y_j\), including its finite and real Petersson factors.
Let \(r_j>0\) be the quotient-measure ratio. In the convention of this comparison the two pairings are related by
\[
 \langle\ ,\ \rangle_{Y_j,\kappa_j}=(\kappa_j/r_j)\,\langle\ ,\ \rangle_{\rm transported\ model}.
\] The Haar ratios include transport from the definite
top. Evaluate the multiplicity map on the chosen tensor and
multiply the transported cusp function by
\(s_j=(c_jr_j/\kappa_j)^{1/2}\). Schur orthogonality and rapid
decay give the original \(M_j\) inner product, because the multiplier for two readings is
\[
 \frac{\kappa_j}{r_j}s_j^2
 =\frac{\kappa_j}{r_j}\frac{c_jr_j}{\kappa_j}=c_j .
\] Denote this isometric reading by \(R_j\).

On the edge with upper space \(W=Y_{j+1}\), let \(r_W^I\)
be the transported quotient ratio for its integration variable.
Use the oriented Weil kernel from \(W\) to \(Y_j\), with its
recorded splitting characters. Multiply its test by
\[
 a_j=\frac{c_jr_j r_W^I}{\kappa_j\kappa_{j+1}s_js_{j+1}} .
\]
The integral against \(R_{j+1}y\), followed by the lower
Petersson pairing against \(R_jx\), gives the model pairing
with scalar
\[
 (\kappa_{j+1}/r_W^I)\, a_js_{j+1}s_j\,(\kappa_j/r_j)=c_j .
\]
The real intertwiner coefficient in this pairing is one.
The model chain identity now gives \(\langle D_jy,x\rangle\)
for all \(x,y\). This argument uses the convergent upper
theta integral of the marked cuspidal chain. It includes the
last edge by transporting its Weil action as well as its Haar.

Use the unique marked infinity data, the standard character of
sign \(-1\), and a unitary Weil realization for that character.
For a finite tuple character \(\chi\) there is an integer \(n\ge1\) such that, writing \(F\) for the totally real subfield of the auxiliary CM field,
\[
 \chi^n=1,\qquad \chi|_{\mathbb A_F^\times}=1 .
\] Their splitting characters
agree with the marked ones by the character identities on
ideles. On unused tuples the specified \(D_j\) are zero, so zero
row functions and the original nonzero tests give both required
edge integrals. The row isometry is required on used tuples.
In the theta realization the quotients and their Petersson
multipliers are exactly the ones just used. Equality of these
integrals therefore gives the required downward maps under
pairing, and the construction has the prescribed Weil
equivariance and finite characters on the same realization.
\label{cor:level-four-core-2}
\end{proof}

\begin{proposition}[Divisor with the marked row and finite support]\label{prop:core-divisor-support}
Use the refined theta chain of
Proposition~\ref{prop:isometric-refined-theta-chain} with
its sharp and bottom row comparisons at the same specialization.
The finite prime support of the divisor, denoted \(S_D\), is
\[
 S_D=\{p:\text{some component of }\xi_\mu\text{ above }p\text{ is ramified}\},
\]
where the character is the actual tuple character of the divisor. It has the same marked
coefficient-place valuations and row norm at the
initial flag datum.
\end{proposition}

\begin{proof}
Use the simultaneous theta kernels of
Proposition~\ref{prop:isometric-refined-theta-chain}.
Let \(E/F\) be the marked CM extension, \(P(v)\) the chosen prime
of \(E\) over each \(v\) of \(F\), and \(\xi_\mu\) the
original character of the column tuple. Use the rational lift
spaces, bases, integral modules and Gram matrices of the fitted
row. The sharp scalar is fixed before the finite type is chosen. For each type \(\sigma\) used here it is
\[
 \kappa(\sigma)=\kappa_{\rm sharp}.
\]

For each spectral label choose a reference finite type
above it if one occurs. A fallback type exists by positivity
of the reference primitive block. At the distinguished embedding
the logarithmic value is \(\log\lambda_{\mu_0,j,\sigma}\) from
the row. At another embedding the value at a tested lift is zero
on a zero isotypic component and is the logarithm of the chosen
type's reference reading otherwise. One choice of type for each
label works at every lift.

\proofstep{Identify finite support by contracting the same ramification places.}
Put \(S_\mu=\{v:\xi_\mu\text{ ramifies at some }P\mid v\}\).
The character pins give ramification at the chosen \(P(v)\) if
\(v\in S_\mu\), so the rational set is
\[
 Q_\mu=\{v\cap\mathbb Z:v\in S_\mu\}
      =\{P(v)\cap\mathbb Z:v\in S_\mu\}.
\]
For the reference \(\mu\) it is empty. For every tuple it
excludes the two type primes, and for \(\mu(w)\) it excludes
the residue characteristic at \(w\), by the local twist choice
on the support and the trivial tuple elsewhere.

If \(p=P(v)\cap\mathbb Z\) is a ramified character place,
the pin there gives \({\rm N}P(v)=p\) and
\(p\in P(v)\setminus P(v)^2\). The norm determines the residue degree, while \(p\notin P(v)^2\) determines the ramification index:
\[
 f(P(v)/p)=1,\qquad e(P(v)/p)=v_{P(v)}(p)=1 .
\]
Both indices are multiplicative in the tower, so the norm
downstairs is the same rational prime. This gives the
degree-one local field used for the tame determinant character
\(\delta\) and the unramified unitary quotient characters \(\nu_i\).

\label{cor:level-four-core-3}

Here is the local normalization. Write \(q\) for the residue
cardinality, \(A\) for the local ring of integers, and
\(\varpi\) for its uniformizer. On the edge out of stage \(j\)
use the Levi with blocks \(r,jt\). Write \(P\) for the parabolic in the second block with \(j\) blocks of size \(t\). Its spherical representation is obtained from normalized induction of the ordered determinant characters, whose inducing representation is
\[
 (g_1,\ldots,g_j)\longmapsto\prod_{i=1}^j\nu_i(\det g_i).
\] Take the algebraic cyclic span of its smooth spherical unit vector
\(v_0\). The maximal compact \(K\) has a one-dimensional fixed space
\(\mathbb C v_0\) by Iwasawa decomposition.

If a nonzero invariant
subspace contains \(u\ne0\), there is a translate \(gu\) with
\(\langle gu,v_0\rangle\ne0\), since the translates of \(v_0\)
span. With probability measure on \(K\),
\[
  \int_K \pi(k)\pi(g)u\,dk=\langle \pi(g)u,v_0\rangle v_0 ,
\]
using the linear-first Hermitian convention and \(\|v_0\|=1\).
Smoothness makes the average a finite linear combination, so
the invariant subspace contains \(v_0\) and equals its cyclic span.
For a diagonal entry \(z\) in block \(a=0,\ldots,j-1\) the
Borel law has factor
\(\nu_{a+1}(z)|z|^{(jt-t)/2-at}\) by the half modulus of
this parabolic. The first Levi block acts trivially in this
spherical model. For \(j=0\) the second block is trivial.
\label{cor:level-four-core-4}

Let \(J\) be the parahoric reducing to the stated upper block
parabolic \(P\) modulo \(\varpi\). For \(g\in J\) let \(d(g)\) be its lower diagonal block. Multiplication respects this block modulo \(\varpi\), so for the character of \(J\) one obtains
\[
 \chi_J(gg')=\delta(\det d(g))\delta(\det d(g'))=\chi_J(g)\chi_J(g').
\]
On \(P\cap J\) it agrees on the fixed unit vector with the
inducing law, since the diagonal blocks have unit determinant
norm. Extend this vector to a section on \(PJ\) by the two
transformation laws and set it to zero elsewhere. On the
maximal compact \(K\) its squared norm is the indicator of
\(K\cap PJ\). Give \(J\) volume one and scale the section by
\(\operatorname{vol}(K\cap PJ)^{-1/2}\). The matrix coefficient
at the identity is then one.

Use the dimension of this type
space and the doubling integral of the section in the divisor
data. The first-edge value can be read directly. The two-block test in these coordinates is supported on integral invertible \(X\) and \(\varpi\)-divisible \(Y\). There it has the value
\[
 (X,Y)\longmapsto\delta(\det X)^{-1},\qquad X\in\mathrm{GL}_r(A),\quad Y\in\varpi M_{t,r}(A).
\] Testing also \((Xg,Yg)\) forces
\(g\in{\rm GL}_r(A)\). On this compact group the ratio of the
tests cancels \(\delta(\det g)\) in the Weil action, and the
norm and unramified factors equal one. With the two full
integral matrix boxes of additive volume one, the first
value is therefore
\[
                q^{-tr}\prod_{i=1}^{r}(1-q^{-i}).
\]
If the residue field has cardinality \(q\) and the block has size \(d\), this product is the invertible-matrix density. Counting linearly independent columns gives
\[
 \frac{|\mathrm{GL}_d(\mathbb F_q)|}{q^{d^2}}=\prod_{a=0}^{d-1}\frac{q^d-q^a}{q^d}=\prod_{j=1}^d(1-q^{-j}).
\] All of these values refer to the primes of
\(Q_\mu\) just selected. \label{cor:level-four-core-5}

\proofstep{Realize the original coefficient valuations at the constant-field flag.}
For the initial flag datum use the constant function field equal
to the row field, with full section space there, one auxiliary
level supported on the two type primes, zero curve degrees and
zero maps from the row to sections. At each marked place it is the actual coefficient valuation that has to be recovered. For an embedding \(e\) into a valued field with valuation ring \(O\), the test is
\[
 a\in\mathcal O_{\rm coeff,\,w}\quad\Longleftrightarrow\quad e(a)\in O
\]
on the coefficient field. Embed the Witt fraction
field for each place in \(\C\) using a transcendence basis and
extend its valuation ring. The Witt field has at most continuum
cardinality. Conjugacy of the extensions of the \(p\)-adic place
to \(\overline{\Q}\) permits an embedding of the coefficient
number field inducing the prescribed place on restriction.
A complex field automorphism taking that embedding to the
distinguished one transports the Witt embedding and its
extended valuation ring together.

If \(e\) denotes one of the transported embeddings and \(O\) its chosen extended valuation ring, the local reading on the coefficient field is recovered by inverse image:
\[
 e^{-1}(O)=\{a:e(a)\in O\}.
\] Adjoin
this flag datum to the divisor, retaining its original row
basis, lattice, norm, logarithmic values, sharp factor and
quotient-prime readings. This is the base case of the geometric
condition closed under changing division lattices and heights.
The later geometric degree estimates use the parameter models
constructed on a completed point.
\label{cor:level-four-core-6}
\end{proof}

\begin{proposition}[The integral kernel on the selected divisor]\label{prop:level-four-integral-kernel}
On the selected divisor of Proposition~\ref{prop:core-divisor-support}
and the prescribed stage lattices,
let \(C_j\) be evaluation on the class charts, \(\Lambda_j\) the integral row, \(O\) the upper valuation ring and \(M_j^{\rm int}\) the module of equivariant functions with integral tensor Weyl values. At each valid positive stage, \eqref{eq:core-upper-image} reads
\[
 O\,C_j(\Lambda_j)=M_j^{\rm int}.
\] At each selected pair of
class representatives the transported Taylor kernel is
the prescribed entry series in the adapted bases.
The first maps and all these entries use the same divisor
and the ramified word lattice.
\end{proposition}

\begin{proof}
Fix a word \(u\) and a coefficient place \(w\) satisfying the
local orbit condition of Section~\ref{subsec:column-orbit-length}.
Use its marked column and the tuple \(\mu(w)\). Write \(k\) for
the starting column rank, \(t\) for the block size added on each
edge, and \(\Delta=\lceil x^{c_D}\rceil\). The bound
\(c=16k\Delta\) works for the column entries and the first
filling weight simultaneously on all sequences. The interior Harish--Chandra entries \(u\) and the two extreme entries off the top have the respective bounds and values
\[
 3\Delta-\tfrac12\le u\le4k\Delta,\qquad u_{\max}=4(k-2)\Delta+4k\Delta+\tfrac12,\qquad u_{\min}=-\tfrac12 .
\]
The rank shift in passing to column entries has absolute value
at most \(k\).

At the first edge the gap above the lower
endpoint \(b_1\) is at least \((t+3)/2\). The prescribed
top floor, lower bound for the larger weights \(b_1\ge\lceil b\rceil\), and thin
choice \(b_1=\lceil w''\rceil\) each give \(b_1+1\ge w''\),
where \(w''=b'/(xn)\). Uniformly in \(k\le x^5\), \(n\le x^3\),
we have eventually
\[
 c+k+2\le35x^{c_D+5}\le x^{c_D+6}\le w''\le b.
\]
Thus even after subtracting the shift \((k+t+1)/2\) the first
edge weight is at least \(c\). The size conditions give
\(c\le c_G(L)\) on every compatible size record \(L\).
This is a single bound chosen before those records.

In each orbit of the local step \(h\), take the residue of
least numerical label. If \(b\) is the initial label in an orbit and \(d\) the order of the step, the arithmetic and geometric rules both take the minimum of the same finite set:
\[
 \min_{0\le s<d}\,\operatorname{label}(b+sh),\qquad d=\operatorname{ord}(h).
\] Send its translate by \(sh\) to the
corresponding marked embedding \(\tau\) at slot \(s\).

At top slots the endpoints are the first and the negative of
the last column entry. Put zero in these two positions off
the top. After subtracting the rank shift \(a_j\), the added weight at lift \(j\) and embedding \(\tau\) is
\[
 w_j(\tau)-a_j .
\] These are the
negatives of the first entries of the successive added blocks
of a dominant column. They increase weakly, and the first
dominates the negative last entry of the original block.
For every compatible canonical completion the orbit, the
positions, the top endpoints and the shifted weights just
computed are therefore exactly its prescribed filling. This
comparison concerns the columns and weights alone, so it
persists for each choice of the subsequent level and maps.
\label{cor:level-four-core-7}

Take the standard primitive lift chain at the \emph{actual}
\(Q_\mu\). Use its nonzero unfolding factors, integral jet
identities and adapted bases. In the basis of an entry place one
rescales the stage \(j\) frame at \(\tau\) by the prescribed unit
\(a_\tau^j\). A product of minors of shape \(\lambda\) has total degree \(|\lambda|\). Thus stage-\(j\) scaling by the unit \(a_\tau^j\) gives the multiplier
\[
 (a_\tau^j)^{|\lambda|}=a_\tau^{j|\lambda|}.
\] The scalar substitution in the kernel polynomial contributes a further
\(a_\tau^{|\lambda|}\). The homogeneous factors on the edge are therefore
\[
 a_\tau^{j|\lambda|}a_\tau^{|\lambda|}
        =a_\tau^{(j+1)|\lambda|},
\]
the factor for the target frame. The source and target determinant
normalizations occur in the unfolding scalar with the
corresponding powers. Dividing by this nonzero scalar cancels
them. The unit preserves both integral tests and the local norms.
Thus the entry coefficient in the adapted bases is the Taylor
coefficient of the kernel for these lift objects.
\label{cor:level-four-core-8}

For each marked local component \(v\) let \(A_v\) be its ring
of integers and \(L_v\) the prescribed full lattice at one
stage. The first split embedding \(\iota_v:E\to F_v\) on \(F\)
is its completion at \(v\), and these places of \(F\) are
distinct. Put
\[
 \Lambda=\{z\in E^d:\ \iota_v(z)\in L_v\text{ for all marked }v\}.
\]
It is an \(\mathcal O_E\)-module. Its image spans \(L_v\) over
\(A_v\): approximate a basis matrix for \(L_v\) within its
open integral basis coset at \(v\), and zero within the other
lattices at all the other marked places, by a matrix over \(F\).
Its columns belong to \(\Lambda\) and are a local basis at \(v\).
This weak approximation argument asserts a span at the specified
places. The global finite lattice with all its other localizations
prescribed was constructed separately above.

For the local class map first fix an away-place representative
\(a\), then a local matrix carrying the reference lattice to
a tuple \(L'\). Every admissible adelic class is obtained in
this manner. For a fixed away class \(a\), let \(\Gamma_a\) denote its rational stabilizer. Two presentations with local tuples \(L'_1,L'_2\) give the same class exactly when
\[
 L'_2\in\Gamma_a\cdot L'_1 .
\] This is the double
coset equivalence after fixing \(a\). Stage zero uses the given
word-column identification. At a positive stage the reference
tuple uses its member charts, with just the end chart at stage
\(m\).

For a nonreference tuple use the symmetric chain of its
own character. For the frame matrix \(J=C^*C\), the reference basis \(B=C^{-1}\) therefore satisfies
\[
 B^*JB=(CB)^*(CB)=1.
\] Its lattice and the
level lattice agree upon completion at \(v\), since the latter
is between the lattice and its multiple by the word modulus.
In the nonreference frame \(J=B=1\) the last-coordinate
condition disappears there since \(\pi_T\bar\pi_T=T\) is a
unit. Moreover the quotient-prime set for \(\mu(w)\)
excludes its own characteristic. Thus on both kinds of charts
the local level is the stabilizer of the indicated lattice.
\label{cor:level-four-core-9}

Here is the integral image assertion needed for the upper maps.
Let \(e:H\to K_w\) be the coefficient embedding, \(O\) the upper valuation ring and \(A\) the valuation ring of \(H\) at \(w\). The compatibility of the two integrality tests is
\[
 e^{-1}(O)=A,\qquad h\in A\ \Longleftrightarrow\ e(h)\in O .
\]
Let \(V_j,\Lambda_j\) be the rational row and its integral
lattice at a valid positive stage. Let \(M_j^{\rm int}\) be
the local module of equivariant functions on the class
presentations, taking values in the tensor Weyl lattices of
the filled weights. Then the chart reading \(C_j\) satisfies
\[
                O\,C_j(\Lambda_j)=M_j^{\rm int}.
                \tag{*}\label{eq:core-upper-image}
\]

To see the two integral tests involved, write
\(d=\lfloor j(t+1)/2\rfloor\) for the determinant shift. At
a class \(y\), the normalized evaluation multiplies the
original frame evaluation by
\[
             \chi_D(\det y)^d\prod_\tau
                   \det\bigl(\tau(B)^{-1}z_\tau\bigr)^d .
\]
In a rational basis \(B'\) of the local dual lattice the change
of determinant weight multiplies the coefficient polynomial by
\[
             s=\chi_D(\det y)^d
                \prod_\tau\tau(\det B'/\det B)^d.
\]
The finite determinant-character values are rational over the
coefficient field by the marked character reading and its
values on uniformizers. The projections of the bases are
the marked embeddings of the host. To check the valuation of
\(s\), use finiteness of the class set to write
\(y^{b+l}=\gamma y^b k'\) with \(l>0\), \(\gamma\) rational
unitary, and \(k'\) in the level.

On each first split component,
both \(B'^{-1}yB\) and \(B'^{-1}yk'B\) have unit determinant
by the local basis equality and self-duality. Taking determinants
in the power relation shows that
\(z=(\det B'/\det B)^l/\det\gamma\) is a unit there.
Also \(z\bar z=\det(B'^*JB')^l\) is a unit by self-duality,
so the second split component gives a unit. Applying the character law at \(\gamma\) and \(k'\) now yields
\[
 s^l=\prod_\tau\tau(z)^d .
\] Hence \(s\) and its inverse are
integral. This calculation uses both projections even when the
slot reads the conjugate embedding. \label{cor:level-four-core-10}

The coefficient change at a translated lattice uses
\((B^{-1}\gamma B')^{-1}B^{-1}=B'^{-1}\gamma^{-1}\).
On an ordinary block this substitutes the translated bases
in products of minors. On a dual block the complementary-minor
identity for the inverse transpose, with the column order
reversed, gives integral signs and the determinant correction.
The inverse substitutions give both directions. Tensor over
the slots. The level character multiplies by a root of unity,
and descent along \(e^{-1}O=A\) tests coefficients on a basis
of the rational Weyl lattice. Thus for \(x\in V_j\)
\[
             C_jx\in M_j^{\rm int}\quad\Longleftrightarrow\quad
             x\in A\Lambda_j .
\]
For surjectivity after taking the span, read a local function
at representatives of the finitely many class orbits.

At each
representative its values obey the stabilizer equations with
the inverse character of the glued away-level element. These
are linear equations on a finite-dimensional rational Weyl
space, hence finitely many equations suffice over every scalar
extension. Conversely each rational solution gives a complex
form on that adelic class. Indeed if \(y=\gamma yk'\) the
determinant shift contributes one to the stabilizer equation
by the character law. At the chosen orbit representative the stabilizer is conjugate at level to \(k'^{-1}\). Since its character is constant on conjugacy classes, inverting the glued character recovers
\[
 \chi(k'^{-1})^{-1}=\chi(k').
\] Ordinary positions use the
marked entry of \(\gamma\), dual positions its conjugate in
the reversed dual frame. The resulting frame polynomials are
independent whenever their restrictions to invertible blocks
are independent, since these blocks are Zariski dense and the
frame changes are invertible there. Extend by equivariance and
by zero on other classes.

Independence after restriction to the open set of invertible blocks therefore bounds the integral rank by
\[
 \operatorname{rank} M_j^{\rm int}\le \dim_{\mathbb C}(\text{complex row at }j).
\] In the other direction, the original
frame evaluations separate \(V_j\). A nonzero minor of finitely
many rational coefficient readings remains nonzero over \(K_w\),
giving the opposite bound. Equal rank and that same minor give
a rational preimage for every rational target in the span.
\proofstep{Obtain the integral generators after imposing the stabilizer equations.}
Finally choose bases of all marked local lattices by simultaneous
weak approximation at the distinct places of \(F\). Their local
Weyl lattices extend rational \(A\)-lattices. Intersecting with
the stabilizer equations commutes with extension to \(O\) by
flatness over the DVR \(A\). The resulting rational integral
generators extend by equivariance to the finitely many orbits.
Their rational preimages pass the two-way integral test. This
proves \eqref{eq:core-upper-image} with the original lattice at
each stage. \cite{companion-selected-core}\label{cor:level-four-core-11}

\proofstep{Pull the entry back by the inverse of the exact coordinate substitution.}
For a pair of adjacent chosen bases with image matrices \(A_j\)
and \(A_{j+1}\), the forward substitution of a kernel block is
\[
 X\longmapsto A_{j+1}XA_j^{-1}
 \quad\hbox{or}\quad
 X\longmapsto A_{j+1}^{-\mathsf T}X A_j^{\mathsf T}
\]
according to the marked ordinary or dual position. Pull the
fixed entry series for the pair of global classes back by the
inverse linear substitution. Composing with the forward map
returns it coefficient by coefficient, since the matrices
cancel and all substitutions preserve the degree filtration.
The split idempotents identify the embeddings with those of
the chosen column, and the orbit-position equality identifies
the variables. Attach these series, the local class maps and
the image lattices \eqref{eq:core-upper-image} to the fixed
filling. Use their first comparison maps and the ramified
submodule specified by the word column. This gives the kernel
at every selected pair with the adapted entry identity and the
same divisor throughout. \label{cor:level-four-core-12}
\end{proof}

\paragraph{Assembling the theta maps at the specified level.}
Use the section and kernel of Proposition~\ref{prop:level-four-integral-kernel}
with the theta maps of Proposition~\ref{prop:isometric-refined-theta-chain}.
The two ambient identifications, the quotient primes and the
theta realization were obtained for the same level-four data.
Form the object whose sharp and bottom ambient modules are those
just identified, whose divisor has those primes, and whose
kernel and lattice comparisons are those just constructed for
that divisor. Include its original degree-zero column and
level, normalized character, signed row, and a CM gauge for the
host. The divisor comparison applies to every object with
this divisor, base row and row reading. All three equalities
hold here by construction. Thus the logarithmic type choice,
the rational roots on the base row, the stable multiplicity
reading and the sharp normalization belong to the assembled
object.

In particular, writing \(S\) for the selected twist support, nontriviality on \(S\) gives
\[
 \mu(w)=\mu_0\quad\Longrightarrow\quad w\notin S .
\] At its conjugate the reference tuple is again
assigned. The untwisted identities therefore give zero for the
signed twisting order on this tuple. For any admissible size record of the same scale, a nonzero forced content would itself put the place in the support. Thus for each such record
\[
 v\notin\operatorname{supp}\quad\Longrightarrow\quad\text{forced content at }v=0 .
\] The comparison of
the local first maps and the orbit weights above holds for
each such record with its canonical coordinates. This gives
the required row and column data at this level.
\label{cor:level-four-core-13}

Perform this construction first on the initial member. Its
signed export below uses the rational and integral theta
readings of this member. After transporting its auxiliary
character, index and series to the new field, construct the
second member using its own measured row and columns by this
same argument. On that member the transported norm choice
gives both the original norm condition and the condition on
all twist places. The signed tests descend along the kept
identifications. The second member consequently supplies a
single row, kernel and divisor over the chosen levels.
\label{cor:level-four-core-14}


\begin{auxlemma}[The theta generator on a semistable split Hecke branch]\label{lem:marked-hecke-semistable-unit}
Let \(R\) be a complete discrete valuation ring of mixed
characteristic \((0,\ell)\) with finite residue field. Use a
principally polarized semistable PEL fibre \((A,\lambda)\),
with \(\mathcal O_E\)-action of Rosati involution \(a\mapsto\bar a\).
Write \(\mathcal P_A\) for the rigidified Poincar\'e bundle on
\(A\times A^\vee\). The marked theta line on this fibre is
\[
 L_A=(1,\lambda w_0)^*\mathcal P_A,\qquad
            w_0\in\mathcal O_E,\quad w_0+\bar w_0=1 .
\]

Extend \(L_A\) cubically to the identity component of the
semistable model. For a split Hecke prime \(p\ne\ell\) away
from the marking, the \(\mathcal O_E\)-linear branch and its
target line satisfy
\[
       f:A\longrightarrow B,\qquad f^\vee\lambda_B f=p^2\lambda,
              \qquad L_B=(1,\lambda_Bw_0)^*\mathcal P_B .
\] The theta trace on this branch carries an integral
generator to an integral generator after the finite extension
defining the branch. This holds also for \(p=2\) when this split
prime is admitted.
\end{auxlemma}
\begin{proof}
\proofstep{Identify the model norm for this theta line.}
After finite extension take split Raynaud uniformization with period
lattice \(Y\), character lattice \(X\) and principal polarization
\(\Phi:Y\overset{\sim}{\to}X\). Write \(t(y,x)\) for the
biextension trivialization on \(Y\times X\) and set
\[
     b(y,z)=-\log\|t(y,\Phi z)\| .
\]
This is the positive definite monodromy pairing.

The Poincar\'e trivializations on the two period axes are rigidified.
Translating both arguments multiplies the axis factors by the
biextension value at the pair of periods. At the rigidified origin
the axis factors are one, so the graph trivialization is
\[
                     c(y)=t(y,\Phi w_0y)\qquad(y\in Y).
\] The resulting line on the abelian part is the same graph
pullback of its Poincar\'e bundle. Bilinearity identifies the
rigidified biextension with its pullback by \((-1,-1)\), so both
graph lines are symmetric. Functoriality with respect to
\(w_0\) and the adjoint identity give
\[
          -\log\|c(y)\|=b(y,w_0y)=b(y,\bar w_0 y)
                      =\tfrac12 b(y,y).
\]
The graph description thus specifies the characteristic of the
actual line. \cite[preprint, Section~3.4 and Theorems~3.10, 4.10]{frss-nonarch-theta}

For a nonzero principal theta section \(s\), the coefficient at
\(\Phi y\) in the Raynaud expansion is its coefficient at zero
translated on the good abelian part and multiplied by \(c(y)\)
in the corresponding Poincar\'e fibres. Translation preserves the
model norm there. Write \(\gamma\) for the negative logarithm of
the norm of the zero coefficient at the generic point of the good
special fibre. In skeleton coordinates \(z\in Y_{\mathbb R}\),
the Gauss formula and the quadratic correction for the cubical
metric give
\[
 \begin{aligned}
   -\log\|s(\xi_z)\|_{\rm can}
    &=\gamma+\tfrac12 b(z,z)
       +\min_{y\in Y}\{\tfrac12 b(y,y)+b(y,z)\}\\
    &=\gamma+\tfrac12\min_{y\in Y}b(y+z,y+z).
 \end{aligned}
\]
Indeed the correction \(\tfrac12 b(z,z)\) cancels the change
\(b(y,z)+b(y,y)/2\) in the period functional equation. It gives the
rigidified canonical metric by compatibility with pullback by
\([2]\), which multiplies the quadratic form by four.
\cite[Proposition~9.1 and its proof]{dejong-shokrieh-theta}

The displayed minimum is zero at the origin of the skeleton.
The supremum on each retraction fibre is attained at its Gauss point.
At the identity component the negative logarithm there is the
vertical order of the section in the cubical extension, using the
valuation with \(-\log|\pi|=1\). Hence
\[
       -\log\sup_{x\in A^{\rm an}}\|s(x)\|_{\rm can}
                    =\operatorname{ord}_{0}(s).
\]
A section on the generic fibre extends to the identity model
precisely when this order is nonnegative: the model is smooth,
the special fibre is integral, and all its other divisorial points
lie on the generic fibre. Thus a generator has supremum norm one.
The same calculation applies to \(B\). \cite[Proposition~10.3 and
Corollary~10.4]{dejong-shokrieh-theta}

\proofstep{Compare the two kernel actions at the actual level.}
Let \(C=\ker f\), \(C_0=C\cap A[p]\) and \(\mathcal L=f^*L_B\).
Write \(\mathcal B_A=(1\times\lambda)^*\mathcal P_A\) on
\(A\times A\), with \(\mathcal B_B\) defined similarly.
Functoriality and the polarization relation give rigidified
isomorphisms of biextensions and of their restrictions to the graph:
\[
 (f,f)^*\mathcal B_B\simeq([p],[p])^*\mathcal B_A,\qquad
             \mathcal L\simeq[p]^*L_A.
\]
For \(d\in C_0\), the two lifts of translation by \((d,0)\)
differ by an invertible function on the proper connected
\(A\times A\), hence by a scalar. Its restriction to
\(A\times\{0\}\) is one by rigidification. Apply the same
argument on \(\{0\}\times A\) to \((0,w_0d)\), using
\(w_0d\in C_0\) by \(\mathcal O_E\)-linearity. Composing along
the graph identifies the splittings \(U\) from \(f\) and \(V\)
from \([p]\):
\[
                    U_d=V_d\qquad(d\in C\cap A[p]).
\]

Put \(S=[p]^*s\). The trace is \(C\)-invariant and therefore
defines a section \(s'\) on \(B\) by
\[
                       f^*s'=\sum_{c\in C/C_0}U_cS .
\]
The rigidified pullback isomorphism preserves the canonical metrics
by their multiplication law. Thus \(U_c,V_a\) are isometries for
the supremum norm: each is a deck action on its own pullback line.
The summands have distinct \(A[p]\)-characters by the perfect
commutator pairing on \(A[p^2]\), since \(A[p]^\perp=A[p]\) there.
Over a finite extension splitting the torsion, the zero-character
projector therefore gives
\[
        |A[p]|^{-1}\sum_{a\in A[p]}V_a(f^*s')=S .
\]
Its denominator is an \(\ell\)-adic unit. Apply the ultrametric
inequality first to the trace and then to the projector:
\[
 \|f^*s'\|_{\sup}\le\|S\|_{\sup}\le\|f^*s'\|_{\sup}.
\]
Surjectivity of \(f\) and \([p]\) now gives
\(\|s'\|_{\sup}=\|s\|_{\sup}\). The first part of the proof now identifies the
integral generator on the target.

The complementary isogeny
\(f'=\lambda^{-1}f^\vee\lambda_B\) satisfies \(f'f=ff'=[p^2]\).
After extension both maps act on the identity components of the
N\'eron models. Since \(p\) is a unit, \(p^{-1}df\) and
\(p^{-1}df'\) are mutually inverse integral maps on each saturated
eigensummand of the Lie lattice. Thus the tangent frame and theta
generator remain integral on every such branch, including traits
with toric reduction. Counting the branches gives an integral sum
on these frames. Identifying this operator with a spectral Hecke
operator on a specified rational subspace further requires the
equivariance of its realization in the marked sections.
\end{proof}

\begin{auxlemma}[The geometric local unit ball and its scalar extensions]
\label{lem:marked-semistable-local-ball}
Let \(v\mid\ell\) be a coefficient place of the embedded small field
\(F\). At a fixed marked level and weight, take \(V\) in the space of
\(F\)-rational cusp sections.
For every complete discrete valuation ring \(R\) with finite residue
field and fraction field \(K\), take all embeddings
\(F\to K\) over \(v\) and all framed semistable tests \(y\)
over \(R\). The tangent vectors are bases of the saturated
\(\mathcal O_E\)-eigenlattices and the theta section is a generator
on the identity component. Use the extended valuation with
\(v_R(\pi_v)=1\), where \(\pi_v\) is a uniformizer at \(v\). Set
\[
 M_v(V)=\{s\in V:\ v_R(s(y))\ge0\quad
                \text{for every }(R,F\to K,y)\}.
\]

Let \(P\) be an admitted split place away from the level and
splitting-character conductor, of rational characteristic \(p\ne\ell\).
Write \(q_P=|\mathcal O_E/P|=\mathrm NP\). Fix a Hecke degree
\(0\le a\le r\) and let \(n_V\) be the even theta weight.
The geometric branch sum with scalar
\[
 c=q_P^{-a n_V/2},\qquad u_{\rm geom}=c(\mathrm NP)^{-a n_V/2}
\]
preserves this ball if it preserves \(V\). For a field extension
\(F'/F\) and any valuation ring \(O\) of \(F'\) inducing \(v\), it
preserves
\[
  \langle1\otimes M_v(V)\rangle_O
       \ \cap\ (F'\otimes_F V_0)
\]
for every invariant subspace \(V_0\subseteq V\). The same conclusion
holds for an invariant subspace over a larger coefficient field in
the corresponding scalar-extended ambient space.
\end{auxlemma}
\begin{proof}
On a test \(y\) each geometric kernel and its quotient may be defined
over a finite extension. By Lemma~\ref{lem:marked-hecke-semistable-unit}
the transported frame is integral there. Replacing its eigenbasis by
an integral basis of the same flag and multiplying its theta generator
by a unit change the weight evaluation by a unit. With the valuation
still normalized on \(F\), each term in the branch sum satisfies
\[
    v_{R'}(s(y_C))\ge0,\qquad
    v_{R'}\!\left(u_{\rm geom}\sum_C s(y_C)\right)\ge0 .
\]
Both integers \(q_P\) and \(\mathrm NP\) are powers of \(p\).
The estimate descends to \(R\) on the original evaluation.
Preservation of the \(O\)-span follows by linearity. Intersecting
two invariant modules gives the stated restriction for valuation
rings of arbitrary rank.
\end{proof}

In geometric coordinates the normalized trace contributes
\((\mathrm NP)^{-a n_V/2}\). Writing its adelic evaluation as
a sum of right translations introduces the splitting character
of the reading. With \(c=q_P^{-a n_V/2}\) the respective
coefficients are
\[
 u_{\rm geom}=c(\mathrm NP)^{-a n_V/2},\qquad
 u_{\rm ad}=c\,\mu_0(\varpi_P)^{-a}.
\]
To compute the valuation of the character, choose \(h>0\) with
\(P^h=(\alpha)\) and \(\alpha\equiv1\) modulo its conductor.
The algebraic infinity type expresses a fixed positive power of
\(\mu_0(\varpi_P)^h\) as a root of unity times a monomial in the
conjugates of \(\alpha\). Each conjugate has all its finite prime
divisors over \(p\), so
\[
                    v(\mu_0(\varpi_P))=0\qquad(v\nmid p).
\]

A finite family of CM evaluations can be used to bound a prescribed
collection of vectors inside a full \(M_v(V)\). Its fixed inverse
matrix costs the power of \(\pi_v\) in
Lemma~\ref{lem:marked-bounded-denominator}. If the realized
spectral span is the quotient of the same class-input space by the
kernel of the geometric evaluation, this construction gives a
specific integral module in that quotient, by inverse image under
the injection into cusp sections. Using its ball for an independently
specified rational row requires the specified reading to identify
the two evaluations and their normalizations.

\begin{auxlemma}[Finite harmonic jets in a semistable theta frame]
\label{lem:marked-bounded-jet-frame}
Fix the embedded CM field \(E\), the dimensions, an order \(d\)
of differentiation and a coefficient place \(v\mid\ell\). Use the
principal graph line of Lemma~\ref{lem:marked-hecke-semistable-unit}.
On each of the semistable tests in
Lemma~\ref{lem:marked-semistable-local-ball}, evaluate the harmonic
jet of order \(d\) in any integral eigenframe. Start with sections
of canonical supremum norm at most one, and form a fixed finite
linear combination of their rigidified isogeny pullbacks and
finite-order theta translations, of the marked tensor weight.
Writing \(J_P\) for this contracted jet, there is a constant \(b\)
depending only on the order, the polynomial, the tensor weight and
the finite operations such that
\[
             v_R(J_P(y))\ge-b\quad\text{for every test }y,
                         \qquad v_R(\pi_v)=1 .
\]
\end{auxlemma}
\begin{proof}
\proofstep{Trivialize the biextension to finite order.}
Let \(G\) be the identity model. The biextension
\(\mathcal B=(1,\lambda)^*\mathcal P\) on the generic fibre extends
to \(G\times G\) as the polar biextension of its cubical line.
In particular the two rigidified restrictions have the generic identification
\[
            \mathcal B|_{(1,w_0)}\simeq L_A .
\]
It extends to the identity model: the
generic isomorphism can have vertical divisor only a multiple
of the integral special fibre, and comparison at the identity
makes its order zero. The same argument extends the rigidified
adjoint isomorphisms for the \(\mathcal O_E\)-action.

Use logarithmic coordinates \(z,w\) on the two formal groups over
the fraction field. In smooth integral parameters the logarithm
is obtained by integrating the invariant differentials with identity
linear term. Its coefficients through degree \(D=\max(2,d)\) and
those of its inverse have bounded denominators depending only on
\(D\): integration divides by integers at most \(D\), and inversion
with identity linear term uses finitely many additions and products.
Thus for the logarithm \(l\) and exponential \(e\) of the formal group there is an integer \(c_D>0\), independent of the test, with
\[
    c_D [z^\alpha]l,\quad c_D[z^\alpha]e\ \in R^{\dim G}
                   \qquad(1\le|\alpha|\le D).
\]
Substitution along integral tangent eigenvectors preserves this bound.

Pick an integral formal frame for \(\mathcal B\) trivialized on
both axes. Denote by \(h_1,h_2\) the logarithms of its discrepancies
under doubling the first or second variable, using the biextension
law to identify the resulting line with its tensor square.
The discrepancy of addition in the first variable is one if either
summand is zero. On putting the summands equal, its logarithm
therefore has first-variable degree at least two. Likewise \(h_2\)
has second-variable degree at least two. Commutativity of the two
doublings gives in bidegree \((a,b)\)
\[
 (2^b-2)(h_1)_{ab}=(2^a-2)(h_2)_{ab}\qquad(a,b\ge1).
\]

Multiplying the frame by \(\exp H\) changes the two logarithms by
\((2^a-2)H_{ab}\) and \((2^b-2)H_{ab}\), respectively. Thus one
can make both vanish by setting
\[
 H_{ab}=\begin{cases}
 -(h_1)_{ab}/(2^a-2),&a\ge2,\\
 -(h_2)_{ab}/(2^b-2),&a=1,\ b\ge2,\\
 0,&a=b=1.
 \end{cases}
\]
For the addition discrepancy of the new frame, simultaneous
doubling of the two summands in the first variable doubles its
logarithm. It has total degree at least two in those summands
and consequently vanishes. The second variable gives the same
conclusion. This constructs a frame compatible with both addition
laws; at order \(D\) only finitely many of the indicated integers
and exponential factorials occur in denominators.

\proofstep{Impose the adjoint identity.}
For \(t\in\mathcal O_E\), the discrepancy between the frame at
\((tz,w)\) and at \((z,\bar t w)\) is the exponential of a bilinear
form \(B_t(z,w)\), since both frames respect both additions.
For a pair of eigenlabels \(\tau,\sigma\), put
\(\sigma^*(t)=\sigma(\bar t)\). On this pair, addition and composition
of endomorphisms give
\[
 B_{t+s}=B_t+B_s,\qquad
 B_{ts}=\tau(s)B_t+\sigma^*(t)B_s,\qquad B_m=0\quad(m\in\mathbb Z).
\]

If \(\tau=\sigma^*\), these identities give a derivation of the
separable number field \(E/\mathbb Q\) with values in the coefficient
space. It vanishes, as follows by differentiating a separable
minimal polynomial. Otherwise choose one primitive integral element
\(t_0\) of \(E\). Replacing the frame by the exponential of a
bilinear form changes \(B_t\) by
\((\tau(t)-\sigma^*(t))\) times that form. Choose its coefficient
to kill \(B_{t_0}\). Commuting \(t\) and \(t_0\) in the adjusted frame gives
\[
       (\tau(t_0)-\sigma^*(t_0))B_t
           =(\tau(t)-\sigma^*(t))B_{t_0}=0 .
\]
The first factor is nonzero, so every \(B_t\) vanishes.

The extra denominator is one of the finitely many nonzero
differences of conjugates of \(t_0\). Each coefficient of
\(B_{t_0}\) is bounded in the actual eigenvectors, since it
uses only the integral adjoint isomorphism and the preceding
finite-order frame. The adjusted frame has the required uniform
denominator when evaluated in those vectors.

The graph of \(w_0\) now gives a formal frame of the theta line.
Two such adjoint-compatible frames differ on the graph by the
exponential of a quadratic form coupling only opposite eigenlabels.
Over all the columns its coefficients give the contractions
\[
       \sum_h\partial_{\tau,i,h}\partial_{\bar\tau,j,h} P=0
\]
of the specified harmonic symbol \(P\). Multiplication by this
exponential therefore leaves its contracted jet unchanged. This
identifies the computed harmonic jet with the one in the marked
formula, including the one read in complex coordinates.

\proofstep{Bound the translated germs.}
A rigidified pullback between the corresponding symmetric lines
is isometric for their canonical metrics by the multiplication
law. A finite-order theta lift above torsion of order \(h\) also
preserves the metric. Indeed take \(m>1\) with \(m\equiv1\bmod h\).
The discrepancy of the two metrics, as a real logarithm \(b\) on
the compact analytic abelian variety, obeys
\[
                  b([m]x)=m^2 b(x)+c
\]
for a constant \(c\), since the two rigid line isomorphisms differ
by a scalar. Surjectivity of \([m]\) forces the oscillation of \(b\)
to equal \(m^2\) times itself. Thus \(b\) is constant, and a power
of the lift is the identity, so \(b=0\).

By the identification at the identity Gauss point in
Lemma~\ref{lem:marked-hecke-semistable-unit}, each resulting
norm-bounded section has vertical order at least zero there.
For exterior products and symmetric powers use the tensor powers
of the same cubical metric at this point. The section consequently
has an integral germ in any integral local frame.
Converting its \(d\)-jet to the graph frame, applying the harmonic
differential and summing the fixed coefficients cost only the
finitely many denominators above. This proves the bound.
\end{proof}

\begin{auxlemma}[Fullness on the rational span of theta classes]
\label{cor:marked-full-class-ball}
Let \(V\) be the image of a finite-dimensional space of fixed
definite class inputs under the marked algebraic theta map. At
every coefficient place \(v\), its local unit ball \(M_v(V)\)
satisfies
\[
      M_v(V)\simeq\mathcal O_{F,v}^{\dim_F V},\qquad
                  F\otimes_{\mathcal O_{F,v}}M_v(V)=V .
\]
Its inverse
image on the class space modulo the kernel gives the same lattice
in that rational realization of the spectral span.
\end{auxlemma}
\begin{proof}
Choose a basis \(s_1,\ldots,s_m\) of \(V\) among the class images.
The principal theta generator has norm one by
Lemma~\ref{lem:marked-hecke-semistable-unit}. Its fixed class
formula consists of the finite operations of
Lemma~\ref{lem:marked-bounded-jet-frame}. Thus a single \(b\ge0\)
makes \(\pi_v^b s_i\) integral on every test. Use the separating
marked CM evaluations on \(V\), at integral frames after one finite
extension of the local field. Their fixed matrix is invertible on
these basis sections. Every vector in the unit ball has integral
readings there, so the inverse matrix bounds its coordinates:
\[
  \pi_v^b\sum_i\mathcal O_{F,v}s_i\ \subseteq\ M_v(V)
       \ \subseteq\ \pi_v^{-c}\sum_i\mathcal O_{F,v}s_i
                         \qquad(c\ge0).
\]
Over the discrete valuation ring this submodule is finite free of
rank \(m\). The assertion includes \(m=0\). Pulling back to the
quotient by the kernel preserves both inclusions. For a different
prescribed realization one needs its identification with this
particular quotient before transporting the lattice.
\end{proof}

\begin{proposition}[Signed theta coefficients on the original specialization]\label{prop:signed-theta-data}
On the initial member with the prescribed row, tuple, and
coefficient embeddings, the normalized signed theta coefficients
are rational in the original coefficient coordinates. Fix an admitted
index \(i\), lift, sign and theta datum \((s,h,P)\) consisting of support,
adelic argument and harmonic polynomial. Let \(L_i^{\rm sup}\) be the
supremum lattice over the coefficient DVR \(O_i\), and \(R\) the signed
series map on the rational row. Write \(x_{Q,M,\delta}\) for the
normalized row vector of the coefficient at \([\delta]\) of the universal
series with auxiliary level \(Q\) and depth \(M\). There is one integer
\(N\ge0\), depending on the fixed data, such that, with
\(x=x_{Q,M,\delta}\),
\[
 \ell_i^Nx\in L_i^{\rm sup},\qquad
 \ell_i^N R(x)\in R(L_i^{\rm sup}).
\]

This holds for every allowed \(Q\) away from the norm product of
\(s\), every admitted depth \(M\), and each \(\delta\) in the resulting
finite ray quotient. The same coefficient
choices give the prescribed residual Galois representations on the
averaged row.
\end{proposition}

\begin{proof}
\proofstep{Compare the real and theta-series realizations of the rational row.}
The real and series realizations of the initial member use the same marked lattice, point evaluations and finite Hecke action. At index \(i\), lift \(j\) and sign \(b\), let
\(O_i\) be the coefficient DVR, \(F_i\) its fraction field
and \(\ell_i\) its residue characteristic. Denote the signed
rational series map by \(R\), with values in the complex
coefficient functions. The matched rational ray-level test
applies to the base, twisted and universal series over the
split deletion places.

The ray-level test uses an orthonormal frame on the
chosen test space and a common nonzero normalization scalar.
One can move the frame to obtain this test. The moved and
original readings agree after the phase and coordinate changes,
so rationality of the three series transfers to the original
reading. In the zero cusp space every test reads the zero vector. Thus if \(A,B\) denote the two readings,
\[
 A=B=0\quad\Longrightarrow\quad A=cB\quad\text{for every }c\ne0,
\]
so the normalization can still be chosen nonzero.
\label{cor:automorphic-realization-46}

\proofstep{Clear the universal coefficients before changing the level or depth.}
A single denominator must work across the allowed levels and depths.

Fix a nonzero index, a lift, a sign and one admitted theta datum
\((s,h,P)\), consisting of its support, adelic argument and harmonic
polynomial. The coefficient estimate requires an integer \(c\)
for this datum. It must apply to every auxiliary level \(Q\) in the
allowed range whose places lie away from the integer norm product
of \(s\), every auxiliary exponent \(M\), and every group-algebra
coefficient.

For a harmonic
polynomial whose evaluation is a generator of the small-field
coefficient lattice, the integral core calculation supplies a
nonzero factor \(u\) and a power of \(\ell_i\). Multiplying
the coefficient by their product gives a sum of evaluations
of integral core forms in the supremum lattice over \(O_i\).
The two constants have different dependencies:
\[
 u=u(i,j,b,s,h)\ne0,\qquad
 c_\gamma=c_\gamma(i,j,b,s,h,P_\gamma).
\]
They are uniform in the admissible auxiliary level, the auxiliary
exponent and the coefficient of the group algebra.
Lift the finite integral expansion of the evaluation of \(P\).
Injectivity of harmonic evaluation gives
\[
 P=\sum_\gamma a_\gamma P_\gamma,\qquad
 c=\max_\gamma c_\gamma,\qquad
 \ell_i^c u\,x(P)=\sum_\gamma a_\gamma\ell_i^{c-c_\gamma}
                              \bigl(\ell_i^{c_\gamma}u\,x(P_\gamma)\bigr).
\]
Here \(x\) denotes the coefficient evaluation before normalization.
The right side is integral for each allowed \(Q,M\) and coefficient.
For \(P=0\) no clearing is necessary. \label{cor:automorphic-realization-47}

Dividing by the common normalization \(\nu\ne0\) gives
\[
 \lambda R(x)=R(x'),\qquad \lambda=\ell_i^c u/\nu,\qquad
 x'\in L_{i}^{\rm sup}.
\]
The \(R\)-images of an \(F_i\)-basis are complex linearly
independent. Hence \(x'_k=\lambda x_k\) in each coordinate.
For any nonzero coordinate this gives
\[
                 \lambda=x'_k/x_k\in F_i
\]
in the chosen embedding. In this case choose \(N\) with
\(\ell_i^N/\lambda\in O_i\); then \(\ell_i^Nx=(\ell_i^N/\lambda)x'\)
is integral for every admissible coefficient vector \(x\). If \(\lambda\notin F_i\), every such vector vanishes. Thus the exponent may be chosen as
\[
 N=0\qquad(\lambda\notin F_i).
\] The exponent depends
on the datum, index, lift and sign and is common to all \(Q,M\)
and coefficients. \label{cor:automorphic-realization-48}

For the full Hecke preservation, take an allowed split prime \(P\)
away from the two type primes and its normalized finite double coset
on the lower row. Use a coefficient valuation of rational
characteristic different from that of \(P\). The framed evaluation
identifies the double-coset action with the corresponding transform
of cusp sections, since their evaluations at the test points agree and
the evaluation map is injective.

For the degree \(a\) double coset, let \(c=q_P^{-a n_V/2}\).
In the geometric frame the scalar and its adelic reading are
\[
  u_{\rm geom}=c(\mathrm NP)^{-a n_V/2},\qquad
  u_{\rm ad}=c\,\mu_0(\varpi_P)^{-a}.
\]
Lemma~\ref{lem:marked-semistable-local-ball} proves the pointwise bound in the geometric frame, using the transported integral generator on each branch. The finite action commutes with the Hecke sum, since it changes the markings on its source and target simultaneously. Intersecting the extended ball with the type or good space gives the asserted preservation.

\label{cor:automorphic-realization-49}

The Galois input begins with the regular algebraic polarized
attachment just used in block uniqueness. For a regular
algebraic real norm shift \(\nu\) of a conjugate self-dual
cusp \(\pi\), use the representation of \(\pi\nu\).
Its labelled de Rham weights are those of the shifted infinity
type with the rank normalization. At a place \(w\mid\ell\) with spherical local component, the unramified case of the stated attachment identifies its local Galois representation as crystalline:
\[
 \dim D_{\rm cris}(\rho|_{G_w})=\dim\rho ,
\]
with the dimensions taken over the respective base coefficient fields. For the potential crystallinity
test, suppose that the local parameter has zero monodromy
and the stipulated local type. First take a conjugate
self-dual algebraic shift of \(\pi\). For an integral norm
exponent take the trivial character.

For a half-integral norm
exponent use the following construction on the marked CM
extension \(E/F\) with the chosen CM type \(\Phi\).
This extension is unramified at every finite place. Let
\(\eta\) be its quadratic idele class character on \(F\). Its local values in the unramified finite normalization are
\[
 \eta_v(\varpi_v)=
 \begin{cases}1,&v\text{ split},\\-1,&v\text{ inert} .\end{cases}
\]
The real factors are sign characters.
The CM type and field supply, for every root of unity \(\zeta\) in
\(E\), an element \(\delta\in E^\times\) for which
\[
 \frac{\delta}{\bar\delta}=\zeta,\qquad
 v_P(\delta)\text{ is even for }P\text{ inert over }F,\qquad
 \prod_{\phi\in\Phi}\frac{\phi(\delta)}{|\phi(\delta)|}=1 .
\]
For \(\zeta\ne-1\), take \(\delta=1+\zeta\); explicitly
\[
                   \frac{1+\zeta}{1+\zeta^{-1}}=\zeta .
\] The marked CM type has a fixed-point-free involution
induced by a nonidentity Galois involution in its indexing
subgroup that inverts the roots of unity. The phases in each
pair are complex conjugates, so their product is one.

For \(\zeta=-1\) take the marked square root of \(-d\).
It has phase \(i\) at each member of the CM type, whose
cardinality is a multiple of four. Both choices have
valuation zero away from the rational conductor of the host.
For the square root this follows from \(d\) dividing the
conductor. For \(\zeta\ne-1\), a prime away from the conductor cannot divide \(1+\zeta\). Otherwise injectivity of reduction on the relevant roots of unity would give
\[
 \bar\zeta=\overline{-1}\quad\Longrightarrow\quad\zeta=-1 .
\] The prime divisors of the
conductor split in \(E/F\), giving the inert valuation condition.
\label{cor:automorphic-realization-50}

On ideles with unit finite components prescribe the product of
the phases at \(\Phi\). On principal ideles prescribe one, and on
finite lifts of uniformizers of \(F\) prescribe the values of
\(\eta\). To verify compatibility, let \(k\in E^\times\) and let
\(m_v\) be the exponents of the chosen uniformizers. Since all ramification indices here are one, the valuation of the product at \(P\mid v\) is the sum of the scalar valuation and the prescribed exponent. The unit condition is
\[
 v_P(k)+m_v=0\quad\Longleftrightarrow\quad v_P(k)=-m_v\qquad(P\mid v).
\] Thus \(k/\bar k\) is a unit of absolute value
one at every embedding. Kronecker's theorem and the preceding
choice of \(\delta\) give
\[
 k/\bar k=\delta/\bar\delta,\qquad a=k/\delta\in F^\times,\qquad
 m_v\equiv-v(a)\pmod 2\quad(v\text{ inert}).
\]
The split uniformizers have assigned value one.

Write
\(\operatorname{ph}_\Phi(x)=\prod_{\phi\in\Phi}\phi(x)/|\phi(x)|\).
The intersection test becomes the quadratic product formula
\[
 \prod_v\eta_v(\varpi_v)^{m_v}
 =\eta_{\rm fin}(a)^{-1}=\eta_\infty(a)
 =\operatorname{ph}_\Phi(a)=\operatorname{ph}_\Phi(k).
\]
The character so defined on the generated open subgroup extends
to the idele group by divisibility of the circle group. It is
continuous because its restriction to the open subgroup is
continuous. It is trivial on principal ideles, unramified at
finite places and restricts to \(\eta\) on the ideles of \(F\).
On an idele \(x\) the product evaluates the character on a norm; triviality of \(\eta\) there gives
\[
 (\xi\xi^c)(x)=\xi(x)\xi(x^c)=\xi(xx^c)=1 .
\]
The factor \(z/|z|\) at a member of \(\Phi\) has exponent
\(1/2\) there and \(-1/2\) at its conjugate, in the convention
of labelled exponents.
\label{cor:automorphic-realization-51}

Call the chosen character in either parity \(\xi\).
Its exponents differ from those of \(\nu\) by integers at each
label. It follows that \(\pi\xi\) is regular algebraic, since
the change is a common integral shift within each infinity
parameter. The quotient \(\nu/\xi\) is an algebraic Hecke
character unramified at finite places. Its rank-one Galois
character is crystalline above \(\ell\).
At the specified local type in rank at least two, local-global
compatibility gives zero monodromy on the de Rham attachment
to \(\pi\xi\). It is potentially semistable and hence potentially
crystalline when the monodromy vanishes. The twist back to the
original attachment satisfies
\[
 r_\ell(\pi\nu)\simeq r_\ell(\pi\xi)\otimes r_\ell(\nu/\xi).
\]
Indeed the two semisimple sides have equal characteristic
polynomials at dense Frobenius classes. The second factor on the
right is crystalline, so the tensor product is potentially crystalline.
\cite{blggt-attachment,caraiani-monodromy,berger-differential}\label{cor:automorphic-realization-52}

The remaining Galois inputs are read at the actual indices.
The full degree label gives the trivial wild inertia
required for the level type. At conjugate places the local
conjugate self-dual comparison transports the dual parameter.
Inverting the eigenvalues at a transported tame generator
preserves separation. In rank one the split Frobenius values identify the row character with the reciprocity character. Tameness makes its restriction to the principal units trivial:
\[
 \chi_w|_{1+\mathfrak p_w}=1 .
\] The marked signature comparison and the
toric datum join these properties in the cohomological and
toric hypotheses of the export. Here the original rank and
the two type primes satisfy \(r\ge100\) and \(g_k\ge2r\).

The toric denominator uses
integral normalized representatives at the tested prime and
the finite norm calculation of
Section~\ref{subsec:fitted-square-ratios}. The same class-zeta
bound applies to each normalized class term and each class
sum twisted by a ring-class character. The added level consists
of split places prime to the base conductor and away from the
tested characteristic. This gives the threshold at the
given index before choosing the added level. The required roots of unity
come from the joint coefficient choice for the two type
primes. Thus the signed conditions use the row, series and
character identifications of the averaged bases.
\label{cor:automorphic-realization-53}
\end{proof}

\subsubsection*{Coefficient fields and character data for the second twist}

The averaged theta bases already realize the assigned characters for the initial rational rows. We now compare the coefficient fields and character data required by the second twist, using those rows and their period values.
For an index \(i\) over a prime \(\ell_i\), a lift
\(1\le j\le m\) and a sign \(b\), the relevant theta vectors
are the class averages of the signed theta series in the
type-selected Hecke space. The rank used to choose the finite
set of vectors is the rank of this span. The twist theorem uses
the full-rank assertion to identify the span with the rational
reading of the type-selected subspace.

\begin{proposition}[Finite basis for the retained retwisting row]\label{prop:retwisting-row-basis}
For an index \(i\) over \(\ell_i\), a lift \(1\le j\le m\) and
a sign \(b\) in Proposition~\ref{prop:signed-theta-data}, the class averages
of the signed theta series generate a finite free lattice
\(L_{i,j,b}\) over \(O_i\). One can select class averages \(a_1,\ldots,a_d\) that give an integral basis:
\[
 L_{i,j,b}=\bigoplus_{k=1}^{d} O_i a_k.
\] For the nonbase tuple, the tied
range and faithfulness hypotheses identify its rational
span with the good subspace used by the retwisting theorem.
\end{proposition}

\begin{proof}
\proofstep{Use the integral comparison before extending the fraction field.}
The
finite-class construction gives the integral comparison of the
class sums with the Hecke coefficient module. The roots of
unity of orders \(g_0(g_0-1)\) and \(g_1(g_1-1)\) are in the
embedded field \(K_0\) for the coefficient values at the type
primes. Write \(O_i\) for the local coefficient DVR. Take the
integral comparison at trivial twist and empty auxiliary level.
The \(O_i\)-span of the raw series injects via \(u\mapsto 1\otimes u\)
into the image of the finitely generated multiplicity lattice under
the theta map over the trivial-twist extension ring.

For the finite faithfully flat extension ring \(B/O_i\), the descent map is injective:
\[
 U\hookrightarrow B\otimes_{O_i}U,\qquad u\longmapsto1\otimes u .
\]
The multiplicity image in which the raw series are read is finitely generated over \(O_i\). Since \(O_i\) is
Noetherian, the raw span is finitely generated. Its image under
the integral averaging operator is consequently a finitely generated
torsion-free \(O_i\)-module \(L_{i,j,b}\).

\proofstep{Choose a basis among the reductions, then lift by Nakayama.}
Write \(L=L_{i,j,b}\), \(k_i=O_i/\mathfrak m_i\) and
\(d=\dim_{k_i} L/\mathfrak m_iL\). Choose \(d\) averaged vectors
whose reductions form a basis. Their column map satisfies
\[
 O_i^d\xrightarrow{\,C\,}L,\qquad
 k_i^d\xrightarrow[\sim]{\,\bar C\,}L/\mathfrak m_iL .
\]
Nakayama gives surjectivity of \(C\). Since \(L\) is finite and torsion-free over the DVR, it is free of
rank \(d\). The surjection \(C\) fits into the exact sequence
\[
       0\longrightarrow \ker C\longrightarrow O_i^d
              \xrightarrow{C} L\longrightarrow0 .
\]
If \(K\) denotes the kernel over the coefficient DVR \(O\), its rank-zero and torsion-free conditions together give
\[
 K\hookrightarrow K\otimes_O\operatorname{Frac}(O)=0 .
\]
This also gives the empty basis when \(d=0\).

\proofstep{Identify the rational span using the actual signed theta range.}
The second argument determines the rational rank of \(L_{i,j,b}\)
inside the signed Hecke space. The real-side map from each signed
finite model is injective under the rational row comparison.
The nonzero scalar multiplying its theta reading and the chosen
real vector give the same coefficient reading on the rational
form.

At a nonbase index, let \(V_{\rm good}\) denote the good complex space at the indicated base reading. The range hypothesis and the finite class sum give 
\[
 V_{\rm good}=\operatorname{span}_{\mathbb C}(\text{projected lifts})\subseteq\operatorname{span}_{\mathbb C}(\text{averaged series}).
\]
 For the base reading the projected image is the good row by Proposition~\ref{prop:span-original-theta} in its stated stable range. 

At a nonbase reading the marked comparison must identify the projected lifts and their signed class sums on the same space. On a tied good eigenspace \(V\), with lower summand \(U\) and lift \(T:U\to V\), Proposition~\ref{prop:downward-on-good-space} gives
\[
 V\ominus T(U)=\ker(T^*|_V)=0
\]
in the rank range of Lemma~\ref{hyp:tied-adjoint-level}, using \(T^*V\subseteq U\). Thus no choice of a basis of the integral lattice by itself provides the rational spanning assertion. 

Over the rational form giving
the complex reading, the dimension of this span is the largest
size of a nonzero minor of a coefficient matrix. For the coefficient matrix \(C\) and local embedding \(\tau\), apply \(\tau\) to each minor:
\[
 \det\tau(C_{I,J})=\tau(\det C_{I,J})\quad\Longrightarrow\quad\operatorname{rank}\tau(C)=\operatorname{rank}C .
\] This proves full rank
of \(L_{i,j,b}\). The integral index in the Hecke lattice
continues to enter the twist estimates.
\label{cor:automorphic-realization-54}

Use these bases in the all-coprime-powers choice of
Section~\ref{subsec:retained-twist}. The type comparison
above supplies the universal coefficient membership. Use the
clearing powers and full Hecke preservation with the Galois
and toric hypotheses on this same initial system. For a
nonzero index \(i\) served by a nonbase tuple, the output is
\(N_0(i)\), a level \(Q_i\) and a character with values in
\(\mu_{\ell_i^{k_i}}\), \(k_i\ge1\). The quantifier on the signed
character and valuation tests is
\[
 \exists\, Q_i,k_i,\varphi_i\quad
 \forall a\in\mathbb Z,\quad (a,\ell_i)=1
       \ \Longrightarrow\ \varphi_i^a\ \text{satisfies both tests}.
\]
Before each choice put all host primes over the exclusions into
the avoidance set, including the two deletion primes, the sequence
primes, the residual test prime and the original level primes.

For successive indices also exclude all rational primes used
earlier. The auxiliary places have degree one; this gives disjoint
levels even after taking their rational norms.
To choose the index at a coefficient place, group the contracted
primes of \(K_0\) into orbits under conjugation. The marked
two-place covering provides the original index and signed
valuation on the fixed coefficient subfield at each place. Let \(\mathcal O\) be an orbit of contracted primes meeting the nonbase support. Choose one supported place \(r\) over the orbit, then use the single index
\[
 i(w)=i_{\rm orig}(r)\qquad(w\text{ supported},\ w|_{K_0}\in\mathcal O).
\] The
export gives the full rank and the basis on that very index,
since \(r\) is assigned a nonbase tuple.

At places contracting
to the same prime as \(r\), its original side recovers the same
valuation tests on \(\mathcal O_{K_0}\). At the conjugate prime
use the conjugate side. Its embedding is the first embedding
composed with conjugation. If the contracted prime is fixed,
the representative's side already gives the valuation tests.
At each place off the support keep its original index and
signed reading, with trivial auxiliary character. Thus there
is a reading at every place. The index is constant under
conjugation and on each contraction fibre on the support.
\end{proof}

\begin{auxlemma}[Transport of the reference coefficients]\label{lem:reference-field-transport}
Let \(K_s\subset K_0\subset K\) be the three coefficient fields of
the reference calculation, embedded by \(\iota:K\hookrightarrow\mathbb C\).
Let \(\Xi\) be its finite set of selected Hecke characters. Suppose a
finite set \(S(\Xi)\subset\overline{\mathbb Q}\subset\mathbb C\) of
additional character values, roots of unity and ray-field generators
is required. There is a finite extension \(K'/K\), with embedding
\(\iota'\) extending \(\iota\), containing these prescribed complex
values. The embedded copies of \(K_s,K_0\) may be kept equal. If
conjugation stability is required, assume that \(\iota K\) is stable
under complex conjugation. In this case a possible image is
\[
    \iota'K'=\iota K\bigl(S(\Xi),\overline{S(\Xi)}\bigr)
                           \ \subset\mathbb C .
\]

Let \(\mathcal H\) be the Hecke algebra away from the excluded
places, let \(V\) be the reference space over \(K\), and give it and
its two marked Hecke actions their original embeddings. If
\(e\in\mathcal H\otimes K\) is either prescribed reference idempotent
and \(\lambda\) a character acting on a reference vector, their copies
on \(V\otimes_K K'\) satisfy
\[
  (eV)\otimes_K K'=(e\otimes1)(V\otimes_KK'),\qquad
  \lambda'(T)=\iota'^{-1}\iota(\lambda(T))
                  \quad(T\in\mathcal H).
\]
These identities apply to the reference summands carried to the new
construction.
\end{auxlemma}
\begin{proof}
Inside \(\mathbb C\) take the compositum of \(\iota K\) and the fields
generated by the finitely many specified elements, and also by their
conjugates when required. Give the two smaller fields their original
embeddings. The required inclusion is the unique one satisfying
\[
    j=(\iota')^{-1}\circ\iota:K\longrightarrow K'.
\]
Tensor the direct sum \(V=eV\oplus(1-e)V\) with \(K'\). On a
reference eigenvector \(v\) one has
\((T\otimes1)(v\otimes1)=(v\otimes1)j(\lambda(T))\).
A reference evaluation expressed in a fixed cusp basis has the same
complex coordinates after replacing each coefficient \(a\) by \(j(a)\).
When \(w\) lies above \(v\), contraction of prime ideals gives
\[
    j^{-1}(\mathfrak p_w)=\mathfrak p_v,\qquad
    \operatorname{ord}_w(j(a))=e(w/v)\operatorname{ord}_v(a).
\]
Consequently an original vanishing condition modulo \(v\), or a
nonzero residue, pulls back to the corresponding condition modulo
every \(w\) above it. For a marked set specified by these original
conditions, the resulting set of places is the entire inverse image.
\end{proof}

\begin{auxlemma}[The reference datum after adjoining the ray characters]
\label{lem:second-core-reference}
Fix the Hermitian spaces, lattices, levels at the word primes, split
halves and deletion characters at one translate as in
Definition~\ref{hyp:initial-adelic-frame}. Suppose the preliminary
realization at this translate has been constructed. Write
\[
 K_s\subset K_0\subset K\subset\mathbb C,\qquad
 \iota:K\hookrightarrow K'\subset\mathbb C
\]
for the distinguished fields and the finite enlargement of
Proposition~\ref{prop:second-specialization-row}. Keep the images
of \(K_s,K_0\) fixed in \(\mathbb C\). The geometric and adelic
data of the reference tuple admit a presentation over this
enlargement with the original small field, marked points, class
representatives, oscillator, measures and analytic theta maps.

For a coefficient-valued reference class form \(\phi\), denote the
resulting form by \(\phi^\iota\). For its Hecke operators away
from the levels and conductors one has
\[
 \phi^\iota(h)=\iota_*\phi(h),\qquad
 T\phi^\iota=(T\phi)^\iota .
\]
Here \(\iota_*\) applies the inclusion to the coefficients of the
value polynomial. These assertions hold on the reference row and
on each copy of a common test at its own level.
\end{auxlemma}
\begin{proof}
\proofstep{Retain the geometric model.}
The PEL equations, integral structures, marked base point, component
maps and the correspondence with the analytic quotient are already
defined in the preliminary presentation. If \(F\subset K_0\)
denotes its embedded geometric coefficient field, use the map
\[
                 F\longrightarrow K\xrightarrow{\iota}K'.
\]
Both embeddings of \(F\) into \(\mathbb C\) agree. Keep the
model over \(F\), its ring of integers, its polarization, frames,
and all the maps at its original levels. Its presentation over
\(K'\) is obtained by pullback. The Hermitian groups and their
adelic quotients depend on the original host and Hermitian
spaces. We can use the original fundamental sets and Haar measures.
Thus both their orbit masses and the Petersson integrals on the
old complex forms have the original values.

\proofstep{Identify the reference level and the copied forms.}
Let \(S_{\rm type}\) be the prescribed set of type places and
\(Q(\mu)\) the conductor support used for a tuple \(\mu\).
For the reference tuple \(\mu_0\),
\[
              Q(\mu_0)=\varnothing,\qquad
              S_{\rm type}\cup Q(\mu_0)=S_{\rm type}.
\]
At these places a pair consisting of a level subgroup \(J\) and
its character \(\lambda\) becomes \((J,\iota\lambda)\).
For each definite class represented by \(h\), apply \(\iota\)
to the coefficients of \(\phi(h)\). The stabilizer and
right-\(J\) equations commute with \(\iota_*\), since the
weight action and the level character extend coefficientwise.
Thus these polynomials define the form \(\phi^\iota\). On
evaluation in \(\mathbb C\) at the distinguished embeddings,
\[
              (\iota_*\phi(h))_{\mathbb C}=\phi(h)_{\mathbb C}.
\]
For any fixed integral class lattice given by these finite
stabilizer and level equations in the same value lattices, scalar
extension to a finite extension of discrete valuation rings is
flat. It commutes with their common kernel. This observation
applies to the copied equations and their coefficient lattices.

\proofstep{Transport the operators.}
With the same oscillator vector and integration domain, the
reference theta operator is given by the same integral on the
same smooth inputs. The normalized algebraic reading agrees by
Lemmas~\ref{lem:marked-theta-fibre-comparison}
and~\ref{lem:marked-vacuum-scalar}. At a good prime a Hecke
double coset has the same right-coset representatives on the
copied level. If \(A_d\) is the coefficient action in the
chosen convention, then at each representative \(h\)
\[
 (T\phi^\iota)(h)
   =\sum_d \iota_*(A_d\phi(hd))
   =\iota_*\!\left((T\phi)(h)\right).
\]
In particular a copied nonzero eigenvector of eigenvalue \(a_T\)
has eigenvalue \(\iota(a_T)\). If a supported tuple requires a
finer level, use its forgetful map to the reference level for
the pullback of the old form. At a good prime the local factors
of the two level subgroups coincide, so the same computation
intertwines pullback and \(T\).

\proofstep{Insert a moving character at its actual level.}
For a carried tuple \(\mu\), the prescribed local level uses
\(S_{\rm type}\cup Q(\mu)\), its split halves, and the finite
level character specified there. Realize a level character on
the cover by its kernel, where it acts by deck transformations.
For an allowed fixed ray test with universal polynomial
\(\sum_{\delta\in\Delta} a_\delta[\delta]\), choose the same
test on this cover and pull back its value frames. Then for any
allowed character \(\varphi\), with values read in a common
extension,
\[
       \sum_{\delta\in\Delta}\iota_*(a_\delta)
                   (\iota\varphi)(\delta)
        =\iota_*\!\left(\sum_{\delta\in\Delta}
                         a_\delta\varphi(\delta)\right).
\]

Here \(\iota\) also denotes the chosen extension to the field
of the character values. Using a different primitive-root
identification at a place above it replaces \(\varphi\) by
\(\varphi^a\) for one \(a\) prime to its order.
Proposition~\ref{prop:retwisting-row-basis} supplies the
selected tests for all such powers. The equality compares the
same fixed test and character in two coefficient presentations.
The reference eigenvalue identity concerns the copied
eigenvector. A newly twisted eigenvector is tested with its
own finite character and level.
\end{proof}

The construction therefore starts with one preliminary realization.
It exports the chosen finite indices, ray supports, characters and
their all-powers tests. After adjoining their values the reference
data can be carried as in Lemma~\ref{lem:second-core-reference}.
At a place \(w\) of the larger field the index is obtained by
contraction to the preceding field. The prescribed character
and all-powers comparison then apply at \(w\). The final row
must realize the new local types, the corresponding integral
and real comparisons, and the estimates at its actual levels.
Once it has been so constructed, write \(D'\) for that final
member, with \(M_{w,D'}\) its column lattices.
Both division estimates use the lattices of \(D'\); the
preliminary member supplies the selection of characters. This
order uses no comparison between the division errors of two
different members.

\begin{auxlemma}[A single period on the copied rational row]
\label{lem:copied-reference-period}
Fix a served coefficient index \(i\), a lift \(1\le j\le m\)
and a sign \(b\). Use the reference realization of
Lemma~\ref{lem:second-core-reference} and its fixed extremal
vector \(v_0\). Suppose the first signed reading on its rational
row is given by a period \(\nu_{ijb}\ne0\). On the copy one may
take the very same \(\nu_{ijb}\), including for the common ray
tests. Under the fixed-field identification the resulting
reading at \(K'\) restricts to the old reading.
Writing \(R,R'\) for the two signed readings and \(j\) for transport of coefficients, this gives on every common test
\[
                         R'(jz)(g)=R(z)(g).
\]
\end{auxlemma}
\begin{proof}
Write \(\sigma\) and \(\sigma'\) for the signed coefficient
embeddings into the common complex realization, with
\(\sigma'\iota=\sigma\). On a pure tensor \(z\otimes f\) and a
marked representative \(g\), the normalized value on the copy is
\[
   \frac{\sigma'(\iota z)}{\nu_{ijb}}\,
            \operatorname{ev}_{v_0}(f^\iota)(g)
       = \frac{\sigma(z)}{\nu_{ijb}}\,
            \operatorname{ev}_{v_0}(f)(g).
\]

The evaluation vectors and the analytic kernels are the copied
ones. The identity on tensors gives equality on the original
rational span and determines its scalar extension. To compare a
common test supported at \(Q\), use the local section of each
ray label, the original Fourier weights, and the determinant
translation at \(g\). For its coordinate \(a\in\Delta_Q\), the
translated local section is indexed by
\[
                      a-d_Q(g),\qquad d_Q(g)\in\Delta_Q .
\]
This label, the input polynomial and the placed finite
representative coincide in the copied formulas. The stabilizer
and Haar factor at each matching definite source class coincide
as well. Their weighted sums give the same value for that
coordinate. Finite specialization and the normalized value
therefore commute at every eligible common test in this copy.
\end{proof}

\begin{auxlemma}[The character line at an added split prime]\label{lem:new-ray-character-line}
At an added split prime \(v\in Q\) of a retained tuple, let
\(F_v\) have ring of integers \(A\) and residue field \(k\).
Write \(\xi:k^\times\to\mathbb C^\times\) for the nontrivial tame
restriction of the chosen ray character, in the selected split half.
Use a row of rank \(r+d\), where \(d=jt\), and the inverse image \(J\)
in \(K=\mathrm{GL}_{r+d}(A)\) of the upper \((r,d)\) block parabolic.
On the conjugate automorphic input the suffix is twisted by \(\xi\).
After putting the determinant multiplier of the Weil action into the
group coefficient, the normalized induced representation and the
character to be tested there are
\[
 \pi'=\operatorname{Ind}_{P_{r,d}}^{\mathrm{GL}_{r+d}(F_v)}
       \big((\tau_0\otimes\widetilde\xi^{-1}\circ\det)
                           \boxtimes\sigma_0\big),
 \qquad \eta(j)=\xi^{-1}(\det\bar j_{00}).
\]

Here \(\widetilde\xi\) extends \(\xi\) to \(F_v^\times\), with its
unramified part absorbed in the spherical factors \(\tau_0,\sigma_0\).
The representation on the rank-\(r\) block is spherical before the
displayed twist. The rank-\(d\) factor is the spherical representation
of the preceding \(j\) determinant strings. Then
\[
               \dim_{\mathbb C}\operatorname{Hom}_{J}(\eta,\pi')=1.
\]
The formula includes \(d=0\) and every power of the retained character
prime to its order.
\end{auxlemma}
\begin{proof}
Write \(X=\left(\begin{smallmatrix}A_0&B_0\\C_0&D_0\end{smallmatrix}\right)\)
in the split oscillator model, with \(A_0\) square of size \(r\).
On integral matrices take the test
\[
 f(X)={\bf1}_{\{A_0\in\mathrm{GL}_r(A),\,C_0\equiv0\bmod\varpi\}}
                   \xi(\det\bar A_0).
\]
Right multiplication by \(j\in J\) multiplies this test by
\(\xi(\det\bar j_{00})\). The scalar of the oscillator action on units
is \(\xi(\det\bar j)^{-1}\). Their product is
\(\xi(\det\bar j_{11})^{-1}\), the suffix character in the level.
After conjugating the automorphic coefficient and moving the
oscillator scalar to that coefficient the character is
\[
     \xi(\det\bar j_{11})\xi(\det\bar j)^{-1}=\eta(j).
\]
The same computation with left inverse multiplication gives the
inverse level law on the other leg.

Take unit spherical vectors in \(\tau_0,\sigma_0\), with product \(u\).
The normalized induction on \(P_{r,d}\cap J\) acts on \(u\) by
\(\eta\), since the modulus character is one there. Thus
\[
 v(pj)=\delta_{P_{r,d}}(p)^{1/2}
             (\tau\boxtimes\sigma_0)(p)\eta(j)u,\qquad
 \tau=\tau_0\otimes(\widetilde\xi^{-1}\circ\det),
\]
with value zero off \(P_{r,d}J\), is well defined and \(v(1)=u\).

To bound the dimension, take unramified principal-series
presentations surjecting onto the two spherical factors and induce
in stages. Taking a \(J\)-character is exact by compact averaging.
The cells \(B\backslash \mathrm{GL}_{r+d}(F_v)/J\), for the upper Borel
\(B\), are indexed by \(r\)-element subsets of the diagonal positions.
On diagonal units the characters of the principal series are
\[
        \underbrace{\xi^{-1},\ldots,\xi^{-1}}_{r},
        \underbrace{1,\ldots,1}_{d}.
\]
Only the subset of the first \(r\) positions can match \(\eta\),
because \(\xi\ne1\). On that cell a section is determined by a single
scalar. This also proves the assertion when the suffix is empty.
A power prime to the order leaves the inertia character nontrivial.
\end{proof}

\begin{auxlemma}[Occurrence on the signed row after specializing the ray character]\label{lem:new-ray-type-occurrence}
Use the signed row of Proposition~\ref{prop:retwisting-row-basis} and the
coefficient tests of Lemma~\ref{hyp:simultaneous-retained-twist}.
Fix a served index, a lift \(1\le j\le m\) and either side. Let
\(\varphi\) be its retained character and \(a\) prime to \(\ell\).
Let \(V\) be the good rational span in that row. On the prescribed
finite-class realization at the tuple \(\varphi^a\), denote by \(X_a\)
the character space of source forms and by \(L_a:X_a\to V\) the theta
class sum followed by the good projection, extending scalars as needed.

Then \(L_a\) is onto. In its complex Petersson realization, for
each joint bottom eigencharacter \(\lambda\) of the spherical operators
outside the exceptional places where the attached local theta
transfer is specified,
\[
 L_a^*(V[\lambda])\ \subseteq\ X_a[\lambda_a],\qquad
 \dim X_a[\lambda_a]\ \geq\ \dim V[\lambda],
\]
where \(\lambda_a\) is obtained by that transfer. The assertion also
holds after imposing the common finite type idempotent intertwined by
the two theta actions. Thus every occurring bottom eigensystem has a
nonzero vector on the prescribed upper character space.
\end{auxlemma}
\begin{proof}
Put \(n=\dim V\). Proposition~\ref{prop:row-pair-removal} supplies a basis
of the row lattice and a common clearing scalar \(x=\ell^c\).
For the \(n\) selected source class generators \(z_i(a)\) at this
tuple the specialized matrix in that basis satisfies
\[
       C_Q(u^a-1)=
         [\,xL_a z_1(a)\ \cdots\ xL_a z_n(a)\,],\qquad
       \det C_Q(u^a-1)=f_M(u^a-1)\ne0 .
\]
Use the inverse tuple on the side defined with the inverse character.
The nonvanishing for both sides and all these powers is the simultaneous
choice of Proposition~\ref{prop:one-retained-character}.
Lemma~\ref{lem:universal-theta-type} gives the signed type on each
specialized column. For the source character law, translate the
class representative in its finite sum by a level element: the
change in its determinant index multiplies its coefficient by the
prescribed character, and its polynomial is transported by the
coefficient representation. These are precisely the transformation
laws defining \(X_a\). The columns span \(V\), including the empty
case \(n=0\).

For a tied arrow let \(h\) be an upper spherical operator and
\(\mathcal T(h)\) its lower transferred operator. The good projection
is orthogonal and commutes with this algebra. The downward
intertwining law on the kernel and the Petersson adjoint give
\[
 L_a h^*=\mathcal T(h^*)L_a,\qquad
 h L_a^*=L_a^*\mathcal T(h).
\]

Indeed the local transfer respects the involution. In normalized
Satake coordinates the involution conjugates and inverts the
parameters. The splitting characters are unitary and the appended
string is fixed under this operation, by \(j'\leftrightarrow t-1-j'\).
At an unramified place the local transfer on normalized parameters of
one upward arrow, with its oriented splitting characters
\(\chi_S,\chi_T\), is the multiset map
\[
 B_w\longmapsto
     (\chi_S/\chi_T)(\varpi_w)B_w
       \ \sqcup\
     \{\chi_S(\varpi_w)({\rm N}w)^{\,j'-(t-1)/2}:0\le j'<t\}.
\]
It is injective by multiset cancellation. At inert hyperspecial
places this formula reads the base-change parameter over \(E_w\);
the Hecke operator is the relative spherical operator over \(F_v\).

Iterate the formula along the tied arrows. At the residue-type
places simultaneous translation of the two copies of the first
block in the Weil kernel intertwines the common finite projector.
The spherical operators commute with this projector by disjoint
support. Equivalently the latter statements can be imposed first
in forming \(X_a,V\).

If \(v\in V[\lambda]\), the intertwining identity gives
\(hL_a^*v=\lambda(\mathcal T(h))L_a^*v\). Surjectivity gives
\[
          \ker L_a^*=(\operatorname{im}L_a)^\perp=0 .
\]
This proves the asserted lower bound with the same finite test
for each vector of the bottom eigenspace. When several full labels
have the same restriction to the tested operators, the argument
applies to their direct sum. It supplies a bound on the whole
restricted eigenspace. It gives no numerical multiplicity for a
particular representation inside that eigenspace without a projector
separating that representation. In the comparison of the two signed
coefficient models alone, the tested operators have split primes
\(P\notin S_0(i)\) and degrees \(1\le k\le r\); the signed
normalization multiplies their degree-\(k\) eigenvalue by
\(c_{i,b,P}^{\,k}\).
\end{proof}

\begin{proposition}[Preservation of the row under the second specialization]\label{prop:second-specialization-row}
In the two-pass construction, retain the embedded fields
\(K_s,K_0\), the selected character data and the assigned
word levels. Extension of the coefficient field and the second
specialization preserve the transported row and the tuple
character at each marked place and the grid tests at every
admitted depth. If \(w\) lies over an original marked place \(v\), its tuple character is inherited from \(v\). In particular
\[
 \{\chi_{v(w)}:w\text{ over the marked places}\}=\{\chi_v:v\text{ marked}\},
\]
and the same tuple list includes the distinguished characters \(1,\delta,\delta'\).
\end{proposition}

\begin{proof}
Carry the averaged bases of Proposition~\ref{prop:retwisting-row-basis}
into this comparison, with their level type from
Lemma~\ref{lem:universal-theta-type}.

Keep the embedded fields \(K_s,K_0\) and the character frame.
For each selected index choose a complex identification of its
cyclic \(\ell_i^{k_i}\)-th roots and use class field theory for
its ray quotient to form the complex finite-order Hecke
character. At a prime outside \(Q_i\), let \(\gamma\) be
the image of its Frobenius in the ray quotient, and let \(j_i\)
be the chosen embedding of its character values. In the given
Artin convention,
\[
 \chi_i(\mathfrak p)=j_i(\varphi_i(\gamma))^{-1},\qquad
 \chi_i^c=\chi_i^{-1}.
\]
The second equality follows from conjugation on the ray character.
The kernel is open because the quotient is finite and continuous.
Injectivity of \(j_i\) preserves the nontrivial inertia at each
place of \(Q_i\); the character is unramified elsewhere, and
\(Q_i\) avoids the primes over \(\ell_i\).

For every original
coefficient place, use the character for its assigned index
on the prescribed nonbase support, and use the trivial
character on the complement. A new place of any eventual
coefficient extension then inherits the complex character,
level and root order from the place below it. If \(w\) is an upper prime and \(K\) the lower field in the contraction, compatibility of the field inclusion with conjugation gives
\[
 (\bar w)\cap K=\overline{w\cap K}.
\]
Under contraction,
places with the same prime of \(K_0\) have contractions with the
same prime of \(K_0\). The saturation of the initial support
and the constancy of the index on this support show that
these assignments agree there, including on the complement.
The conjugate pair of local readings inherits the same index.
The rational prime set of the marked data is unchanged.

Let \(S\) be the set of the three distinguished characters and the
deletion-character products from the assigned places. Contraction
\(\pi\) is onto, so for the assignment \(v\mapsto\chi_v\) one has
\[
 \{\chi_{\pi(w)}:w\text{ in the extension}\}
                 =\{\chi_v:v\text{ in the original field}\}.
\]
Thus \(\Xi'=S\) is known before adjoining roots. Adjoin to \(K\)
the roots of orders \(\ell_i^{k_i}\), the values required by
\(\Xi'\), and \(S_2(K_s,K_0,\Xi')\). Only finitely many places
lie over the prescribed rational primes; at the other places
the assigned root order is one. Choose finitely many generators
of \(K/\mathbb Q\).

Let \(A\) consist of these generators, the requested entries and
roots, and \(i\). The required enlargement is
\[
 K'=\operatorname{Spl}_{\mathbb Q}
          \left(\prod_{a\in A}\operatorname{minpol}_{\mathbb Q}(a)\right)
           \subset\mathbb C .
\]
It is finite and normal, and complex conjugation permutes the roots
of every factor. Hence it acts on \(K'\).
The inclusions retain both coefficient embeddings: if \(i'\) denotes the inclusion of the enlarged field, then
\[
 i'|_{K_s}=i_s,\qquad i'|_{K_0}=i_0,
\]
where \(i_s,i_0\) are the original maps into \(\mathbb C\). If an additional field requirement concerns the
reconstruction of the adelic datum or finite coefficient matrices,
its finite list must be included in this step, with its
additional avoidance set imposed before selecting the ray
characters. The interval caller itself uses empty additional
avoidance sets. \label{cor:automorphic-realization-55}

There is a root identification to check when comparing the
old character with the rebuilt one at \(w\) above \(v\).
Write \(i=i(v)\), \(n=\ell_i^{k_i}\) and use the transported ray
quotient with both marked maps on \(K_0\). In the common local
fraction field let \(j_1,j_2\) read the old and new roots. Both embeddings account for all \(n\) distinct roots of \(X^n-1\) in the common local fraction field \(L\). Thus
\[
 j_1(\mu_n)=\mu_n(L)=j_2(\mu_n).
\] For one
primitive \(\zeta\) there is therefore a unit \(a\bmod n\) with
\[
 j_2(\zeta)=j_1(\zeta)^a,\qquad
 j_2(\zeta^k)=j_1(\zeta^k)^a\quad(k\in\mathbb Z).
\]
The same exponent works on the whole ray-character image.

Both the new arithmetic character and the old ray character
read the same complex Hecke character at Frobenius outside
\(Q_i\) and the residue characteristic. For an arbitrary element
\(\gamma\) of the finite ray quotient, surjectivity of the Artin
map gives a preimage in the Galois group of the ray extension.
The abelian Chebotarev theorem realizes \(\gamma\) by a
Frobenius at a place outside those exclusions.

At that place
each reading uses the inverse of the complex Hecke value, so
the common exponent gives
\[
             \varphi^{\mathrm{new}}_{w}(\gamma)
                         =\varphi_i(\gamma)^a
\]
after identifying the common coefficient rings. Equality in
the fraction field gives equality in the root ring by
injectivity. Hence the rebuilt tuple of index, level, root order
and raw character is the original tuple with its character
raised to a power prime to \(\ell_i\). The exponent may vary
with \(w\), and the all-powers twist bound covers each new place.
For the base support use the trivial auxiliary character.
\par
For the scalar extension of a fixed coefficient lattice, the
length comparisons use normalized valuations. If a finite
length module over the first local integer ring has elementary
divisors \(\pi_v^{a_1},\ldots,\pi_v^{a_h}\), their orders after
extension are \(e(w/v)a_1,\ldots,e(w/v)a_h\). Transitivity of ramification indices gives the equality of
normalized lengths explicitly:
\[
 \frac{\sum_j e(w/v)a_j}{e(w/\ell_i)}
       =\frac{\sum_j a_j}{e(v/\ell_i)} .
\] The same calculation applies to the
valuation of a scalar. Applying this observation to an
automorphic comparison requires identifying the actual
transported lattice and period on both sides.
\label{cor:automorphic-realization-56}

The rebuilt adelic datum uses the same column frame,
choice of split places and character coordinates and the
extended finite readings. Its bare row is built on the
extended fine datum. The same-Galois construction and the
rational and cuspidal readings give the intermediate state
there. The first pass fixes the moduli point, its signed integral series and its eigensystem. For each good Hecke operator \(T\) at a split place away from both levels, write \(H_T\) for the finite level correspondence there. For a form \(a\) the test on a marked point \(z\) has the double-coset description
\[
             (Ta)(z)=\sum_{h\in H_T}a(zh)
\]
with the same normalization (and coefficient action) on both levels. Thus pullback of the values along the marked correspondence intertwines the two Hecke operators. If the marked value tests separate the form, equality of these evaluations gives equality of the carried eigensystems. 

Lemma~\ref{lem:placed-row-theta} gives the marked points. The tests and their integral evaluations on the reference space are transported as in Lemma~\ref{lem:constructed-signed-coefficient-row}. At the chosen new type, Lemma~\ref{lem:new-ray-type-occurrence} supplies the eigencharacter on the tested algebra. The two bare Galois
frames agree by transitivity, through the comparison
with the bare source in each construction.
\label{cor:automorphic-realization-57}

Write \(i',i\) for the two fixed-field embeddings. Equality of their images defines the identification by
\[
 K'_0\xrightarrow{\ i^{-1}\circ i'\ }K_0 .
\] This identification respects
conjugation and integer rings. For the place map \(\pi\) and an
integer \(a\in\mathcal O_{K_0}\) the local comparison reads
\[
 a\in\mathfrak p_{\pi(v')}\ \Longleftrightarrow\
   a\in\mathfrak p_{v'}\cap\mathcal O_{K_0},\qquad
 S_{\rm new}=\pi^{-1}(S_{\rm old}) .
\]
If \(S_i\) denotes the old support at index \(i\), the transported
support is \(\pi^{-1}(S_i)\). Since contraction is surjective,
\[
                \pi^{-1}(S_i)\ne\varnothing
                     \quad\Longleftrightarrow\quad S_i\ne\varnothing .
\] \label{cor:automorphic-realization-58}

These comparisons give a copy of the original
eigensystem in the new intermediate state. It transports
the twisting, scalar and theta data after the finite
group and trivial-tuple projector have been identified.
The series on the copy uses the comparison of the
Hecke readings after restriction at the new types.
The positive multiplicity and signed level comparisons
are made on this series and give the new second
intermediate level. They refer to the copied
eigenvalues, the same range of lifts and the copied
unipotent character. They also supply the conductor
and level data for the two signed readings of the
appropriate local twist.
\label{cor:automorphic-realization-59}

\proofstep{Keep the selected place and transport the old level tests using the same support.}
The field comparison must preserve the selected place at which the row tests are imposed.

Let \(p_*\) be the least split auxiliary prime specified for the grid and \(v_*\) its prescribed base place. This choice meets both exclusions simultaneously:
\[
 v_*\notin S_{\rm grid}\cup \bigcup_\mu\operatorname{supp}(\operatorname{cond}\chi_\mu),
\]
where \(S_{\rm grid}\) denotes the bad places of the grid.
Use the compatible
place on the enlarged frame, with the same selected
host place above it. The field values and required
roots lift through the inverse of the fixed-field
identification. Since the distinguished and deleted
character coordinates agree, the formula for each
grid character agrees. The old level test at a supported
index applies at its new presentation by the equality
of support just proved. This gives transport of the list
at every depth for which the old list is given.
\label{cor:automorphic-realization-60}

Let \(\pi\) denote contraction on coefficient places and \(i(v)\)
the exported index. At the geometric point of the carried series
the character choice is
\[
                 i_{\rm new}(w)=i(\pi w),\qquad
                 i_{\rm new}(\bar w)=i_{\rm new}(w).
\]
The first equality uses the pulled-back support and its allowed
reading; the second uses conjugation on the original assignments.
The geometric comparison carries the grid, and spectral
specialization preserves this index. Hence its signed tests use
the exported signature and basis. The first specialization supplies
these data from the original real and series readings and the
second receives them on the transported series.
\label{cor:automorphic-realization-61}

For the tuple set \(U\) on this system, injectivity of the tuple-to-character map and the description of its image give the bijection
\[
 U\xrightarrow{\ \sim\ }\Xi',\qquad \mu\longmapsto\chi_\mu .
\] Its distinguished tuples read
\(1,\delta,\delta'\), and its assigned nonbase tuples read the
products just constructed. Off the selected support the
auxiliary character is trivial. On the base tuple the two
conjugate readings have twisting order zero. The mass
formula of Section~\ref{subsec:selected-word-level} returns
the original \(d_{u,*}\) and \(D_*\) on this assembled member.
The volume-to-mass identity uses its measured row, and its
signatures and Weyl weights are the prescribed ones.

For the
stabilizer divisor use its single value across words and
tuples, supplied by the column comparison. The congruence
orders off the twist support vanish as above. These assignments exhaust the tuple range. On the twist support this follows from the tuple character formula; elsewhere its auxiliary factor is one. Since the tuple-to-character map is injective, equality of the two readings also identifies their labels:
\[
 \chi_{\mu_1}=\chi_{\mu_2}\quad\Longrightarrow\quad \mu_1=\mu_2 .
\] These identities extend the member
by its numerical conditions. Forgetting those conditions
returns exactly its constructed object. In particular the
pre-level masses are valid for the member obtained after both
passes at any permitted levels at the primes assigned to the sequences.
\label{cor:automorphic-realization-62}
\end{proof}

\paragraph{Substitution of the division lattices.}

\begin{lemma}[Simultaneous substitution of the division lattices]
At a fixed good member let \(M_u\) denote its rational ramified word columns. The divisor, arithmetic and period comparisons may be transported simultaneously for any assignment
\[
 u\longmapsto(N_u,h_u),\qquad N_u\subseteq M_u\text{ an integral submodule},\quad h_u\in\mathbb R.
\] For a compatible completion these are its local
lattices and heights.
\end{lemma}
\begin{proof}
Write \(M_u\) for the rational space of the ramified word
column. The substitution is made independently in every word column. Its entire input is
\[
 (N_u,h_u)_u,\qquad N_u\subset M_u\ \text{a coefficient-integer submodule},\quad h_u\in\mathbb R .
\] Substitute \(N_u,h_u\) in the division
coordinates of the member. Its Hermitian forms, local
comparison bases, positions, split levels and degree-zero column
are carried in the same coordinates. In particular their joint
reconstruction can use the old matrices, local coefficient
embeddings, bases and weight arrays with the new division
coordinates inserted in the original column datum. Its kernel
family on each selected pair uses the same class, series and
substitution matrices. The expression of its entries from
the chosen places and adapted bases is carried with the series.
The downward maps and their paired integral identities persist
with the same test, function, character and Haar measure.

Let \(\pi\) pull an index of the substituted flag datum back to
its original flag. At a flag \(a\) the quotient-prime set and
model degree satisfy
\[
 Q^{(N,h)}(a)=Q(\pi a),\qquad
 \deg^{(N,h)}(a)=\deg(\pi a).
\]
The marked chart, weights and model fibre are read at the same
\(\pi a\). Closure of the geometric condition under a single
\((N,h)\) substitution thus applies both to the initial flags
and to flags resulting from earlier substitutions. \label{cor:automorphic-realization-63}

For the arithmetic comparison use the two row realizations with their
rational coordinates, lattices, finite types and signatures.
In the revised norm choice retain the normalized character and
signed basis. The same-row comparison transfers its grid list.
If \(a\) identifies the earlier field and series with the present origin and \(b\) identifies that origin with the new one, the induced identification is
\[
 x\longmapsto b(a(x)).
\]

In the measured comparison use the saturation for the
degree-zero column just transported and its original scale and
evaluation. Its spectral data therefore transport on these
columns with the prescribed \(N,h\). The sharp normalization
and the type choices defining the logarithmic values transfer
on the carried row. The two finite sets in the twist test and the
same signed theta data give its transported test. The error
model is the one read from the carried Hermitian column.
For a fixed \((N,h)\) the transported column, comparison map, twist form and kernel are identified together. Thus the spectral and error estimates concern the same lattices and heights. \label{cor:automorphic-realization-64}

Now fix a compatible completion \(p\). Take the actual local
prime exponents prescribed by \(p\) and its heights \(h_u(p)\).
By the primary-component construction the submodule \(N_u\)
realizing these conditions satisfies, at each marked place \(v\),
\[
              (N_u)_v=N_{u,v}(p).
\]
Insert this joint choice in the preceding transport. Every
coordinate, including the ramified column in the adapted entry,
is the coordinate of \(p\). The transported column thus has the prescribed
localizations and heights \(h_u(p)\). The
twist and tuple laws give its local divisor and character tests.
The scale, mass and size comparisons fixed the threshold before
\(p\) was chosen. \label{cor:automorphic-realization-65}

\end{proof}

\subsection{The period operator at the selected subspace}
\label{subsec:selected-completion-period}

The localizations and heights now come from a single completed point \(q\). We must identify the operator that acts on its selected subspace. We first obtain the
reciprocity data there from rational coordinates and the
rational-pairing law of the chosen norm extension. We then
identify the period of Section~\ref{sec:selected-row-gram}
in the ramified coordinates of that same point. The measured
row, the word cosets and the evaluation translate are those
already chosen on \(q\).


\begin{lemma}[The period pairing in rational coordinates]\label{lem:completed-rational-period-reading}
Let \(D_c(Z)\) be the definite column data at an allowed subspace
of the completion, with its chosen levels, jets, tuple characters
and integral kernel. For a word \(u\) and tuple \(a\) in use, write
\(\mathcal P^{\rm col}_{u,a}\) for the complex period pairing between
its column forms and the row of that tuple. For a rational column
coordinate \(x\) and face coordinate \(F\), the embedded scalar of \(D_c(Z)\) equals
\[
 \mathcal P^{\rm col}_{u,a}
       \big(j_{N,u}(e_{\nu,u}x),\,j_{A,a}(r_{V,0}F)\big),
\]
where \(r_{V,0}\) restricts the original row at the refined level.
\end{lemma}
\begin{proof}
The rational coordinate comparison accepts any Hilbert space,
period operator and distinguished subspace on each \(D_c(Z)\).
Use \((\mathbb C,0,0)\) provisionally. Its word index is the
identity. The fitted rational column isomorphism satisfies
\[
       e_{\nu,u}\big(e_{\nu,u}^{-1}\Lambda_u\big)=\Lambda_u .
\]
For the other argument, the form recovered from \(r_{V,0}F\)
and the restriction of the form defining the face have the same
evaluation. Evaluation on restricted forms is injective, so their
embeddings in the tuple row agree. Applying the definite column
pairing gives the stated scalar, using the twisted form at a
nonreference tuple. \label{cor:automorphic-realization-66}
\end{proof}

\begin{proposition}[Rationality of a nonzero period factor]\label{prop:rational-period-factor}
In the coordinates of Lemma~\ref{lem:completed-rational-period-reading}, the reference sharp
period descends to the rational good quotient. At a used tuple write \(P_\mu\) for the period and \(B_\mu\) for the indicated quotient pairing. For the coefficient embedding \(e:H\hookrightarrow\mathbb C\),
\[
 0\ne P_\mu=c B_\mu\quad\Longrightarrow\quad c\in e(H).
\]
\end{proposition}

\begin{proof}
\proofstep{Keep the provisional periods while forming the rational quotient.}
Put \(H\) for the coefficient field and \(e:H\to\C\) for its
given embedding. For this calculation keep the provisional
periods and the original size coordinates. The new real errors
are those read on the resulting face. The original Petersson
scalar and row scalars still give its primitive, separation
and character identities, since its \(l_{6m,j}\) have their
original values. The tuple labels, sharp and bottom embeddings
and all the coordinate and lattice readings in the displayed
period identity are therefore still available.

On the reference tuple let \(T=j_A r_{V,0}\) be the sharp
coordinate map and \(P_q\) the rational matrix to the good row.
Its range is the entire good coordinate space by the primitive
comparison. On the sharp coordinates the reference row identification makes the two composites agree:
\[
 \iota T=j_{EE}e_RP_q,\qquad T=j_A r_{V,0}.
\]
The complex span of the image of \(T\) is the degree-zero
space \(V_0\) by equality of the two sharp projector ranges.
The subspace removed by the orthogonal projection is
\[
 U=V_0\cap\ker\iota=\operatorname{span}_{\mathbb C}T(\ker P_q).
\] Indeed the images of a rational row basis
under \(j_{EE}\) are complex linearly independent, so ordinary
elimination in the rational matrix \(P_q\) describes its
kernel after extension of scalars. The same computation
applied to \(T\) gives the asserted image in \(V_0\).

Choose a rational \(G\) with \(e_RP_qG=v\). Pick a basis
\(u_a=TG_a\) of \(U\) with \(G_a\in\ker P_q\). The projection
coefficients of \(TG\) satisfy the Gram system
\[
 \left\langle u_a,TG-\sum_b c_b u_b\right\rangle=0
 \quad(\text{each }a).
\]
The Gram matrix is invertible and its coefficients and the
right sides lie in \(e(H)\), a field stable under conjugation.
Hence all \(c_b\) lie in \(e(H)\). Their rational preimages
give \(G'=\sum_b e^{-1}(c_b)G_b\in\ker P_q\). Now
\[
 T(G-G')\in V_0\cap U^\perp,\qquad \iota T(G-G')=j_{EE}v .
\]
This is the representative used by the quotient pairing.
For every rational \(\phi,v\) it follows that
\[
   B_u^{\C}(j_N\phi,j_{EE}v)
    = \mathcal P^{\rm col}_{u,0}
       \big(j_N\phi,T(G-G')\big)\ \in e(H),
\]
by the rational scalar of the definite column with
\(x=e_{\nu,u}^{-1}\phi\). This holds even when the orthogonal
projection of \(TG\) is nonzero.

The division construction requires a family of forms over \(H\)
before rationality has been proved. Put
\(b_u(\phi,v)=B_u^{\C}(j_N\phi,j_{EE}v)\) on rational coordinates.
Use the single test \(b_{u'}(\phi',v')\in e(H)\) for all words and
all rational pairs. Define, simultaneously for every \(u\),
\[
 B_{k,u}(\phi,v)=
 \begin{cases}
  e^{-1}(b_u(\phi,v)),&\text{if the test holds for every pair and word},\\
  0,&\text{otherwise}.
 \end{cases}
\]
In the first case the two coordinate maps are semilinear over \(e\).

For \(a,b\in H\), complex bilinearity gives
\[
 b_u(a\phi,bv)=e(a)e(b)b_u(\phi,v).
\]
Applying \(e^{-1}\) on the image of \(H\) gives scalar bilinearity
there, and additivity follows in each argument in the same way.
The other case gives the zero forms. This defines the full family
before the rationality proof. Localization carries the
resulting conditional identity from the initial state, and
changing \(N,h\) preserves the underlying rational spaces,
maps and this identity. The calculation above satisfies its
hypothesis on the completed data. Thus for all these pairs
the displayed value is \(e(B_{k,u}(\phi,v))\).
\label{cor:automorphic-realization-67}

\proofstep{Choose a nonzero pair of rational coordinates before taking the ratio.}
A rational pairing determines the comparison of two period normalizations from any pair of coordinates on which the denominator is nonzero. The choice below avoids dividing by an individual period that could vanish.

Fix a word \(u\) and tuple \(\mu\). On \(V_0(\mu)\) suppose the period is given by the scalar \(d\) and the projected reading
\[
 v\longmapsto d\,B_u^{\C}\bigl(\iota(p_\perp v)\bigr),
\]
where \(p_\perp\) projects orthogonally onto \(V_0(\mu)\cap U(\mu)^\perp\). Suppose some period there is
nonzero. Write \(p_\perp\) for that orthogonal projection. The kernel calculation identifies the two inputs to the period:
\[
 v-p_\perp v\in U(\mu)\quad\Longrightarrow\quad \iota(v)=\iota(p_\perp v)\qquad(v\in V_0(\mu)).
\]
An unused tuple has zero row image and could not have such a
period. For a used tuple the sharp coordinate vectors project
to a spanning set in \(V_0(\mu)\), and the rational column
vectors span the complex column.

If the right-hand pairing
vanished on every pair of those vectors, complex linearity in
each variable would give vanishing throughout. Thus there
are rational coordinates \(x,F\) with nonzero denominator.
The quotient image of the row vector is \(j_{EE}v\) for a
rational \(v\) by the identification of the row space at the refined level at that tuple.
Consequently
\[
        d=e\!\left(
          \frac{D_c(Z).\operatorname{Pt}_{u,\mu}(x,F)}
               {B_{k,u}(e_{\nu,u}x,v)}\right)
\]
at the selected row subspace \(Z\) of the rational reading.
The original row chooses its scalar to agree with any such
nonzero factor lying in \(e(H)\). This proves that agreement
for every complex factor meeting the stated hypotheses. Here
the factor multiplies the quotient pairing to give the column
period. On a nonzero reading this is the reciprocal convention
to the factor multiplying the column period in
Section~\ref{subsec:selected-period}. \label{cor:automorphic-realization-68}
\end{proof}

\paragraph{The compatible reciprocity data.}

\begin{proposition}[Compatibility of the marked reciprocity data]
Let \(E\) be the original marked CM host. The spectral comparison and standard kernel at the adapted entry admit compatible Galois and automorphic embeddings with
\[
 \iota_{\rm Galois}|_E=\iota_{\rm aut}|_E .
\]
They retain the prescribed ray field, auxiliary extension and local choices.
\end{proposition}
\begin{proof}
The host level in the completed point is supported on the two
type primes, the two distinguished primes and the specified
base auxiliary prime. Its level realization follows from the
word system over this marked host. Retain the existing full and
degree-zero word forms and their local comparison maps, the
normalized character, the divisor geometry and the theta action.
Use the twist supplied by the good member in the local Galois
and norm data. Its choice agrees with the retained selection,
so its character frame and tuple characters, local types,
cover and Galois frame are the prescribed ones.

Thus, writing \(E\) for the CM host, the two comparison embeddings restrict to the same map on \(E\):
\[
 \iota_{\rm Galois}|_E=\iota_{\rm aut}|_E .
\] Artin reciprocity and the Weber
ray-field comparison on the marked imaginary quadratic field
identify the two prescribed field pieces (the ray field and
the chosen auxiliary extension) in this frame. The row
coefficient comparison just proved supplies the scalar
condition on the same rows. These identities assemble the
spectral data with the original geometry, tuple choices and
word columns. \label{cor:automorphic-realization-69}

Use the transported adapted entry to choose the standard
kernel with its original class, frame, characteristic and
support readings. In the standard coordinates the two Hermitian forms have the same
Gram matrix and the two lattices impose the same local membership
test. At each word, stage and admissible frame the identity matrix
therefore satisfies both comparisons:
\[
           I^*H_{\rm std}I=H_{\rm std},\qquad
           I\Lambda_{\rm std}=\Lambda_{\rm std}.
\] The
frame equality uses the ordinary comparison at an ordinary
slot and the conjugate comparison at a dual slot. The
original and twisted lower maps give the two row comparisons
in these coordinates. Use the compatibility of the divisor
character, the additive character of sign \(-1\), and the
common Weil action to retain the joint comparisons.

Let \(S\) contain the assigned sequence primes and all contractions
of nonreference quotient primes. For a base ideal class represented
by \(\mathfrak a\), choose \(b\in\mathfrak a^{-1}\) by approximation
with \(v_P(b)=-v_P(\mathfrak a)\) for \(P\in S\). Then
\[
 b\mathfrak a\subset\mathcal O_E,\qquad
 v_P(b\mathfrak a)=0\ (P\in S),\qquad [b\mathfrak a]=[\mathfrak a].
\]
Over the totally complex field \(E\), Artin reciprocity
identifies the given Hilbert class field with the maximal unramified
abelian extension. Its integers lie in the prescribed field by the
marked integral containment. Transport through the specified
embedding preserves the Galois property of the coefficient extension
over \(\mathbb Q\). These choices give the required realization on
the original host. \label{cor:automorphic-realization-70}
\end{proof}

Let \(Z_*\) be the subspace determined by the actual real row and
its face. Use the positive definite counting Hilbert space
\(\mathcal H_X\) of that measured comparison. Its character
level and its good projector \(G\) were fixed on \(q\).
Transport of the degree-zero word forms identifies them with the
automorphic column functions at this level. Write \(J_u\) for
the resulting bijection from the automorphic column to the
abstract full column. On the former, let \(\mathcal P_u\) be
the Poincar\'e sum using the specified coset representatives for the sequence.
Its membership in \(\mathcal H_X\) and its linearity are the
period identities for the sequence in the measured model.

Write \(e_i\) for the rational coordinate vector and
\(e_{\mu,u}\) for the complex column evaluation.
In the ramified
subspace of the full column put
\[
       b_{u,i}=e_{\mu,u}\,j_{N,u}(e_{\nu,u}e_i).
\]
Identify the word
at a column label by the inverse of the bijection between the chosen sequence indices.
Let \(\sigma\) be identity or complex conjugation according to
the comparison of the coefficient embedding with the selected
archimedean place. Use the Hilbert space with scalar action
twisted by \(\sigma\). On its underlying additive space the
operator at \(Z_*\) is
\[
   \operatorname{out}\,P_{u,*}(z)
     =\mathcal P_u\,J_u^{-1}
        \left(\sum_i\sigma(z_i)b_{u,i}\right).
\]
\begin{lemma}[The operator on the selected row]
With the \(\sigma\)-twisted target action, the underlying operator satisfies
\[
 P_{u,*}(av+bv')=\sigma(a)P_{u,*}(v)+\sigma(b)P_{u,*}(v')\qquad(a,b\in\mathbb C).
\]
Thus it is complex linear into that target. On the good row its restriction is the ramified column reading at \(Z_*\).
\end{lemma}
\begin{proof}
On the underlying Hilbert space, substitution of \(az+bz'\) in
the displayed formula gives \(\sigma(a)\operatorname{out}P_{u,*}(z)
+\sigma(b)\operatorname{out}P_{u,*}(z')\), which is the required
linearity in the twisted space. Membership in the distinguished
row is tested by \(Gv=v\) after inclusion into \(L^2\).
Consequently unfolding \(\mathcal P_u\) on these word functions
restricts that very operator to the ramified columns.
\label{cor:automorphic-realization-71}
\end{proof}

Replace the provisional period record at \(Z_*\) by this one.
Keep the fitted words and their maps \(e_{\nu,u}\). The selected
map into the row space at the refined level \(s_E\) is the measured evaluation at the
original translate and Haar scale times the positive
normalization chosen for the placement. On the sharp subspace
it satisfies \(n_P\|s_Ev\|^2=\|v\|^2\).
Read the size coordinates with the actual \(\Delta,h_0\) and
\(N_c\) from the completion, where \(\eta_{\rm sz}=mh_0\).
For the given real coefficient \(c'\), the update acts simultaneously
on the size coordinates:
\[
 \widetilde l_{6m,j}=l_{6m,j}-c'l_{6u},\qquad \widetilde l_{6u}=0 .
\]
Every subtraction uses the original \(l_{6u}\).

\paragraph{Placement on the updated face.}
The row and divisor comparisons identify the two ambient realizations,
the intermediate modules and the finite group identification
after this replacement. For each nonzero base isotypic row labelled by \(\sigma\), the occurrence condition on the real row has two parts. At each required tuple \(\mu\), and at each other infinite place \(v\), it requires
\[
 V_{\mu,\sigma}\ne0,\qquad \pi_{\sigma,v}\simeq\pi_{\sigma,v}^{\rm prescribed},
\]
where \(V_{\mu,\sigma}\) is the corresponding tuple summand. This is the multiplicity comparison
on the real row, applied with the generic series character and
the chosen type data. With these inputs the column construction
permits the new period and scale at fixed rational readings.

Apply the definite-column comparison to the resulting face.
It uses the equality of kernels and definite columns, the
fitted period map for the sequence and the original rational coordinates. It then
identifies its entire family of period spaces with the
\emph{updated} family, including \(P_{u,*}\) at its selected
subspace. The same comparison gives the actual weights, finite
character labels and compatible column characters.
\label{cor:automorphic-realization-72}

\subsection{The adelic division pair on the completed point}
\label{subsec:completed-division}

The preceding period comparison identifies the operator in the ramified coordinates. For the division argument this identification must also respect the degree-zero forms, their metrics and the refinements of level at every word.

Fix a completed member \(q\) of Definition~\ref{def:completed-member}
and its compatible reciprocity data \(R\). For each word \(u\) the
degree-zero column forms of \(R\) agree with those of \(q\). Apply this equality to the measured
saturation and its remaining spectral data. This gives the
measured saturation for \(R\) with the same group, measure,
evaluation translate and comparison of the full and good rows.
The finite part of the saturation puts each period indexed by the sequence in the
level space of its indicated character. Thus the period
construction above uses the prescribed level projector on
every word.

\begin{proposition}[The completed division pair]\label{prop:completed-division-pair}
On a completed member in the sense of Definition~\ref{def:completed-member}, the real, Hecke, and lower
comparisons identify the square-integrable and geometric
realizations of the definite columns, periods, and size
functions at all subspaces. The identified point carries
the original tuple characters and reciprocity data,
including the prescribed side at each selected prime.
In particular, writing \(q,R\) for the two identified realizations and \(F_u^0\) for the degree-zero forms of word \(u\),
\[
 F_u^0(q)=F_u^0(R)\qquad\text{for every }u.
\]
Both parameter runs are evaluated on this same completed point.
\end{proposition}

\begin{proof}
Use the two ambient identifications on \(q\), the intermediate
modules and the finite group comparison of the preceding
construction. Denote the resulting adelic row data and face by
\(x\). The adelic division pair is formed from the base and
sharp form readings of \(x\) and its word columns. Its
intersection comparison over the coefficient field is the
base change of the one on the original row. The rational word
isomorphisms and the sharp and base ambient identifications
are those in the construction of \(x\).

Let \(n_{\rm sh}\) be the sharp-level index. Write \(E\) for
sharp evaluation with the measured translate and Petersson
comparison, using the conjugate scalar action at a conjugate
reading. The placement chosen in the preceding construction
is \(s_{\rm sh}=\kappa E\) with \(\kappa>0\) and
\[
                 n_{\rm sh}\|s_{\rm sh}v\|^2=\|v\|^2
\]
on the refined row. In particular its phase is \(1\). Writing \(F_u^0\) for the degree-zero form space, the identification on the two realizations is wordwise:
\[
 F_u^0(q)=F_u^0(R)\qquad\text{for each word }u .
\]
The saturation therefore uses the same level projector, character,
evaluation translate and Haar measure in the norm identity.

Let \(e_F,\mathrm{fi}\) be the two coefficient maps in the normal
expansion. A substitution of face sizes or error changes the
arguments of the expansion. Its form and polynomial coefficients
are read in the original rational frame, so the coefficient maps
are still \(e_F,\mathrm{fi}\). The diagonal forms, positive
terms, standard Hermitian frames and their lattices are those
of \(q\); uniqueness gives the standard additive character of
sign \(-1\) on both sides. Use the bounded real comparison on
this core with its frame weights, extremal vectors, level action
and local uniformization maps. All its phase factors are the
one just computed. This identifies the intersection, norm and
period comparisons on a single pair. \label{cor:completed-division-1}

\proofstep{Identify the actual Hecke eigenvalues and tuple characters.}

On the base row write the decomposition by complete Hecke systems as
\[
 V_{\rm base}=\bigoplus_{\lambda}V_{\rm base}[\lambda].
\] Fix one
occurring eigensystem \(\lambda\) and a required tuple \(\mu\).
Multiplicity gives its nonzero constituent at \(\mu\) with
the prescribed infinite components off the distinguished place.
The generic series choice makes the product of the finite
character and level projectors nonzero on the required labels.
At each tuple in use the adelic and marked characters satisfy
\[
 \chi^{\rm adelic}_\mu=\chi^{\rm marked}_\mu=\chi^R_\mu .
\]
The first equality uses the tuple map and the second uses
compatibility with \(R\). In particular the character depends only on the tuple label. For two positions \(a,b\) the implication is
\[
 \mu(a)=\mu(b)\quad\Longrightarrow\quad \chi_{\mu(a)}=\chi_{\mu(b)} .
\] Discrete-series reciprocity at the actual
weights gives the relative parameter as the logarithm of the
positive lift parameter; intersection at the refined level gives
both rational realizations of this constituent. \label{cor:completed-division-2}

\proofstep{Identify the lower maps in both ambient realizations.}

At the base tuple use the reference lower map of the
standard columns. At a nonreference tuple the lower map uses
the integral columns of the actual twist. To see the
polynomial comparison there, let \(P\) be the highest-weight
polynomial in the span for the prescribed top weight. Take
the standard original and twisted columns and a last-stage
frame whose determinant is a unit.
Suppose the evaluation of the specified twisted vector on the
given class agrees at every argument with the function of \(P\)
in that frame, with its prescribed coordinate shift. In the standard last-stage coordinates these evaluations recover the polynomial
\[
 P_{\rm recovered}=\rho_\lambda(A)P,
\]
where \(A\) is the matrix of the last-stage frame and \(\rho_\lambda\) its polynomial action. Applying the lower
map means taking its branching projection. The two true
weights interlace.

In those weights the restriction to the
last block commutes with the projection in the sense
\[
                 \operatorname{br}\big(\rho_{\lambda}(A)P\big)
                  =\operatorname{br}_{\rm tens}(P)
\]
with the lower map expressed in its standard coordinates.
Substitution in the top formula followed by extraction in the standard twisted coordinates recovers \(\rho_\lambda(A)P\). The block branching formula, with the upper and
lower weight shifts, then gives the displayed equality and
the determinant factor of the lower coordinate is \(1\).
This comparison uses the equality of the twist datum, the
standard twisted column maps, and the standard original
columns on \(q\). The unit condition restricts the frame in this statement.

The sharp ambient map is the map obtained from the sharp
identification at \(q\) and the two forms of the tuple
reading. The splitting of the lower modules is the splitting
of the adelic pair. If \(s_1,s_2\) are the stabilizer orders at the two levels and \(C\) the common stabilizer multiple, the original class and level conditions give
\[
 s_1\mid C,\quad s_2\mid C,\qquad \operatorname{lcm}(s_1,s_2)\mid C .
\] These identities supply
both lower maps as well as the two ambient realizations on
\(x\). \label{cor:completed-division-3}

\proofstep{Identify the two sets of marked primes under the prescribed field isomorphism.}

Let \(\theta\) be the auxiliary field isomorphism from the
standard kernel to the column model. Write \(Q_{\ker}(\mu)\)
for its finite set of quotient primes. The two prime comparisons
in the construction give
\[
 \theta_*Q_{\ker}(\mu)=Q(\mu),\qquad
 \theta_*(\operatorname{im}\mathfrak P_{\ker})=
                         \operatorname{im}\mathfrak P .
\]
Here \(\mathfrak P(v)\) is the selected prime of \(E\) over
\(v\) in \(E^+\). A prime contracting under \(\theta\) to
\(P\in Q_{\ker}(\mu)\) is necessarily \(\theta_*P\). The
level comparison at a used tuple consequently puts that very
prime in the adelic quotient-prime set. For
\(\ell\in\{p_d,p_d'\}\), if \(\ell\in\mathfrak P(v)\), then
\(v\mid\ell\). The marked side law at \(v\) identifies the
contraction of \(\mathfrak P(v)\) to the imaginary quadratic
field with the \(\ell\)-prime specified by \(R\).
\label{cor:completed-division-4}

\proofstep{Substitute the common columns before comparing the periods and sizes.}
Compare periods and heights on these identified columns.


For any subspace \(U\), identify its square-integrable column
with its geometric column by the definite comparison. The
resulting period spaces on this one domain agree, including the
Poincar\'e period in the updated component. In symbols, writing
\(C,\mathcal P,H\) for the column, its period space and its size,
the successive identifications give
\[
 C_{L^2}(U)=C_{\rm geom}(U),\qquad
 \mathcal P_{L^2}(U)=\mathcal P_{\rm geom}(U),\qquad
 H_{L^2}(U)=H_{\rm geom}(U).
\]
For the last equality both sides use the record of \(q\), the
nonnegative integer values assigned to \(\Delta,\eta_{\rm sz}/m\) and the same
replacement of \(l_{6m,j},l_{6u}\) at the given real coefficient.
Transport of the bounded real datum through these identifications
preserves the weights, word placement, side and lower maps, on
the point with record \(R\). \label{cor:completed-division-5}
\end{proof}

\begin{construction}
\begin{tabular}{@{}l@{\hspace{1.2em}}l@{}}
\toprule
\textit{Datum of the completed point} & \textit{Identification used in both runs} \\
\midrule
Word and tuple character & Same rational column and level projector \\
Measured sharp evaluation \(E\) & \(s_{\rm sh}=\kappa E,\ \kappa>0\) \\
Sharp-level index \(n_{\rm sh}\) & \(n_{\rm sh}\|s_{\rm sh}v\|^2=\|v\|^2\) \\
Base and sharp forms & One adelic division pair at the chosen translate \\
\bottomrule
\end{tabular}
\caption{The comparison data carried into both parameter runs by Proposition~\ref{prop:completed-division-pair}.}\label{fig:analytic-completed-pair}
\end{construction}

\subsection{Two parameter runs on one division pair}
\label{subsec:completed-prime-runs}

Choose the point of Proposition~\ref{prop:completed-division-pair}.
The depth estimates and the exceptional-prime estimates call for two choices of numerical parameters. Their comparison uses the same row, level and period identification at this point.

\paragraph{Row normalization and the numerical point.}
The numerical point used for the runs has the face, word
lattices, period spaces and scale of the bundle. To form it,
take a recursive realization with the base and refined level
characters, a normalization of the real rows, and the local
row comparison between that normalization and the recursive
realization. This is the choice of comparison denoted by a
summand of the bundle. Denote the row spaces before and after this comparison by superscripts \(1,2\). The same identifications apply simultaneously to both rows and the word columns:
\[
 V_{\rm full}^{(1)}=V_{\rm full}^{(2)},\qquad V_{\rm good}^{(1)}=V_{\rm good}^{(2)},\qquad M_u^{(1)}=M_u^{(2)}\quad(\text{each word }u).
\]

The row normalization uses the equality of the character
frames, the identification of the type primes with the
kernel's type places and the conjugate-level and character
identities of the finite translate. It keeps the standard
frames and the coefficient values on the standard word
vectors. For the recursive comparison use the standard
additive character on these levels and the local row
reading of the marked Galois level. The degree-three
normalization uses the equality of the zero section with
the one in the original scale datum.

For \(z=v\) or \(\bar v\), write \(I_j(z),J_j(z),Z_j(z)\) for the
base-row index, lower-map index and value of the zero-section test at lift \(j\);
nonzero means that the corresponding test detects the place. The twist
support has the following intrinsic description:
\[
 \mathcal S_{\rm tw}
 =\{v:\ (I_j(z),J_j(z),Z_j(z))\ne(0,0,0)
            \text{ for some }j\text{ and }z\in\{v,\bar v\}\}.
\]
All three tests precede the choice of the height and error constants.
Let \(S\) be the set selected by the three tests and \(\mu_0\) the base tuple. The tests identify the tuple map on conjugate places by
\[
 S=\bar S,\qquad v\in S\ \Longleftrightarrow\ \mu(v)\ne\mu_0,\qquad \mu(\bar v)=\mu(v).
\] In particular each
single-index place is twisted.
These are comparisons on \(x\). Substitution of the three
family equalities of Section~\ref{subsec:completed-division}
transports them to the bundle with \(R\) fixed.
\label{cor:completed-division-6}

\paragraph{The row at the chosen comparison.}
\label{par:linked-family-rows}
Fix a completed member \(q\) and its division pair from
Proposition~\ref{prop:completed-division-pair}. Choose a recursive
realization, a normalization of its real rows and the identification
between these realizations on the pair. Write \(P\) for the resulting point. The row,
periods, finite groups, coefficient field and numerical functions
at \(P\) are those of the pair. In particular its two level rows
are the ones in Section~\ref{subsec:completed-division}. The
following verifications hold for each compatible comparison on
the bundle. Thus they also hold if an additional condition, such
as the one on the evaluation grid, selects the comparison.

\begin{proposition}[The two rows at a compatible choice of comparison]\label{prop:linked-comparison-rows}
At every compatible comparison on the completed bundle, let \(q_0\) be projection to the sharp good row and \(P_{u,0}\) the reference pairing of word \(u\). The factorization is characterized by
\[
 F-F'\in\ker q_0\quad\Longrightarrow\quad P_{u,0}(\phi,F)=P_{u,0}(\phi,F')\quad\text{for all }\phi.
\] At either level the rational row and its integral
submodule have the corresponding divisor comparisons of
the original division pair, on all tuples used by the real
row. Any later choice of a comparison satisfying an extra
condition therefore still gives these rows and their pairings.
\end{proposition}

\begin{proof}
\proofstep{Use the sharp reference period to test the kernel of the good projection.}
Fix a word \(u\), and use its selected Poincar\'e operator from
Section~\ref{subsec:selected-completion-period}. Let \(R_P\) be the
rational row, \(R_P^\sharp\) its common sharp good row and
\(q_0:R_P\to R_P^\sharp\) the good projection. In the complex
realization write \(I_u\) for the column isometry and \(I_R\) for
the isometry on the good row at the refined level. Unfolding by
Lemma~\ref{lem:joint-period} gives, for the sharp tests of this
operator,
\[
 \iota\bigl(P_{u,0}(\phi,F)\bigr)
  =\mathcal P_{u,0}\bigl(I_u\phi, I_R(q_0 F)\bigr),
\]
where \(\mathcal P_{u,0}\) denotes the resulting complex theta
pairing.

The normalized factor is one at the reference tuple
of Section~\ref{sec:transport}. The reference tuple is obtained
from the full diagonal and the dimension of the ramified space
there, before applying either division estimate. If \(F-F'\in\ker q_0\),
the difference of the right sides vanishes for every \(\phi\).
The coefficient embedding is injective, so the vanishing difference implies
\[
 F-F'\in\ker q_0\quad\Longrightarrow\quad P_{u,0}(\phi,F)=P_{u,0}(\phi,F')\qquad(\text{every }\phi).
\] This defines the rational
reference pairing on the quotient without choosing an integral
representative.

\proofstep{Read the quotient pairing on the transported integral lattices.}

The rational pairing obtained this way is the one characterized by
the Gram equations in Proposition~\ref{prop:rational-period-factor}.
For each of the two levels, the rational row comparison identifies
the down-lift and the theta map in the same coordinates. To see what
is integral, let \(M_w\) be the localization of the glued lattice
at a place \(w\).

For the coordinate identification \(a\) coming
from the completed pair, its lattice and form are transported by
\[
 M_w\longmapsto aM_w,\qquad
 B_w(\phi,F)\longmapsto B_w(a^{-1}\phi,a^{-1}F).
\]
Here \(a\) denotes the respective coordinate map in each argument.
Thus the localized orders, lattice indices and dual indices are
computed on the identified modules. The reference form descends
because of the kernel calculation above. A nonreference tuple of
the real row occurs as the tuple of a place obtained by the inverse
place bijection of the host. Hence for each used tuple \(\mu\) one of two comparisons is available:
\[
 \mu=\mu_0\quad\text{or}\quad \mu=\mu(w)\ \text{for a place }w\text{ on the nonreference support}.
\]
The measured intersection comparison supplies the surjections
onto the good coordinates and the separation of the primitive types.
The equality of the bounded real realizations identifies the real
datum and its face in both calculations.
\label{cor:linked-family-rows-1}

On a nonreference summand, let \(c\) be an integral level coordinate
reading to an admitted source value and \(\pi_{\rm src}(c)\) its
projection to the lift source. Write \(x=\Theta\pi_{\rm src}(c)\)
in form coordinates and choose \(x=\sum_h a_h f_h\), with \(f_h\)
rational forms in the expansion model and \(a_h\) in the row field.
If \({\rm ev}_w\) denotes the local expansion map, the divisor theta
identity and the signed series comparison give the same value
\[
        {\rm ev}_w(x)=\sum_h\iota_w(a_h)\,{\rm ev}_w(f_h).
\]

To check independence of the expansion, subtract the two expressions for \(x\) in the extended space and use linearity:
\[
 \sum_h b_h f_h=0\quad\Longrightarrow\quad \sum_h b_h\,{\rm ev}_w(f_h)=0 .
\] The completed
division comparison transports \(\pi_{\rm src}\), the rational form
coordinates and \({\rm ev}_w\) through the update of periods and
sizes. Consequently the displayed equality can be checked either
before or after the update. In the chosen rational basis, evaluation on the row is the composite 
\[
 x\longmapsto(a_h)_h\longmapsto(\iota_w(a_h))_h\longmapsto\sum_h \iota_w(a_h)\,{\rm ev}_w(f_h).
\]
 Its restrictions to the integral expansion lattice and theta sublattice give the integral evaluation and the assigned series, respectively.
\label{cor:linked-family-rows-2}

Forget the grid in the arithmetic member of
Section~\ref{subsec:automorphic-realization}. The norm choice,
the signed series used for the exponent, the two type maps, the
base and local level groups, the kernel field and the selected
half primes remain those of this member. The changes in the
Petersson scalar, in the lift scalars and in the error have no
effect on the maps of finite ideles or integral host rings.

Fix a preliminary frame and its finite support data. If \(S\)
is an allowed support, put
\[
    M(S)=g_0g_1\!\!\prod_{\substack{q\text{ a residue cardinality}\\
                         q\text{ of a place in }S}}\!q ,
\]
with each distinct cardinality used once, as in the level of the
member. By the support bound the possible \(M(S)\) belong to a
finite list. If \(A_M\) generates the embedded ray field at level \(M\), one finite request suffices:
\[
 A=\bigcup_{M\text{ in the level list}} A_M,\qquad |A|\le\sum_M |A_M|<\infty .
\] The Artin comparison of
Section~\ref{subsec:automorphic-realization} identifies these ray
fields and the type and character values in the chosen embedding.
The two canonical frames on the same core give the same half-level
ideals as the divisor, including the half-error divisibility.

The integral field reading identifies the coefficient rings. After
passing to fractions it identifies the degree and absolute
discriminant used in the size bound with those of the recorded core.
The single-index support test uses the case with zero error and
height constants. These comparisons give the field, rank, series
and level requirements in the two runs. A coordinate isomorphism preserves each of the module equalities in the comparison. For submodules \(U,V\) in an identified row and its isomorphism \(b\),
\[
 U=V\quad\Longleftrightarrow\quad b(U)=b(V).
\]
\label{cor:linked-family-rows-3}

At infinity the intertwiner is defined on the prescribed lowest
compact type of the polynomial Weil representation. Regularity,
admissibility and finite multiplicity give the Harish--Chandra
modules. Apply Casselman--Wallach globalization to this intertwiner
and cuspidal rapid decay to its integrals, as in
Section~\ref{sec:transport}. Let \(T\) be the difference of two such intertwiners,
and let \(V_0\) be the span of the compact translates of the
distinguished lowest vectors. Equality on the test vectors gives
\(T(V_0)=0\).

At the Harish--Chandra module level,
\[
 V=U(\mathfrak g_{\mathbb C})V_0,\qquad
 T(Xv)=X T(v)=0
 \quad\bigl(X\in U(\mathfrak g_{\mathbb C}),\ v\in V_0\bigr).
\]
Thus \(T=0\) on the module and, by continuity, on its globalization. This explains the injective reading used for the geometric
action. The conjugate frame is read by the transported coefficient
map. When \(n_V\ge8fr+4\), the stable theta span and the form
isomorphism of Section~\ref{subsec:selected-cusp-reading}
apply to these weights. Since the normalized map and the extended
intertwiner agree on every vector of the common orthonormal basis,
they agree on the row. In particular evaluation and the two
linear readings assign the same form at either frame.

At a supported finite type, let \(e_{\rm good}\) be the central
idempotent of Section~\ref{subsec:automorphic-realization}. For a
vector \(f\) in its image and a finite type-group element \(g\),
\[
       e_{\rm good}(g f)=g(e_{\rm good}f)=g f .
\]
The chosen nonzero component therefore generates a module for
the supported split factors in the same constituent. Apply
Green's regular cuspidal character and the depth-zero type
correspondence to the factor actions. Normalized induction and
uniqueness of cuspidal support give the multiset of the local
parameter. The marked local Galois identification reads it at
the coefficient choice of \(P\). This applies to each required
label with the threshold on the family of members. The primitive
base lifting isomorphism is rational and bijective in its
phase-one coordinates and has the uniform positive rank \(D_0\).
The tuple map, level subgroups and base character therefore
refer to this same row on the two conjugate sides.
\label{cor:linked-family-rows-4}

\end{proof}



The two runs below use the rows of Proposition~\ref{prop:linked-comparison-rows}
for the selected compatible comparison.

\begin{lemma}[The two prime orientations on the common core]
For \(\ell=p_d,p_d'\) there is a full short-section extension of
the same core \(J\) with distinguished prime \(\ell\) at every
short weight. For a twisted place \(v\), the determinant character and short cover use the marked half prime of the assigned orientation, where
\[
 p(v)=\begin{cases}p_d',&v\mid p_d,\\ p_d,&v\nmid p_d.\end{cases}
\]
\end{lemma}
\begin{proof}
Write \(J\) for the common geometric core at \(P\). It
has its canonical frames, positive ranks and half-level
divisibility of the error data. The two field conditions and
the values of both distinguished characters hold in the
prescribed fields. The half-level ideal \(\mathfrak m_b\)
of the common core is the product of the marked type prime
ideals of the row. In fact the full level is squarefree
and factors into its two indicated halves. The divisor
comparisons and the equality of the core and row half
primes give the same prime factors on the indicated
side. Equality follows from the squarefree factorization.

Fix \(\ell\in\{p_d,p_d'\}\). For this orientation choose
the distinguished character \(\delta_\ell\) of the frame at
\(\ell\), and use the field condition at \(\ell\). Define
the own prime of a coefficient place \(v\) by
\[
  p(v)=
  \begin{cases}
    p_d' & \text{if the residue characteristic of \(v\) is \(p_d\)},\\
    p_d & \text{otherwise}.
  \end{cases}
\]
For twisted places with \(p(v)=\ell\), the marked side
comparison supplies one prime of the imaginary quadratic
field over \(\ell\) for all the quotient primes over
\(\ell\) on the tuples of those places. Extension of that
prime is contained in every one of these quotient primes,
by contraction. If the class of such places is empty, any
prime above \(\ell\) suffices for this orientation.

On the finite ideles the character construction gives
characters \(\xi_\ell,d_\ell\) with
\[
                        \xi_\ell(a/\bar a)=d_\ell(a).
\]
Here \(d_\ell\) is the prescribed finite-idelic restriction of
\(\delta_\ell\). A character of finite order on the ideles
is trivial on the archimedean subgroup of the CM field,
since its archimedean part is a product of copies of \(\C^\times\).
For a character \(\xi\) of order \(n\) and any archimedean idele
\(z\), take an archimedean \(w\) with \(w^n=z\). Then
\[
 \xi(z)=\xi(w^n)=\xi(w)^n=1 .
\] Thus
its restriction to the finite ideles inherits triviality on principal ideles,
triviality on real ideles, and the specified coefficient
values.

The finite character is trivial on integral units congruent to \(1\)
at the prescribed level. At a prime over \(\ell\), choose
a unit with a nontrivial character value, using the
distinguished character's prescribed tame ramification.
An idele equal to this unit at that prime and to \(1\)
elsewhere is congruent to \(1\) at the two type primes
because their rational primes differ from \(\ell\).
This gives the required nontrivial test there.

On finite ideles put \(h(b)=b/\bar b\). Finite idele Hilbert
90 gives \(\operatorname{im}h=\ker N_{E/E^+}\). Since
\(\ker h\) consists of real ideles and \(d_\ell\) is trivial
there, the assignment
\[
 \xi_\ell(h(b))=d_\ell(b)
\]
is well defined. Extend this character to the finite idele
group; its required congruence and finite-order properties
concern the norm-one subgroup. For a principal norm-one
element \(a\), explicit preimages are
\[
 b=1+a\ (a\ne-1),\qquad
 b\in E^\times,\ \bar b=-b\ (a=-1).
\]
Thus \(\xi_\ell(a)=1\) on principal norm-one elements.
The congruence form of Hilbert 90 transports the tame
congruences at the two type primes and at \(\ell\), giving
the stipulated ramification and coefficient values.
\label{cor:completed-division-7}

Extend the core to the upper geometric model, then take the
full short-section extension with the prime and the
character just chosen. These constructions use the laws
of \(J\), the canonical frames, positive ranks and the
host hypotheses at \(\ell\). The full extension has the
same core \(J\), and its short cover at each lift has
distinguished prime \(\ell\) at every short weight.
For a determinant root \(b\) of an adelic element \(g\), let \(b_f\)
be its finite part. The finite determinants and character values satisfy
\[
 \det(g)_f=b_f/\bar b_f,\qquad
 \xi_\ell(\det(g)_f)=d_\ell(b_f).
\] Rational determinant roots give
value \(1\). At a twisted \(v\) with \(p(v)=\ell\), the
same formula gives the specified determinant character and
the chosen short cover above the marked half prime. This
verifies the own-prime and character pinning at that place.
Taking \(\ell=p_d\) and \(\ell=p'_d\) gives the two extensions.
\label{cor:completed-division-8}
\end{proof}

\begin{established}[The split refinement of the marked core]\label{hyp:split-section-refinement}
In the local construction write \(p=\ell\), \(M=g_0g_1\) for
the original full level and \(I=P\mathcal O_E\), where \(P\)
is the chosen prime above \(p\) in the marked imaginary
quadratic field. If \(J\) denotes the old half, then the new level data are
\[
 (\text{level},\text{chosen half},\text{auxiliary level},\text{field})
                       =(Mp,JI,1,F),
\]
with the same coefficient field \(F\). For each \(\sigma\in\operatorname{Aut}(\mathbb C/F)\), the marked CM curve and its pullback are isomorphic by an isomorphism respecting every torsion label at level \(Mp\) and at each divisor of \(Mp\).
\end{established}
\begin{proof}
We have \(I+\overline I=\mathcal O_E\) and
\(I\overline I=(p)\), by splitting. Use at each lift the
polarization and moduli problem of the core \(J\) just constructed.
At each \(j\) refine the split level by \(I,\overline I\).
Writing \(J\) for the old chosen half, the updated parameters are
\[
 (\text{integer},\text{chosen half},\text{auxiliary level})=(Mp,JI,1).
\] The two complementary halves exist by the displayed
factorization and coprimality. Use the moduli space with this
refined marking and the map forgetting its new level.
The coefficient field stays fixed. The field contains the specified roots of the coprime orders \(M\) and \(p\); their product has order
\[
 \operatorname{ord}(\zeta_M\zeta_p)=\operatorname{lcm}(M,p)=Mp .
\]

Write \(F_{\rm rec}\) for the embedded record field and \(F\) for
the cover coefficient field. The pin on the coefficient field
sends the generators of \(F_{\rm rec}\) to their prescribed complex
values in \(F\), so \(F_{\rm rec}\subseteq F\).
For the marked CM lattice \(\Lambda\) the field conditions read
\[
 k\big(j(\Lambda), W(P):P\in (Mp)^{-1}\Lambda/\Lambda\big)\subseteq F .
\]
Here \(k\) denotes the marked imaginary quadratic field and \(W\) its Weber function. 

For each automorphism
\(\sigma\) of \(\C\) fixing \(F\), these values are fixed. By the main theorem of complex multiplication in its ray-class form \cite{cm-reciprocity}, the Weber values at level \(Mp\) lie in its ray field, contained in \(F\) by Lemma~\ref{lem:marked-original-geometric-field}. The equality of \(j\)-invariants gives an isomorphism to the pullback. Since \(F\) contains the CM field, it respects the specified CM action. Let \(P_0\) generate the free rank-one \(\mathcal O_k/(Mp)\)-module of torsion. Write \(f\) for the isomorphism to the pullback and \(W\) for its Weber quotient. Since \(W\) is the quotient by the CM units, the equality of the two readings gives
\[
 W(f(P_0))=W(P_0^\sigma)
 \quad\Longrightarrow\quad f(P_0)=uP_0^\sigma
                  \quad(u\in\mathcal O_k^\times).
\]

The map \(u^{-1}f\) now matches \(aP_0\) for every \(a\in\mathcal O_k/(Mp)\) by CM-linearity. This isomorphism may depend on \(\sigma\). For any \(m\mid Mp\) the
same isomorphism works at level \(m\), since the point with
label \(a/m\) has label \(((Mp/m)a)/(Mp)\) at the original
level.
\label{cor:split-level-sections-1}
\end{proof}

\begin{proposition}[A fixed split refinement for the short theta rules]\label{prop:split-level-sections}
For the split prime and the marked core \(J\) of Lemma~\ref{hyp:split-section-refinement},
one can choose the same cover, coefficient field and weight-zero
determinant character for every number \(n\) of torsion arguments.
For \(n\ge r_j=p_j+q_j\) the Fourier coefficient is obtained from
the basis test below, with the inverse character under a change of
marking. If \(h_n\) denotes the coefficient attached to \(n\) torsion arguments, the rank obstruction gives
\[
 n<r_j\quad\Longrightarrow\quad h_n=0.
\]
\end{proposition}

\begin{proof}
For the component comparison the refined auxiliary level is
one, the rank is positive, and the marked lattice supplies a
rank-one definite datum for the same host. In indefinite
signature the rank is at least three. Under these hypotheses
the CM comparison for the unitary period spaces says that the
pullback isomorphism fixing the torsion labels and the host
root induces the identity permutation of the period
components. Consequently an algebraic constant in the
function field of a component is fixed by every
\(\sigma\in\operatorname{Aut}(\C/F)\) under its generic
embedding. Indeed, if an algebraic constant \(a\) had degree greater than one, separability and extension of embeddings to \(\mathbb C\) would give
\[
 \sigma\in\operatorname{Aut}(\mathbb C/F),\qquad \sigma(a)\ne a,
\]
contrary to the asserted invariance. Thus \(F\) is relatively algebraically closed in
each component function field. The component idempotents and
the prescribed values in \(F\) now give the componentwise
weight-zero algebraic rules.
\label{cor:split-level-sections-2}

\proofstep{The determinant character.}
Here one must put a nontrivial character on that refinement.
Let \(\mathbb A_{E,f}^1\) be the norm-one subgroup of the
finite ideles. Hilbert 90 gives
\[
 1\longrightarrow\mathbb A_{F,f}^\times
 \longrightarrow\mathbb A_{E,f}^\times
 \xrightarrow{\,h:y\mapsto y/\bar y\,}\mathbb A_{E,f}^1
 \longrightarrow1 .
\]
The finite character \(d_p\) read from \(\delta_p\) is trivial on
the left-hand subgroup. In taking the finite part its
archimedean kernel is also trivial, by divisibility of the
connected complex idele factors and finite order of the
character. Hence \(\xi_p(h(y))=d_p(y)\) is independent of the
lift \(y\). A field-valued complex character on this
subgroup extends as a \(\C^\times\)-valued character of the
abelian group, by divisibility of \(\C^\times\). Only its values
on norm-one determinants will be used. A rational norm-one
determinant has a rational lift by Hilbert 90, so \(\xi_p\)
is trivial there.

At principal level the local congruence
version provides an integral lift congruent to one at \(p\).
At the odd split special prime the componentwise involution interchanges factors, so the lift \(z=(u,1)\) gives
\[
 z/\bar z=(u,1)/(1,u)=(u,u^{-1}).
\] The tame
character therefore gives one on the required level kernel,
and its nontrivial residue value at \(p\) survives on the
determinant quotient. Its values on norm-one determinants are
the original finite-order values in the coefficient field.
The function \(g\mapsto\xi_p(\det g)\) is constant on the
corresponding marked components, by the rational and level
invariance. The preceding component comparison gives a
weight-zero rule with these values, whose same positive power
is one on every component.
\label{cor:split-level-sections-3}

\proofstep{The finite Fourier coefficient.}
The finite Fourier expansion must retain the character just obtained on the determinant.

Write \(r_j=p_j+q_j\) and consider \(n\) torsion arguments on
this cover. For \(n\ge r_j\) fix a basis \(y\) of its second
split half module of rank \(r_j\). On an \(n\)-tuple \(x\)
whose first \(r_j\) entries form a basis \(g y\), set
\[
              h_n(x)=\psi(g^\dagger)^{-1}.
\]
Here \(\psi\) is the character on automorphisms of the first
half induced by the nontrivial determinant character, and
\(g^\dagger\) is the adjoint between the two dual halves.
Set \(h_n=0\) elsewhere. Over the finite residue quotient these
modules are free and a spanning \(r_j\)-tuple is a basis,
as can be checked at every local factor by its determinant.
The basis change \(g\) is unique. Thus \(h_n(y,0)=1\). The inverse character under change of marking follows from reversal of composition by adjunction:
\[
 (g'g)^\dagger=g^\dagger(g')^\dagger .
\]

Take the finite Fourier
sum of \(h_n\) against the additive character of the level
pairing in the \(n\) arguments. Its coefficients lie in the
field just specified. This is the finite coefficient of the
short theta rule at this order. When \(n<r_j\) its spanning
condition has empty hypothesis and we take the zero coefficient.
The weight-zero twist section, the underlying cover and its
coefficient field were fixed before \(n\). Only the coefficient
of the theta rule changes. Finite Fourier evaluation and the
harmonic chart transformation compute these short rules on the
same cover at all orders, with the same determinant twist.
\label{cor:split-level-sections-4}
\end{proof}


\begin{lemma}[One level action for both orientations]
Let \(\bar A_g^{(1)},\bar A_g^{(2)}\) be the two induced operators of the finite row group on the core half-level quotient \(\Lambda/\mathfrak m_b\Lambda\) at lift \(j\). With the common level-form identification,
\[
 \bar A_g^{(1)}=\bar A_g^{(2)}\quad\text{on }\Lambda/\mathfrak m_b\Lambda\qquad\text{for every }g.
\]
This action intertwines that on the level forms.
\end{lemma}
\begin{proof}
Write \(\Lambda\) for this lattice and \(A_g\) for the core level
map at \(g\). The congruences of the level comparison give
\[
 \bar A_1=1,\qquad \bar A_g\bar A_h=\bar A_{gh}
       \quad\text{on }\Lambda/\mathfrak m_b\Lambda .
\]
In particular \(\bar A_{g^{-1}}\) is the two-sided inverse
of \(\bar A_g\). This action is defined on \(J\) before either
extension. Two integral lifts of \(\bar A_g\) differ by a map
into \(\mathfrak m_b\Lambda\); on the complementary level
sublattice their difference lands in the full level multiple.
Thus they give the same allowed change of level vector. Pullback
on forms depends only on the quotient map, and the geometric
action in the common frame intertwines it with the row action.
\label{cor:completed-division-9}
\end{proof}

At \(v\) use the extension at \(p(v)\). Its own prime
differs from the residue characteristic, including when
the residue characteristic is \(p_d'\). The place
\(\bar v\) uses the same orientation as \(v\).
Consequently the short-cover laws, the
determinant comparisons and the one level action apply
on the indicated parts of this single numerical point.
No scale cutoff is used in this selection. The linked row
hypotheses and the numerical inequalities for section counts
have their own thresholds.
\label{cor:completed-division-10}


\begin{lemma}[Common form coordinates on the two runs]
At each lift \(j\), write \(I_j\) for the common arithmetic frame identification. The same inverse coordinate map can be used for both orientations:
\[
 z_j^{(1)}=z_j^{(2)}=I_j^{-1}=:z_j.
\] It reads the
same finite action, good projector and depth scalar at \((v,j)\).
The short sections at \(v\) come from its own prime \(p(v)\).
\end{lemma}
\begin{proof}
Use the phase-one base row of \(P\). Its weights are the positive
entries \(a^+_{j,\tau}\) and the reversed negatives of
\(a^-_{j,\tau}\), of signature \((k+j-1,m-j)\) and central
weight \(n_V\). These are the product order and the light and
short weights of the numerical record.

Let \(I_j\) denote the injective comparison of core cusp forms
with the arithmetic frame. The level quotient determines the
pullback \(A_g\) for each element of the mirabolic group.
It suffices to read the representation \(\rho\) at one coefficient
place, because \(I_j A_g=\rho(g)I_j\) in the common frame.
In particular
\[
 I_jA_gA_h=\rho(g)\rho(h)I_j=I_jA_{gh}.
\]
Injectivity gives composition, and gives linearity by the same
argument. The basis comparison makes the scalar extension of
\(I_j\) invertible.

Take \(z_j=I_j^{-1}\) on this extension, as
in Paragraph~\ref{par:selected-form-coordinates}. For the
arithmetic projector \(e_{\rm good}\) and a rational class-sum
map \(T_\mu\) in these coordinates, the corresponding form
readings are \(z_je_{\rm good}I_j\) and \(z_jT_\mu\).
Thus
\[
 I_j(z_je_{\rm good}I_j)=e_{\rm good}I_j,\qquad
 I_j(z_jT_\mu)=T_\mu .
\]
Use this one isomorphism for the ambient lattice and the scaled
and unscaled short sections at \(p(v)\). The section lattice in
the good row is the projector image of the theta lattice at
\(\mu(v)\).

The embedding of the host integers into the point field factors
through the integral coefficient ring of the core, and agrees
with complex conjugation on the host involution. It therefore
reads the same half-level and error generators and the same
scalar in the form and depth comparisons. Its integral exponent
is the depth in the local section estimate. The remaining
parameters \(f,m,r,k,x,\epsilon,b',X''\), the type primes and
the determinant and cutoff data are those of the common numerical
point. Only the local short-section spans depend on the
distinguished prime. \label{cor:completed-division-11}
\end{proof}


\begin{proposition}[One remainder before choosing the cut]
At the common point of Proposition~\ref{prop:completed-division-pair},
assume the linked rows of Proposition~\ref{prop:linked-comparison-rows}
and the section counts of Proposition~\ref{prop:slice-rank-shares}.
Use the simultaneous short basis of
Lemma~\ref{lem:slice-short-basis-before-localization} on the cut by
the selected light columns, with \(C_{n_W}=10\), \(\log g_i\le9L+1\), \(C_P\ge1000\)
and \(C_P^{20}\le frL\). At sufficiently large scale the section
remainder on each lift of either orientation satisfies
\[
                 \langle r\rangle\le R_j fmb\delta_2 .
\]
The scale threshold is common to the orientations and lifts.
At each place the selected remainder precedes the choice of the
cut and its integer depth.
\end{proposition}
\begin{proof}
Use the core coefficient bound \(C_K\) in \(S_\kappa\le C_K kfb\)
in the scale inequalities \eqref{eq:slice-remainder-scales}.
Write \(X_{\rm net}(20,C_K)\) for a scale beyond which the size
inequalities of Proposition~\ref{prop:slice-evaluation-net} hold
with exponent \(20\) and this core bound. Increase the common
scale threshold \(x_1\) to satisfy
\[
 x_1\ge\max\{10^{60},X_{\rm net}(20,C_K)\}.
\]
Read the half-error divisibility on \(J\), and the nonzero short
sections on the selected full extension. The full, short and
light weights and the index depth are those of
Section~\ref{subsec:paired-row-indices}.

At \(p=p(v)\) fix the split marking, once for all three weight spaces. The input level and its special character are
\[
 g_0g_1p,\qquad g\longmapsto\xi_p(\det g),\qquad \xi_p(y/\bar y)=d_p(y).
\] Read them in the core weight spaces at
\(\kappa_j,w_{F,j},w_{L,j}\) by the common form isomorphism.
The incidence chart is the one for the light columns of this
run on the core. Its full level in the light-weight inequalities
is \(g_0g_1\). Select the \(d_s\) scaled short sections before localization by
Lemma~\ref{lem:slice-short-basis-before-localization}, using independent
evaluation on the geometric generic intersection. The shifted
discriminant of this basis is computed in \eqref{eq:slice-discriminant-order}.

If \(A\) is the selected square evaluation minor over the generic
field \(K\), then for each branch field \(K_\zeta/K\) one has
\[
                    0\ne\det A\in K\hookrightarrow K_\zeta .
\]
Thus the same sections give full rank on every branch. Their local
integral index still contributes to \(I_\lambda\). The choice by rank in the
boundary and complementary spaces preserves the light columns
and, before the final lift, the same auxiliary section
\eqref{eq:selected-eta-product}. Hence the finite valuation
test, the two sections for the height calculation and the
determinants concern this one family with the new choice of sections. The evaluation
net gives the second norm on the same rational spaces and
integral section lattices. Changing this norm in the
determinant proof preserves the index on those modules.


Let \(e\) be the branch ramification index. At its reference point the weight-\(w_F\) section \(g_0\) fixes the valuation offset
\[
 c_\zeta=v(g_0(y_\zeta)).
\]
It cancels against the denominator in the quotient reading.
The Gauss value \(\nu_\zeta\) of the generic light row is the
minimum over its light columns, giving the shifted order
\[
              e\bigl(v_{y_\zeta}(F)-T\nu_\zeta\bigr).
\]
This uses the normalized point valuation of the depth estimate.
The vertical, horizontal and archimedean comparisons in
\eqref{eq:slice-height-disc} bound the discriminant and auxiliary
costs on this slice. The remaining two costs are its boundary
determinant and complementary CM determinant heights. Their
ranks give the \(s^2/T\) savings in \eqref{eq:geometric-four-terms}.
At the definite lift the auxiliary cost and complementary rank
are zero. The four bounds give the claimed inequality with the
chosen finite net and the same polynomial bounds at \(frL\)
for every \(j\) and both \(p\).
\label{cor:completed-division-12}
\end{proof}

Extending the form frame to the point field now puts this remainder
estimate in the coordinates of the paired index bound. The
restriction and depth argument uses that bound on the same
lattice and level action. \label{cor:completed-division-13}

\subsection{The joint estimate on the completed family}
\label{subsec:joint-family-estimate}

The two parameter runs now give local inequalities on one realized division pair. Returning to the original period coordinates and summing over places will put them in the same determinant estimate as the cohomological numerator.

At each scale retain the field and the height exponent of
Lemma~\ref{lem:window-parameter-choice}. The original-host estimate in
Proposition~\ref{prop:actual-rank-budget} applies to all the required
twists. Lemma~\ref{lem:strict-budget-after-cut} restricts its sum to
the good types, and Proposition~\ref{prop:weighted-dual-cost} bounds
the localized numerator on the linked row. We now compare its period
lattices on the two covers of that same row.
\begin{definition}[Indices on a pair of conjugate places]\label{def:joint-pair-indices}
For a fixed completed member of Definition~\ref{def:completed-member}, write \(\omega_w=\log\mathrm Nw/[K:\mathbb Q]\) for the coefficient
place weights, \(e_w\) for absolute ramification and
\(\mathcal E=x^\epsilon(B_0+N_c)\).
At a twisted place \(w\mid p\), let \(I_{wj}\) be the sum of the
two signed transport indices of lift \(j\) at \(w,\bar w\). The
error before passing to the restricted lattice \(r^0_{j,w}\), forced term \(F_j(w)\) and integer
least exponent \(J_{j,w}\) are as in Section~\ref{sec:word-division}.
\end{definition}

\begin{lemma}\label{lem:joint-family}
Fix the family with the parameter choice of
Lemma~\ref{lem:window-parameter-choice} and the linked base
and refined rows of Proposition~\ref{prop:linked-comparison-rows}.
Use the good types of Lemma~\ref{lem:strict-budget-after-cut} on the
original host and the integral numerator of
Proposition~\ref{prop:weighted-dual-cost} on these rows. There are \(C_0\ge0\) and \(X_*\) such that, at
every scale \(x\ge X_*\) in the family's domain, the base row of
each compatible completed member of Definition~\ref{def:completed-member}
satisfies, for \(C\ge C_0\),
\begin{align}
 \sum_j I_{wj}
 &\le D_0\left(\sum_j(2e_wF_j(w)+e_w r^0_{j,w})
                         +2e_w/p^2\right),\label{eq:joint-parent-local}\\
 r^0_{j,w}&\ge0,\qquad
 \sum_{w\ {\rm twisted}}\omega_w(e_w r^0_{j,w}+J_{j,w})
                \le C\mathcal E.\label{eq:joint-parent-global}
\end{align}
On the final members of Construction~\ref{hyp:family-prescribed-realization}, use the reduced chart and product lattices of Lemmas~\ref{hyp:slice-chart-branch-lattices} and~\ref{hyp:slice-product-lattices} on the marked rows. The transport of Lemma~\ref{lem:parent-transport-cost} then compares the lattices at each compatible completion.
The constants are uniform in the translates, words, lifts and
places of the family.
\end{lemma}

\begin{proof}
\proofstep{Choose a comparison that reads the parent row.}
Choose the comparison in the completed division pair with
the prescribed identification of coefficient embeddings, series
reading and geometric core, as in Paragraph~\ref{par:linked-family-rows}.
The primitive and divisor identifications hold on each such
summand. Fix a compatible completed member \(P\), and write \(D(P)\) for
its period maps, integral lattices and numerical sizes. The covolume
comparison and the evaluation-grid comparison may choose different
realizations over \(P\). Denote those identifications by \(\iota_{\rm cov}\) and
\(\iota_{\rm ev}\), on the covolume and evaluation realizations
over \(P\), respectively. Their period-and-lattice records satisfy
\[
         \iota_{\rm cov}(D_{\rm cov})=D(P)
                                =\iota_{\rm ev}(D_{\rm ev}).
\]


More explicitly, let \(C_{\rm cov}(P)\) and \(C_{\rm ev}(P)\) denote the choices of auxiliary comparisons inducing the fixed data \(D(P)\). The constructions on the fixed row give an element of each set. Any pair in
\[
                  C_{\rm cov}(P)\times C_{\rm ev}(P)
\]
gives a simultaneous choice over \(P\). Apply the two inequalities there. Their
row, weight, divisor and coefficient variables then coincide. For each further restriction
with threshold \(X_i\), work on \(x\ge\max_i X_i\) in the
domain of the chosen family.
The rank, divisor data and choice of orientations in
Section~\ref{subsec:completed-prime-runs} now apply there.
\label{cor:joint-family-estimate-1}

\proofstep{Descend the class sums to one coefficient field.}
Let \(A(y)\) denote the complex form of a vector in the lift's
coefficient frame over \(K\), and let \(\iota:K\hookrightarrow
\mathbb C\) be its embedding. The images of a frame basis are
complex linearly independent by the form isomorphism and its
injective reading on the generating lowest vector. We need the
class-sum map at the actual tuple level over this field. Its
rationality uses the ray fields from Paragraph~\ref{par:linked-family-rows}.
Write \(K_{2}\) for the imaginary quadratic subfield. Since it has degree two and contains the indicated square root,
\[
 \mathbb Q(\sqrt{-d})\subseteq K_{2},\qquad [\mathbb Q(\sqrt{-d}):\mathbb Q]=[K_{2}:\mathbb Q]=2,
\]
and the two fields coincide.

At tuple level \(M\ge3\), the Weber values of
the level torsion lie in the embedded ray field by complex
multiplication. With the root of \(-d\) fixed, the Weber comparison
gives the CM curve with the required level structure. The level
of the cover divides \(M\), since its auxiliary level is one.
The component comparison using the rank-one CM datum then gives
geometric irreducibility of its components over the coefficient
field, equivalently relative algebraic closure of that field in
their function fields. Use this comparison when choosing the half
levels and their covers on the normalized reference frame.

On
each tuple take the finite class representatives and their
stabilizer masses in the new level. In an integral frame its local factor takes the form
\[
         K_{t,v}(z)=\mathbf1_{A_{t,v}}(z)
                         \prod_\nu \chi_{\nu,t,v}(g_\nu(z))^{a_\nu},
\]
where \(A_{t,v}\) is the selected finite support and the exponents
are integral. The admitted character values lie in \(\iota K\).
Thus each local factor takes its values there, including zero
off the support. The finite chart products therefore have the
same coefficient field at all tuples. The original one-map comparison also extends from
the reference frame to these tuple levels on the prescribed
quotient primes. The change to the split character reading
preserves its levels, individual chart maps and extremal vector.
\label{cor:joint-family-estimate-2}

Fix a lift of polynomial degree \(d_W\). With the same class factor \(c_\alpha\) in both presentations, their ordered pair of factors is
\[
 (u_1c_\alpha,u_2c_\alpha)=c_\alpha(u_1,u_2).
\]
To see the uniformity in this assertion, compare both readings
of a single class in a standard frame. Identifying the adelic
kernel with the short theta series gives an orientation
\(\delta\in\{1,-1\}\). Evaluation of its Fock coefficient on
the chosen extremal vector gives a phase \(a\), and comparison
on the cover gives a phase \(c_0'\). A sign change of the
polynomial argument acts by its total degree. Consequently the
multiplier for a presentation is
\[
                 u=a(-\delta)^{d_W}/c_0' .
\]

The positive Lie multiplier is
\(L=(\sqrt{2\pi\xi_{\mathbb R}})^{-d_W}\), with
\(\xi_{\mathbb R}>0\) in the prescribed vacuum metric.
The remaining factor \(c_\alpha\) compares the normalized
frames on a component. Compute it once on the reference pair
using its rational translation law for the density. Change of
a marking in the same component gives the same value by the
chart transformation rule. The orientation, extremal-vector
phase and cover phase above are fixed before the tuple and
its class vector. Summing the single-class identity with the
stabilizer masses thus gives the same \(u\) for every tuple
and \(\lvert u c_\alpha\rvert=1\).
\label{cor:joint-family-estimate-3}

For a fixed finite argument the two readings of the series
on a source vector give
\[
 u_1c_\alpha L m_g a_g=u_2c_\alpha L m_g a'_g .
\]
The Lie factor \(L\) is positive. Cancelling \(u_1c_\alpha L\)
gives \(m_g a_g=(u_2/u_1)m_g a'_g\), also when the component
density \(m_g\) vanishes. For the second chart, the determinant character on the
cover and the finite twist have finite orders. Their product
of orders clears both multipliers. With the coefficient-field
character values on the half level, the class sum then
satisfies the descent law on the cover. Its weight is dominant
with positive last weight, so the descended section is cuspidal.

By scalar extension for cusp sections there are finitely many base cusp sections \(f_i\) and coefficients \(a_i\) such that the descended section \(f\) satisfies
\[
 f=\sum_i a_i f_i,\qquad a_i\in\iota K .
\] Read each
base section in the form frame and lift its coefficient
through \(\iota\). Translating that section by the level action
gives the finite translate of its adelic reading with the
prescribed character factor. Incorporating the tuple scalar in its
coordinates gives, for every tuple \(t\) and source \(a\),
\[
       \mathrm{ph}\,\Theta_t(a)=\rho A(y_{t,a}),\qquad
       \rho=u_2/u_1,\quad y_{t,a}\in K^{n_W}.
\]
The equality on the lowest-vector evaluations gives the equality
of forms by injectivity on that space. This uses the same
finite group action for the two readings.
\label{cor:joint-family-estimate-4}

To descend the common scalar, use the good projector \(e\) at
the reference tuple. Its rational lift matrix is bijective by
the base comparison, and the target has dimension \(D_0>0\).
Lift the vector with all coordinates one by its inverse. The
resulting good theta reading is nonzero by injectivity of the
target coordinates.

At this tuple the comparison gives rational
coordinates for the phased good lift. Applying the good projector
to the preceding equality therefore gives
\[
           A(y_1)=\rho A(e y)\ne0.
\]
Complex linear independence gives
\(\rho=\iota((y_1)_h/(e y)_h)\) at any coordinate with
\((e y)_h\ne0\). Absorb this field element into each
\(y_{t,a}\). The phase in the selected comparison is one.
For each tuple restrict the resulting rational map to its level
source. Its preimages under \(A\) are unique, hence additive
and \(K\)-linear by semilinearity of both readings. This
produces the rational class-sum operator \(\theta\) at every
assigned pair \((w,j)\), even for a tuple with zero class sum.

\proofstep{Verify equivariance and restriction on the good space.}
Write the central good idempotent as \(\sum_g a_g[g]\), with
actions \(e_X,e\) on the source and target. The descended
class-sum map intertwines each \(g\), so
\[
     \theta e_Xx=\sum_g a_g g\,\theta(x)=e\theta(x).
\]
Centrality gives the same identity after restriction to the level
character. To project in the source level metric, take a basis
\(k_1,\ldots,k_h\) of \(\ker\theta\) and its Gram matrix \(G\),
using the convention that the Hermitian form is linear in its
second argument. The restricted matrix is invertible at the
selected embedding, since the metric is positive definite there.
Thus the correction has coordinates in the field:
\[
        {\rm pr}(x)=x-\sum_{i=1}^h a_i k_i,\qquad
        Ga=(\langle k_j,x\rangle)_j .
\]

The correction satisfies \(\theta({\rm pr}(x))=\theta(x)\). The base lift is surjective
onto the good rational space. For lattice membership, use
Section~\ref{subsec:paired-row-indices} on this same rational
map. On a tuple generator supported on its class the
inverse-character trace over its stabilizer belongs to the
level lattice. Its class-sum image is the chosen generator
image times an integral character unit. The span comparison gives equality on generic fibres, and at each of the indicated places it gives the integral inclusion
\[
 \mathcal X_{3,w}\subseteq(e\theta(M_t))_w .
\] The two places \(w,\bar w\) have
the same tuple and residue characteristic.
\label{cor:joint-family-estimate-5}

\proofstep{Apply the geometric and arithmetic estimates to one record.}
The local transport inequality and the weighted cohomological numerator have now been obtained. Their sum can be used in the determinant only after checking that their ranks, place labels and lattices all refer to the chosen member.

Use the simultaneous choice of the outer, middle and last
parameters in Section~\ref{subsec:window-points} for the
finitely many family conditions. For the resulting family
and \(C_H\), the run and rationality arguments have cutoffs
\(x_R,x_Q\). The geometric, rank, twist, character-share,
bad-type-share, exponent and numerical size estimates each
have a lower bound on their main comparison constant.
Choose \(c_0\) at least the maximum of these seven bounds
and \(20\), and put
\[
 c=(c_0,14,66c_0+4),\qquad C_b=14+(66c_0+4)+4,\qquad
 C_E=4c_0^2+25c_0+C_b+44 .
\]
The last two entries of \(c\) meet the two bad-prime allowances.
Write \(x_i(c)\), \(1\le i\le7\), for the thresholds of these estimates
at the chosen \(c\), and let \(x_{\rm fam}\) dominate the preceding family
cutoffs.

Take
\[
 X_*=\max\{x_{\rm fam},x_R,x_Q,3c_0,x_1(c),\ldots,x_7(c)\}.
\]
At any point beyond this threshold choose the two runs on its common
geometric core. To build the record \(O\) of
Section~\ref{subsec:completed-prime-runs}, use the second run at the
first distinguished characteristic and the first run everywhere else.
Its section maps and lattices are pulled back through the same
cusp form frame. Its depth uses the canonical half ideals, the
actual error generator and the point's weight data. Every
one of the seven estimates applies to every pair of runs and
record satisfying these identifications. We apply each to
these particular runs and this \(O\).
\label{cor:joint-family-estimate-6}

Here are the hypotheses of those applications.
\begin{enumerate}
\item The numerical sizes come from the parameters of this family
at \(x\). After increasing \(X_*\), the numerical conditions used for this point are simultaneous:
\[
 x\ge3c_0,\qquad C_{n_W}=10,\qquad \mathcal E\ge0 .
\]
The two primes and the level in this application are those of the point. The level and
cutoff comparisons of Section~\ref{subsec:completed-prime-runs}
give the bounds on the level squared and the cutoff \(Y_0\)
in the prime sum.
\item Use the geometric section theorem of that section with
the orientation prescribed at each place. Its short sections
span, their scaled lattice lies in the ambient lattice, and
its remainder is chosen before restriction to a subspace.
Pulling its estimate to the selected good subspace gives the
index estimate of Section~\ref{subsec:paired-row-indices}.
Use its four depth bounds at the pair \((C_E,C_H)\) for the
ordinary, ramified, dyadic and level-prime cases. If \(a_w\) is the scaling scalar at the place under consideration, the forced content used in this application is read at the same place and with the same normalization:
\[
 \text{forced content at }w=v_w(a_w).
\]
\item For the rank comparison take the integral lattice \(J\)
of base-character generators of the selected run. The class representatives of
the base theta kernel give the spanning list. Indeed evaluation of a generator on its own finite
representative recovers the class series, and tracing that
representative over its stabilizer supplies every class in
the list. The deletion-character values are in the point's
field. The off-prime comparison for the used run identifies the
localized short theta generators inside \(J\). The latter is
finitely generated and lies in the level character space.
The finite good projector commutes with the action there.

For the character \(\chi\) and a vector \(v\) in its
space, the unnormalized trace and its projected image are
\[
      \sum_{u\in U}\chi(u)^{-1}uv=|U|v,\qquad
      e\sum_{u\in U}\chi(u)^{-1}uv=|U|ev .
\]
The short sections span the character space over the
fraction field. Since \(|U|\ne0\), their projected
traces give the full generic good space of \(eJ\). Apply this also at the conjugate place, using
the same off-prime choice. \label{cor:joint-family-estimate-7}
\item Now use the simultaneous test for the chosen twisting
character with this full-rank baseline. At a place \(v\) of the chosen pair, Section~\ref{subsec:paired-row-indices} bounds the excess over the full-rank baseline:
\[
 [M_g:\mathcal X_3]_v-[M_g:eJ]_v\le e_v\epsilon_{vj}.
\]
On a pair the sum of both excesses over all lifts is at most
\(2/p^2\). Its application uses the integral inclusion above
and the prescribed character at both places.
\item For the character share the common action on the cusp
forms identifies the trace comparison with that for the
actual weight. For the discarded share the two central
finite-type idempotents act in the same representation.
Equation~\eqref{eq:row-bad-share} bounds the discarded rank by
the regular-representation trace times the multiplicity bound.
Together the two estimates give
\[
 (1-y)R_j\le|U|d_0,\qquad (1-y)|U|\le I,\qquad
 (1-\delta)d_0\le D_0,\qquad y\le1/8,\quad\delta\le1/4 .
\]
Here \(I\) is the level index. Hence \(R_j\le3 I D_0\).
\item Finally the integral exponent comparison at \(p\) uses
the actual pair of lift lattices in the record. The inclusion
by the rational-prime power \(p^{B_j(p)}\) gives
\(0\le J_{j,w}\le e_w B_j(p)\). Grouping the places above each rational prime and using \eqref{eq:division-least-sum} therefore gives
\[
 \sum_{w\mid p}e_w\omega_w=\log p,\qquad \sum_{w\ {\rm twisted}}\omega_w J_{j,w}\le32\mathcal E .
\] The depth and exponent
arguments use the same coefficient field as the signed lengths.
\end{enumerate}

The restriction, signed length and conjugation calculations in
Section~\ref{subsec:paired-row-indices} now apply to \(O\) with
the two class-sum identities and the base inclusion. The character
and central projections in that calculation are the projections
in its pulled-back level action. The lengths and rank therefore
belong to the base row of this very member. By Section~\ref{subsec:parent-row-errors}, at every lift the three contributions to the weighted auxiliary error together obey
\[
 \text{weighted auxiliary error at }j\le C_E\mathcal E\qquad(1\le j\le m).
\] On forgetting the split phase-one comparison,
the normalized row has phase one. Substitution in the period
and numerical readings identifies the matrices, both
lattice indices, \(D_0\), the forced weights and the normalized
exponents \(J_{j,w}/e_w\) with their parent values.
They have these values for every choice of the numerical
allowance \(C\).

Write \(\mathfrak e_j(w)\ge0\) for the auxiliary
transport error of the comparison before descent. The local
bound gives
\[
   B_w:=D_0^{-1}\sum_h I_{wh}
         -2\sum_h e_wF_h(w)-2e_w/p^2
       \le \sum_h e_w\mathfrak e_h(w).
\]
The summands and the right-hand side are read at the selected
completed member. Therefore its defining error
\(e_w r^0_{j,w}=[B_w]_+/m\) is nonnegative and
\(B_w\le\sum_j e_w r^0_{j,w}\), proving
\eqref{eq:joint-parent-local}. For the weighted bound, first
discard the nonnegative least exponents in each auxiliary
transport budget. This gives
\[
 \sum_{w\ {\rm twisted}}\omega_w e_w r^0_{j,w}
 \le \frac1m\sum_h\sum_{w\ {\rm twisted}}
                    \omega_w e_w\mathfrak e_h(w)
 \le C_E\mathcal E.
\]
For the fixed lift \(j\), discard instead the transport term
to obtain \(\sum_{w\ {\rm twisted}}\omega_w J_{j,w}\le C_E\mathcal E\).
Their sum proves \eqref{eq:joint-parent-global}, with
\(C_0=2\max(C_E,0)\), at every member above the common threshold. \label{cor:joint-family-estimate-8}
\end{proof}


